\chapter{$\kappa$-non-collapsed ancient solutions}\label{chap:kappa-ancient-solutions}
\label{kappasect}

 In this chapter we discuss the qualitative properties of
$\kappa$-non-collapsed, ancient solutions. One of the most important is the
existence of a gradient shrinking soliton that is asymptotic at $-\infty$ to
the solution. The other main qualitative result is the compactness result (up
to scaling) for these solutions. Also extremely important for us is
classification of $3$-dimensional gradient shrinking solitons -- up to finite
covers there are only two: a shrinking family of round $S^3$'s and a shrinking
family of products of round $S^2$'s with $\Ar$. This leads to a rough
classification of all $3$-dimensional $\kappa$-non-collapsed, ancient
solutions. The $\kappa$-solutions are in turn the models for singularity
development in $3$-dimensional Ricci flows on compact manifolds, and eventually
for singularity development in $3$-dimensional Ricci flows with surgery.

\section{Preliminaries}

Our objects of study are Ricci flows $(M,g(t)),\ -\infty<t\le 0$,
with each $(M,g(t))$ being a complete manifold of bounded
non-negative curvature. The first remark to make is that the
appropriate notion of non-negative curvature is that the Riemann
curvature operator
$$\Rm\colon \wedge^2TM\to \wedge^2TM,$$
which is a symmetric operator, is non-negative. In general, this
implies, but is stronger than, the condition that the sectional
curvatures are all non-negative. In case the dimension of $M$ is at
most three, every element of $\wedge^2TM$ is represented by a
$2$-plane (with area form) and hence the Riemann curvature operator
is non-negative if and only if all the sectional curvatures are
non-negative.  In the case of non-negative curvature operator,
bounded curvature operator is equivalent to bounded scalar
curvature.

It follows immediately from the Ricci flow equation that since the $(M,g(t))$
have non-negative Ricci curvature,  the metric is non-increasing in time in the
sense that for any point $p\in M$ and any $v\in T_pM$ the function
$|v|^2_{g(t)}$ is a non-increasing function of $t$.

There are  stronger results under the assumption of bounded, non-negative
curvature operator. These are consequences of the Harnack inequality (see
\cite{Hamiltonharnack}). As was established in Corollary~\ref{posderiv}, since
the flow exists for $t\in (-\infty,0]$ and since the curvature operator is
non-negative and bounded for each $(q,t)\in M\times (-\infty,0]$, it follows
that $\partial R(q,t)/\partial t\ge 0$ for all $q$ and $t$. That is to say, for
each $q\in M$ the scalar curvature $R(q,t)$  is a non-decreasing function of
$t$.

\subsection{Definition}

Now we turn to the definition of what it means for a Ricci flow to
be $\kappa$-non-collapsed.

\begin{definition}[$\kappa$-non-collapsed Ricci flow]\label{def:kappa-noncollapsed}
\label{kappadefn}
\uses{def:ricci-flow}
Fix $\kappa>0$. Let $(M,g(t)),\ a<t\le b$, be a Ricci flow of
complete $n$-manifolds. Fix $r_0>0$. We say that $(M,g(t))$ is {\em
$\kappa$-non-collapsed on scales at
most $r_0$} if the following holds for any $(p,t)\in M\times (a,b]$
and any $0<r\le r_0$  with the property that $a\le t-r^2$. Whenever
$|\Rm(q,t')|\le r^{-2}$ for all $q\in B(p,t,r)$ and all $t'\in
(t-r^2,t]$, then $\Vol\,B(p,t,r)\ge \kappa r^n$. We say that
$(M,g(t))$ is {\em $\kappa$-non-collapsed}, or equivalently {\em
$\kappa$-non-collapsed on all scales} if it is
$\kappa$-non-collapsed on scales at most $r_0$ for every
$r_0<\infty$.
\end{definition}


\begin{definition}[Ancient solution and $\kappa$-solution]\label{def:ancient-solution}
\uses{def:ricci-flow,def:kappa-noncollapsed}
An {\em ancient solution} is a Ricci flow
 $(M,g(t))$ defined for $-\infty <t \leq 0$ such
that for each $t$, $(M,g(t))$ is a connected, complete, non-flat
Riemannian manifold whose curvature operator is bounded and
nonnegative. For any $\kappa>0$, an ancient solution is {\sl
$\kappa$-non-collapsed} if it is $\kappa$-non-collapsed on all
scales. We also use the terminology $\kappa$-{\sl solution} for a
$\kappa$-non-collapsed, ancient solution.
\end{definition}

Notice that a $\kappa$-solution is a $\kappa'$-solution for any
$0<\kappa'\le \kappa$.



\subsection{Examples}

Here are some examples of $\kappa$-solutions:
\begin{example}[Shrinking round $2$-sphere is a $\kappa$-solution]\label{ex:shrinking-sphere-kappa-solution}
\uses{def:ancient-solution}
Let $(S^{2},g_{0})$ be the standard round $2$-sphere of scalar
curvature $1$ (and hence Ricci tensor $g_0/2$). Set
$g(t)=(1-t)g_{0}$. Then $\partial g(t)/\partial t=-2\Ric(g(t))$, $-\infty<t\le 0$. This Ricci flow is an ancient
solution which is $\kappa$-non-collapsed on all scales for any
$\kappa$ at most the volume of the ball of radius one in the unit
$2$-sphere.
\end{example}

According to a result of Hamilton which we shall prove below
(Corollary~\ref{2DGSS}):

\begin{theorem}[Two-dimensional $\kappa$-solutions are round]\label{thm:two-dimensional-kappa-solution-preview}
\uses{def:ancient-solution,cor:two-dimensional-gradient-shrinking-soliton}
Every orientable, $2$-dimensional $\kappa$-solution is a rescaling
of the previous example, i.e., is a family of shrinking round
$2$-spheres.
\end{theorem}


\begin{example}[Round $n$-sphere and its quotients]\label{ex:round-sphere-quotient-kappa-solution}
\uses{def:ancient-solution}
Let $(S^n,g_0)$ be the standard round $n$-sphere of scalar curvature
$n/2$. Set $g(t)=(1-t)g_0$. This is a $\kappa$-solution for any
$\kappa$ which is at most the volume of the ball of radius one in
the unit $n$-sphere. If $\Gamma$ is a finite subgroup of the
isometries of $S^n$ acting freely on $S^n$, then the quotient
$S^n/\Gamma$ inherits an induced family of metrics $\bar g(t)$
satisfying the Ricci flow equation. The result is a
$\kappa$-solution for any $\kappa$ at most $1/|\Gamma|$ times the
volume of the ball of radius one in the unit sphere.
\end{example}

\smallskip

\begin{example}[Shrinking cylinder $S^2\times\Ar$]\label{ex:cylinder-kappa-solution}
\uses{def:ancient-solution}
Consider the product $S^{2}\times \mathbb{R} $, with the metric $
g(t) =(1-t)g_{0}+ds^{2}.$ This is a $\kappa$-solution for any
$\kappa$ at most the volume of a ball of radius one in the product
of the unit $2$-sphere with $\Ar$.
\end{example}


\smallskip

\begin{example}[Quotient of the shrinking cylinder]\label{ex:cylinder-quotient-kappa-solution}
\uses{def:ancient-solution,ex:cylinder-kappa-solution}
The quotient $S^{2}\times \mathbb{R}/\langle\iota\rangle$, where the
involution $\iota$ is the product of the antipodal map on $S^2$ with
$s\mapsto -s$ on the $\Ar$ factor,  is an orientable
$\kappa$-solution for some $\kappa>0$.
\end{example}

\begin{example}[Product with a circle is not $\kappa$-non-collapsed]\label{ex:cylinder-with-circle-not-kappa-noncollapsed}
\uses{def:ancient-solution,def:kappa-noncollapsed}
Consider the metric product $(S^{2},g_0) \times (S^{1}_{R},ds^2)$
where $(S^1_R,ds^2)$ is the circle of radius $R$. We define
$g(t)=(1-t)g_0+ds^2$. This is an ancient solution to the Ricci flow.
But it is not $\kappa$-non-collapsed for any $\kappa>0$. The reason
is that
$$|\Rm(p,t)|=\frac{1}{1-t},$$ and
$$\frac{\Vol_{g(t)}\,B(p,\sqrt{1-t})}{(1-t)^{3/2}}
\le \frac{\Vol_{g(t)}(S^{2}\times S^{1}_{R})}{(1-t)^{3/2}}
=\frac{2\pi R(1-t)4\pi}{(1-t)^{3/2}}=\frac{8\pi^2R}{\sqrt{1-t}}.$$
Thus, as $t\rightarrow-\infty$ this ratio goes to zero.
\end{example}



\subsection{A consequence of Hamilton's Harnack inequality}

In order to prove the existence of an asymptotic gradient shrinking
soliton associated to every $\kappa$-solution, we need the following
inequality which is a consequence of Hamilton's Harnack
inequality for Ricci flows
with non-negative curvature operator.

\begin{proposition}[Harnack consequence: Lipschitz bound for reduced length]\label{prop:harnack-l-function-lipschitz-bound}
\label{uniformlip}
\uses{def:ricci-flow}
Let $(M,g(t)),\ -\tau_0\le t\le 0$, be an $n$-dimensional Ricci flow
such that for each $t\in [-\tau_0,0]$ the Riemannian manifold
$(M,g(t))$ is complete with non-negative, bounded curvature
operator. Let $\tau=-t$. Fix a point $p\in M$ and let $x=(p,0)\in
M\times [-\tau_0,0]$. Then for any $0<c<1$ and any $\tau\le
(1-c)\tau_0$ we have
$$|\bigtriangledown l_x(q,\tau))|^{2} +R(q,\tau) \leq \frac{(1+2c^{-1})l_x(q,\tau)}{\tau},
\ \ \ \mathrm{and}$$
$$ R(q,\tau)-\frac{(1+c^{-1})l_x(q,\tau)}{\tau}\le
\frac{\partial l_x}{\partial \tau} $$ where these
inequalities hold on the open subset of full measure of $M\times
[-(1-c)\tau_0,0)$ on which $l_x$ is a smooth function.
\end{proposition}

\begin{proof}
Recall that from Definition~\ref{Hdefnn} we have
$$H(X)=-\frac{\partial R}{\partial
\tau}-\frac{R}{\tau}-2\langle \nabla R,X\rangle+2\Ric(X,X).$$
Using Hamilton's Harnack's inequality (Theorem~\ref{Harnack}) with
$\chi=-X$, we have
\begin{equation*}
-\frac{\partial R}{\partial \tau}-\frac{R}{\tau_0-\tau}-2\langle
\nabla R,X\rangle+2\Ric(X,X)\ge 0.\end{equation*}
 Together these imply
 $$H(X)\ge \left(\frac{1}{\tau-\tau_0}-\frac{1}{\tau}\right)R
 =\frac{\tau_0}{\tau(\tau-\tau_0)}R.$$
Restricting to $\tau\le (1-c)\tau_0$ gives
$$H(X)\ge -\frac{c^{-1}}{\tau}R.$$
Take a minimal ${\mathcal L}$-geodesic from $x$ to $(q,\bar\tau)$,
we have
\begin{equation}\label{Kbartaueqn}
K^{\bar\tau}(\gamma)=\int_0^{\bar\tau}\tau^{3/2}H(X)d\tau\ge
-c^{-1}\int_0^{\bar\tau}\sqrt{\tau}Rd\tau\ge
-2c^{-1}\sqrt{\bar\tau}l_x(q,\bar\tau). \end{equation} Together with
the second equality in Theorem~\ref{Ueqn}, this gives
\begin{align*}
4\bar\tau |\bigtriangledown l_x(q,\bar\tau)|^{2} & = -4\bar\tau
R(q,\bar\tau) + 4l_x(q,\bar\tau) -\frac{4}{\sqrt{\bar\tau}}
K^{\bar\tau}(\gamma)\\ & \le -4\bar\tau R(q,\bar\tau)
+4l_x(q,\bar\tau) + 8c^{-1}l_x(q,\bar\tau)
\end{align*}

Dividing through by $4\bar\tau$, and replacing $\bar\tau$ with
$\tau$ yields the first inequality in the statement of the
proposition:
$$|\nabla l_x(q,\tau)|^2+R(q,\tau)\le \frac{(1+2c^{-1})l_x(q,\tau)}{\tau}$$
for all $0<\tau\le (1-c)\tau_0.$ This is an equation of smooth
functions on the open dense subset ${\mathcal U}(\bar \tau)$ but it
extends as an equation of $L^\infty_\mathrm{loc}$-functions on all of
$M$.

As to the second inequality in the statement, by the first equation
in Theorem~\ref{Ueqn} we have
$$\frac{\partial l_x(q,\tau)}{\partial \tau}
=R(q,\tau)-\frac{l_x(q,\tau)}{\tau}+\frac{1}{2\tau^{3/2}}K^{\tau}(\gamma).$$
The estimate on $K^\tau$ in Equation~(\ref{Kbartaueqn}) then gives
$$R(q,\tau)-\frac{(1+c^{-1})l_x(q,\tau)}{\tau}\le
\frac{\partial l_x(q,\tau)}{\partial \tau}.$$ This establishes the
second inequality.
\end{proof}




\begin{corollary}[Uniform reduced-length gradient and time-derivative bound]\label{cor:l-function-gradient-bound}
\label{unianc}
\uses{prop:harnack-l-function-lipschitz-bound}
Let $(M,g(t)),\ -\infty< t\le 0$, be a Ricci flow on a complete,
$n$-dimensional manifold with bounded, non-negative curvature
operator. Fix a point $p\in M$ and let $x=(p,0)\in M\times
(-\infty,0]$. Then for any $\tau>0$ we have
$$|\bigtriangledown l_x(q,\tau))|^{2} +R(q,\tau) \leq \frac{3l_x(q,\tau)}{\tau},$$
$$-\frac{2l_x(q,\tau)}{\tau}\le \frac{\partial l_x(q,\tau)}{\partial \tau}
\le \frac{l_x(q,\tau)}{\tau}.$$ where these inequalities are valid
in the sense of smooth functions on the open subset of full measure
of $M\times \{\tau\}$ on which $l_x$ is a smooth function, and are
valid as inequalities of $L^\infty_\mathrm{loc}$-functions on all of
$M\times \{\tau\}$.
\end{corollary}


\begin{proof}
\uses{prop:harnack-l-function-lipschitz-bound}
Fix $\tau$ and take a sequence of $\tau_0\rightarrow \infty$,
allowing us to take $c\rightarrow 1$, and apply the previous
proposition. This gives the first inequality and gives the lower
bound for $\partial l_x/\partial \tau$ in the second inequality.

 To establish the upper bound in the second inequality we
consider the path that is the concatenation of a minimal ${\mathcal
L}$-geodesic $\gamma$ from $x$ to $(q,\tau)$ followed by the path
$\mu(\tau')=(q,\tau')$ for $\tau'\ge \tau$. Then
$$l_x(\gamma*\mu|_{[\tau,\tau_1]})=\frac{1}{2\sqrt{\tau_1}}\left({\mathcal L}(\gamma)
+\int_{\tau}^{\tau_1} \sqrt{\tau'}R(q,\tau')d\tau'\right).$$
Differentiating at $\tau_1=\tau$ gives
\begin{eqnarray*}
\frac{\partial
l_x(\gamma*\mu)}{\partial\tau}\Bigl|_{\tau_1=\tau}\Bigr. & = &
-\frac{1}{4\tau^{3/2}}{\mathcal L}(\gamma)+
\frac{1}{2\sqrt{\tau}}\sqrt{\tau}R(q,\tau) \\
& = & -\frac{l_x(q,\tau)}{2\tau}+\frac{R(q,\tau)}{2}.\end{eqnarray*}
By the first inequality in this statement, we have
$$-\frac{l_x(q,\tau)}{2\tau}+\frac{R(q,\tau)}{2}\le
 \frac{l_x(q,\tau)}{\tau}.$$ Since $l_x(q,\tau')\le \tilde l(\gamma*\mu|_{[\tau,\tau']})$
 for all $\tau'\ge \tau$, this establishes the claimed upper
bound for $\partial l_x/\partial \tau$.
\end{proof}

\section{The asymptotic gradient shrinking soliton for
$\kappa$-solutions}\label{sect9.2}




We fix $\kappa>0$ and we consider an $n$-dimensional
$\kappa$-solution $(M,g(t)),\ -\infty <t\leq 0.$ Our goal in this
section is to establish the existence of an asymptotic gradient
shrinking soliton associated to this $\kappa$-solution. Fix a
reference point $p \in M$ and set $x=(p,0)\in M\times (-\infty,0]$.
By Theorem~\ref{n/2} for every $\tau>0$ there is a point $q(\tau)\in
M$ at which the function $l_x(\cdot,^\tau)$ achieves its minimum,
and furthermore, we have
$$l_x(q(\tau),\tau) \leq \frac{n}{2}.$$

 For $\bar\tau>0,$ define
$$g_{\bar\tau}(t) = \frac{1}{\bar\tau}g(\bar\tau t),\ \  -\infty<t\le 0.$$

Now we come to one of the main theorems about $\kappa$-solutions, a
result that will eventually provide a qualitative description of all
$\kappa$-solutions.

\begin{theorem}[Existence of the asymptotic gradient shrinking soliton]\label{thm:asymptotic-gradient-shrinking-soliton}
\label{asympGSS}
\uses{def:ancient-solution,def:kappa-noncollapsed,def:gradient-shrinking-soliton}
Let $(M,g(t)),\ -\infty<t\le 0$,
be a $\kappa$-solution of dimension $n$. Fix $x=(p,0)\in
M\times(-\infty,0]$. Suppose that $\{\bar\tau_{k}\}_{k=1}^\infty $
is a sequence tending to $ \infty$ as $k\rightarrow\infty$. Then,
after replacing $\{\bar\tau_k\}$ by a subsequence, the following
holds. For each $k$ denote by $M_k$ the manifold $M$, by $g_k(t)$
the family of metrics $g_{\bar\tau_k}(t)$ on $M_k$, and by $q_k\in
M_k$ the point $q(\bar\tau_k)$. The sequence of pointed flows
$(M_{k}, g_k(t), (q_{k},-1))$ defined for $t \in (-\infty,0)$
converges smoothly to a non-flat based Ricci flow $(M_{\infty},
g_{\infty}(t), (q_{\infty},-1))$ defined for $t\in (-\infty,0)$.
This limiting Ricci flow  satisfies the gradient shrinking
soliton equation in the sense
that there is a smooth function $f\colon
M_\infty\times(-\infty,0)\to \Ar$ such that for every
$t\in(-\infty,0)$ we have
\begin{equation}\label{GSSequation}
 \Ric_{g_{\infty}(t)} + \Hess^{g_{\infty}(t)}(f(t)) +
\frac{1}{2t}g_{\infty}(t) = 0.
\end{equation}
Furthermore, $(M_\infty,g_\infty(t))$ has non-negative curvature operator, is
$\kappa$-non-collapsed, and satisfies $\partial R_{g_\infty}(x,t)/\partial t\ge
0$ for all $x\in M_\infty$ and all $t<0$.
\end{theorem}

\begin{remark}[Caveats about the asymptotic soliton]\label{rem:asymptotic-soliton-caveats}
\uses{thm:asymptotic-gradient-shrinking-soliton}
We are not claiming that the gradient shrinking soliton is a
$\kappa$-solution (or more precisely an extension forward in time of
a time-shifted version of a $\kappa$-solution) because we are not
claiming that the time-slices have bounded curvature operator.
Indeed, we do not know if this is true in general. We shall
establish below (see Corollary~\ref{2DGSS} and
Corollary~\ref{3DGSSkappa}) that in the case $n=2,3$, the gradient
shrinking soliton does indeed have time-slices of bounded curvature,
and hence is an extension of a $\kappa$-solution. We are also not
claiming at this point that the limiting flow is a gradient
shrinking soliton in the sense that there is a one-parameter family
of diffeomorphisms $\varphi_t\colon M_\infty\to M_\infty,\ t<0$,
with the property that $|t|\varphi_t^*g_\infty(-1)=g_\infty(t)$ and
with the property that the $\varphi_t$ are generated by the gradient
vector field of a function. We shall also
 establish this result in dimensions $2$ and $3$ later in this chapter.
\end{remark}



We will divide the proof of Theorem~\ref{thm:asymptotic-gradient-shrinking-soliton} into steps.
First, we will show that the reduced length and norm of the
curvature $|\Rm|$ are bounded throughout the sequence in some
way. Then using the $\kappa$-non-collapsed assumption, by the
compactness theorem (Theorem~\ref{flowlimit}), we conclude that a
subsequence of the sequence of flows converges geometrically to a
limiting flow. Then, using the fact that the limit of the reduced
volumes, denoted $\widetilde
V_{\infty}(M_\infty\times\{t\})$, is
constant we show that the limit flow is a gradient shrinking
soliton. Finally we argue that
the limit is non-flat. The proof occupies the rest of
Section~\ref{sect9.2}.

\subsection{Bounding the reduced length and the curvature}
Now let's carry this procedure out in detail. The first remark is
that since rescaling does not affect the $\kappa$-non-collapsed
hypothesis, all the Ricci flows $(M_k,g_k(t))$ are
$\kappa$-non-collapsed on all scales. Next, we have the effect on
reduced volume.

\begin{claim}[Reduced volume rescaling]\label{claim:reduced-volume-rescaling}
\label{scaleinv2}
\uses{thm:asymptotic-gradient-shrinking-soliton}
For each $k\ge 1$ denote by $x_k\in M_k$ the point $(p,0)\in
M_k$. Let $\widetilde V_{x_k}(\tau)=\widetilde
V_{x_k}(M_k\times\{\tau\})$ denote the reduced volume function for
the Ricci flow $(M_k,g_k(t))$ from the point $x_k$, and let
$\widetilde V_x(\tau)$ denote the reduced volume of
$M\times\{\tau\}$ for the Ricci flow $(M,g(t))$ from the point $x$.
Then
$$\widetilde V_{x_k}(\tau) = \widetilde V_x (\bar \tau_{k}\tau).$$
\end{claim}


\begin{proof}
This is a special case of the reparameterization equation for
reduced volume (Lemma~\ref{Qformula}).
\end{proof}

By Theorem~\ref{4pi} the reduced volume function $\widetilde
V_x(\tau)$ is a non-increasing function of $\tau$ with $\lim_{\tau\rightarrow 0}\widetilde V_x(\tau)=(4\pi)^{\frac{n}{2}}$.
Since the integrand for $\widetilde V_x(\tau)$ is everywhere
positive, it is clear that $\widetilde V_x(\tau)>0$ for all $\tau$.
Hence, $\lim_{\tau\rightarrow\infty}\widetilde V_x(\tau)$
exists. By Corollary~\ref{4picase} either this limit as $\tau$ goes
to infinity is less than $(4\pi)^{n/2}$ or the flow is the constant
flow on flat Euclidean space. The latter is ruled out by our
assumption that the manifolds are non-flat. It follows immediately
from this and Claim~\ref{claim:reduced-volume-rescaling} that:

\begin{corollary}[Reduced volume limit is constant]\label{cor:reduced-volume-limit-constant}
\label{redvolconst}
\uses{claim:reduced-volume-rescaling}
There is a non-negative constant $V_\infty<(4\pi)^{n/2}$ such that
for all $\tau\in(0,\infty)$, we have
\begin{equation}\label{limVeqn}
\lim_{k\rightarrow\infty}\widetilde V_{x_k}(\tau)=V_\infty.
\end{equation}
\end{corollary}

Now let us turn to the length functions $l_{x_k}$.


\begin{claim}[Bound for the reduced length at the minimizing point]\label{claim:qk-length-bound}
\label{tausqtau}
\uses{thm:asymptotic-gradient-shrinking-soliton,cor:l-function-gradient-bound}
For any $\tau>0$ we have
$$l_{x_k}(q_k,\tau)\le \frac{n}{2\tau^2}+\frac{n\tau}{2}.$$
\end{claim}

\begin{proof}
\uses{cor:l-function-gradient-bound}
 By the choice of $q_k$ we have $l_{x_k}(q_k,\tau_k)\le \frac{n}{2}$.
 By the scale invariance of $l$ (Corollary~\ref{COR}) we have
 $l_{x_k}(q_k,-1)\le n/2$ for all $k$. Fix $0<\tau<1$. Integrating the inequality
$$\frac{-2l_x(q_k,\tau)}{\tau}\le \frac{\partial l_{x_k}(q_k,\tau)}{\partial \tau}$$
 from $\tau$ to $1$  yields
$$l_{x_k}(q_k,\tau)\le \frac{n}{2\tau^2}.$$
If $\tau>1$, then
 integrating the second inequality in the second displayed line of Corollary~\ref{cor:l-function-gradient-bound} gives
 $l_{x_k}(q_k,\tau)\le \frac{n\tau}{2}$.
\end{proof}

\begin{corollary}[Quadratic distance bound for the reduced length]\label{cor:l-function-distance-bound}
\label{lxkest}
\uses{claim:qk-length-bound}
There is a  positive continuous function $C_1(\tau)$ defined for
$\tau>0$ such that for any $q\in M_k$ we have:
$$l_{x_k}(q,\tau)\le
\left(\sqrt{\frac{3}{\tau}}d_{g_k(-\tau)}(q_k,q)+C_1(\tau)\right)^2,$$
$$|\nabla l_{x_k}(q,\tau)|\le
\frac{3}{\tau}d_{g_k(-\tau)}(q_k,q)+\sqrt{\frac{3}{\tau}}C_1(\tau).$$
\end{corollary}

\begin{proof}
\uses{cor:l-function-gradient-bound,claim:qk-length-bound,prop:harnack-l-function-lipschitz-bound}
 By Corollary~\ref{cor:l-function-gradient-bound}, for any $q\in M_k$ we have $|\nabla
l_{x_k}(q,\tau)|^2\le 3l_{x_k}(q,\tau)/\tau$. Since
$l_{x_k}(q_k,\tau)\le \frac {n}{2\tau_0^2}+\frac{n\tau}{2}$,
integrating yields
$$l_{x_k}(q,\tau)\le
\left(\sqrt{\frac{3}{\tau}}d_{g_k(-\tau)}(q_k,q)+C_1(\tau)\right)^2,$$
with $C_1(\tau)$ being $\sqrt{(n/2\tau^2)+(n\tau/2)}$. The second
statement follows from this and Proposition~\ref{prop:harnack-l-function-lipschitz-bound}.
\end{proof}

It follows immediately from Corollary~\ref{cor:l-function-distance-bound} that for each
$A<\infty$ and $\tau_0>0$, the functions $l_{x_k}$ are uniformly
bounded (by a bound that is independent of $k$ but depends on
$\tau_0$ and $A$) on the balls $B(q_k,-\tau_0,A)$. Once we know that
the $l_{x_k}$ are uniformly bounded on $B(q_k,-\tau_0,A)$, it
follows from Corollary~\ref{cor:l-function-gradient-bound} that $R_{g_k}$ are also
uniformly bounded on the $B(q_k,-\tau_0,A)$. Invoking
Corollary~\ref{posderiv},
 we see
that for any $A<\infty$ the scalar curvatures of the metrics $g_k$
are uniformly bounded on $B_{g_k}(q_k,-\tau_0,A)\times (-\infty,
-\tau_0]$. Since the metrics have non-negative curvature operator,
this implies that the eigenvalues of this operator are  uniformly
bounded on these regions.
 Since we are assuming
that the original Ricci flows are $\kappa$-non-collapsed on all
scales, it follows from Theorem~\ref{flowlimit} that after passing
to a subsequence there is a geometric limit
$(M_\infty,g_\infty(t),(q_\infty,-1)),\ -\infty<t\le -\tau_0$, which
is a Ricci flow which is $\kappa$-non-collapsed on all scales.

Since this is true for every $\tau_0>0$, by a standard
diagonalization argument passing to a further subsequence we get a
geometric limit flow $(M_\infty,g_\infty(t),(q_\infty,-1)),\
-\infty<t<0$.


Let us summarize our progress to this point.

\begin{corollary}[Existence of the geometric limit flow]\label{cor:kappa-solution-geometric-limit-exists}
\uses{thm:asymptotic-gradient-shrinking-soliton,def:kappa-noncollapsed}
After passing to a subsequence of the $\tau_k$ there is a smooth
limiting flow of the $(M_k,g_k(t),(q_k,-1)),-\infty<t\le 0$,
$$(M_\infty,g_\infty(t),(q_\infty,-1)),$$ defined for  $-\infty<t<0$. For every $t<0$
the Riemannian manifold $(M_\infty,g_\infty(t))$ is complete of
non-negative curvature. The flow is $\kappa$-non-collapsed on all
scales and satisfies $\partial R/\partial t\ge 0$.
\end{corollary}

\begin{proof}
 Since the  flows in the sequence are all
$\kappa$-non-collapsed on all scales and have non-negative curvature
operator, the limiting flow is $\kappa$-non-collapsed on all scales
and has non-negative curvature operator. By the consequence of
Hamilton's Harnack inequality (Corollary~\ref{posderiv}), we have
$\partial R/\partial t\ge 0$ for the original $\kappa$-solution.
This condition also passes to the limit.
\end{proof}

\subsection{The limit function}

The next step in the proof is to construct the limiting function
$l_\infty$ of the $l_{x_k}$ and show that it satisfies the gradient
shrinking soliton equation.

By definition of the geometric limit, for any compact connected set
$K\subset M_\infty$ containing $q_\infty$ and any compact
subinterval $J$ of $(-\infty,0)$ containing $-1$, for all $k$
sufficiently large we have smooth embeddings $\psi_k\colon K\to M_k$
sending $q_\infty$ to $q_k$ so that the pullbacks of the
restrictions of the family of metrics $g_k(t)$ for $t\in J$ to $K$
converge uniformly in the $C^\infty$-topology to the restriction of
$g_\infty(t)$ on $K\times J$. Take an exhausting sequence $K_k\times
J_k$ of such products of compact sets with closed intervals, and
pass to a subsequence so that for all $k$ the diffeomorphism
$\psi_k$ is defined on $K_k\times J_k$. We denote by $l_k$ the
pullback of $l_{x_k}$ under these embeddings and by $h_k(t)$ the
pullback of the family of metrics $g_k(t)$. We denote by $
\nabla^{h_k}$ the gradient with respect to $h_k(t)$, and similarly
$\triangle^{h_k}$ denotes the Laplacian for the metric $h_k(t)$. By
construction, for any compact subset of $M_\infty\times (-\infty,0)$
for all $k$ sufficiently large the function $ l_k$ is defined on the
compact set. We use $\nabla$ and $\triangle$ to refer to the
covariant derivative and the Laplacian in the limiting metric
$g_\infty$.


Now let us consider the functions $l_{x_k}$. According to
Corollary~\ref{cor:l-function-distance-bound}, for any $A<\infty$ and any $0<\tau_0<T$,
both $l_{x_k}$ and $|\nabla l_{x_k}|$ are uniformly bounded on
$B(q_k,-1,A)\times [-T,-\tau_0]$ independent of $k$. Hence, the
$l_{x_k}$ are uniformly Lipschitz on these subspaces. Doing this for
each $A$, $\tau_0$, and $T$ and using a standard diagonalization
argument then shows that, after transferring to the limit, the
functions $l_k$ are uniformly locally bounded and uniformly locally
Lipschitz on $M_\infty\times (-\infty,0)$ with respect to the
limiting metric $g_\infty$.

Fix $0<\alpha<1$. Passing to a further subsequence if necessary, we
can arrange that the $l_k$ converge strongly
in $C^{0,\alpha}_\mathrm{loc}$ to a function $l_\infty$ defined on
$M_\infty\times(-\infty,0)$. Furthermore, it follows that the
restriction of $l_\infty$ is locally Lipschitz, and hence the
function $l_\infty$ is an element of $W^{1,2}_\mathrm{loc}\left(M_\infty\times(-\infty,0)\right)$. Also, by passing to a
further subsequence if necessary, we can assume that the $l_k$
converge weakly in $W^{1,2}_\mathrm{loc}$ to $l_\infty$.


\begin{corollary}[Gradient bound for the limit length function]\label{cor:l-infinity-gradient-bound}
\label{linfest}
\uses{cor:l-function-distance-bound}
For any $\tau>0$ and any $q$ we have
$$|\nabla l_\infty(q,\tau)|\le
\frac{3}{\tau}d_{g_\infty}(-\tau)(q_\infty,q)+\sqrt{\frac{3}{\tau}}C_1(\tau),$$
where $C_1(\tau)$ is the continuous function from
Corollary~\ref{cor:l-function-distance-bound}.
\end{corollary}


\begin{proof}
This is immediate from Corollary~\ref{cor:l-function-distance-bound} and Fatou's lemma.
\end{proof}

\begin{remark}[$l_\infty$ need not be a reduced length function]\label{rem:l-infinity-not-reduced-length}
{\bf N.B.} We are not claiming that $l_\infty$ is the reduced length
function from a point of $M_\infty\times (-\infty,0)$.
\end{remark}

\subsection{Differential inequalities for $l_\infty$}

The next step is to establish differential equalities for $l_\infty$
related to, but stronger than, those that we established in
Chapter~\ref{newcomp2} for $l_x$. Here is a crucial result.

\begin{proposition}[Differential equalities for the limit length function]\label{prop:l-infinity-differential-equations}
\label{weakeqn}
\uses{cor:l-infinity-gradient-bound}
The function $l_\infty$ is a smooth function
on $M\times (-\infty,0)$ and satisfies the following two
differential equalities: \begin{equation}\label{weak1}
\frac{\partial l_\infty}{\partial \tau}+|\nabla
l_\infty|^2-R+\frac{n}{2\tau}-\triangle l_\infty=0\end{equation}
 and
\begin{equation}\label{weak2}2\triangle l_\infty-|\nabla
l_\infty|^2+R+\frac{l_\infty-n}{\tau}=0.\end{equation}
\end{proposition}

The proof of this result is contained in Sections~\ref{sect9.2.3}
through~\ref{sect9.2.5}

\subsection{Preliminary results toward the proof of
Proposition~{\ref{prop:l-infinity-differential-equations}}}\label{sect9.2.3}

In this subsection we shall prove  that the left-hand side of
Equation~(\ref{weak1}) is a distribution and is $\ge 0$ in the
distributional sense. We shall also show that this distribution
extends to a continuous linear functional on compactly supported
functions in $W^{1,2}$.

 The first step in the proof of
this result is the following, somewhat delicate lemma.

\begin{lemma}[Convergence of the gradient-squared term]\label{lem:gradient-convergence-distributions}
\label{nablaconv}
\uses{cor:l-infinity-gradient-bound}
For any $t\in (-\infty,0)$ we have
$$\lim_{k\rightarrow\infty}|\nabla^{h_k} l_k|_{h_k}^2d\vol(h_k)
=|\nabla l_\infty|_{g_\infty}^2 d\vol(g_\infty)$$ in the sense
of distributions on $M_\infty\times \{t\}$.
\end{lemma}

\begin{proof}
It suffices to fix $0<\tau_0<|t|$.
 The inequality in one direction ($\ge$) is a general
result. Here is the argument.  Since the $|\nabla^{g_k} l_{x_k}|_{g_k}$ are
uniformly essentially bounded on every $B(x_k,-\tau_0,A)\times [-T,-\tau_0]$,
the $|\nabla^{h_k}l_k|_{h_k}$ are uniformly essentially bounded on
$B(x_\infty,-\tau_0,A)\times [-T,-\tau_0]$. (Of course,
$\nabla^{h_k}l_k=dl_k=\nabla l_k$.) Since the $h_k$ converge uniformly on
compact sets to $g_\infty$, it is clear that
\begin{equation}\label{firstdist} \lim_{k\rightarrow\infty}\left(|\nabla^{h_k}l_k|^2_{h_k}d\vol(h_k)-|\nabla l_k|_{g_\infty}^2d\vol(g_\infty)\right)=0\end{equation} in the sense of distributions
on $M\times \{t\}$.
 Since the $
l_k$ converge uniformly on compact subsets to $l_\infty$, it follows
immediately from Fatou's lemma that
$$\lim_{k\rightarrow \infty}|\nabla
l_k|^2_{g_\infty}d\vol(g_\infty)\ge|\nabla
l_\infty|^2_{g_\infty}d\vol (g_\infty)$$ in the sense of
distributions on $M_\infty\times \{t\}$. Thus, we have the following
inequality of distributions:
$$\lim_{k\rightarrow\infty}|\nabla^{h_k}l_k|^2_{h_k}d\vol(h_k)\ge
|\nabla l_\infty|^2_{g_\infty}d\vol (g_\infty).$$

 We need to establish the opposite
inequality which is not a general result, but rather relies on  the
bounds on $\triangle^{g_k}l_{x_k}$ (or equivalently on
$\triangle^{h_k}l_k$) given in second inequality in
Theorem~\ref{weak}. We must show that for each $t\le -\tau_0$ and
for any $\varphi$, a non-negative, smooth function with compact
support in $M_\infty\times \{t\}$, we have
$$\lim_{k\rightarrow\infty}\int_{M\times\{t\}}\varphi
\left(|\nabla^{h_k} l_k|_{h_k}^2d\vol(h_k)-|\nabla
l_\infty|_{g_\infty}^2d\vol(g_\infty)\right)\le 0.$$ First,
notice that since, on the support of $\varphi$, the metrics $h_k$
converge uniformly in the $C^\infty$-topology to $g_\infty$ and
since $|\nabla^{h_k}l_k|_{h_k}^2$ and $|\nabla
l_\infty|_{g_\infty}^2$ are essentially bounded on the support of
$\varphi$, we have
\begin{eqnarray}\nonumber
\lefteqn{\lim_{k\rightarrow\infty}\int_{M\times\{t\}}\varphi
\left(|\nabla^{h_k}l_k|_{h_k}^2d\vol(h_k)-|\nabla
l_\infty|_{g_\infty}^2d\vol(g_\infty) \right)} & & \\
&  = & \lim_{k\rightarrow\infty}\int_{M\times\{t\}}\varphi (|
\nabla^{h_k} l_k|_{h_k}^2-|\nabla
l_\infty|_{h_k}^2)d\vol(h_k) \nonumber \\
\nonumber& = & \lim_{k\rightarrow\infty}\int_{M\times\{t\}}\langle\nabla^{h_k}
l_k-\nabla l_\infty),\varphi\nabla^{h_k} l_k\rangle_{h_k}d\vol(h_k)
\\ & & \hskip.55in+\int_{M\times\{t\}}\langle\nabla^{h_k}l_k-\nabla l_\infty),
\varphi\nabla l_\infty\rangle_{h_k}d\vol(h_k).
\label{11}\end{eqnarray} We claim that, in the limit, the last term
in this expression vanishes. Using the fact that the $h_k$ converge
uniformly in the $C^\infty$-topology to $g_\infty$ on the support of
$\varphi$, and $|\nabla l_\infty|$ is bounded on this support we can
rewrite the last term as
\begin{equation}\label{12}
\lim_{k\rightarrow\infty}\int_{M\times\{t\}}\langle\nabla(
l_k-l_\infty),\varphi\nabla l_\infty\rangle_{g_\infty}d\vol(
g_\infty).\end{equation} Since $ l_k-l_\infty$ goes to zero weakly
in $W^{1,2}$ on the support of $\varphi$  whereas $ l_\infty$ is an
element of $W^{1,2}$ of this compact set, we see that the expression
given in~(\ref{12}) vanishes and hence that
$$\lim_{k\rightarrow\infty}\int_{M\times\{t\}}\langle\nabla^{h_k}(
l_k-l_\infty),\varphi\nabla l_\infty\rangle_{h_k}d\vol(h_k)=0.$$

It remains to consider the first term in the last expression in
Equation~(\ref{11}). (This is where we shall need the differential
inequality for the $\triangle^{g_k}l_{x_k}$.) Since the $l_k$
converge uniformly to $l_\infty$ on the support of $\varphi$, we can
choose positive constants $\epsilon_k$ tending to  $0$ as $k$ tends
to $\infty$ so that $l_\infty- l_k+\epsilon_k>0$ on the support of
$\varphi$. We can rewrite
\begin{eqnarray*}\lefteqn{\lim_{k\rightarrow\infty}\int_{M\times\{t\}}\langle\left(\nabla^{h_k}
l_k-\nabla l_\infty\right),\varphi\nabla^{h_k} l_k\rangle_{h_k}
d\vol(h_k) = } \\
& & \lim_{k\rightarrow\infty}\int_{M\times\{t\}}\langle\nabla^{h_k}(
l_k-l_\infty-\epsilon_k),\varphi\nabla^{h_k}l_k\rangle_{h_k} d\vol(h_k).\end{eqnarray*}
\begin{claim}[Vanishing limsup in the gradient identity]\label{claim:nabla-convergence-limsup-bound}
$$\lim_{k\rightarrow\infty}\int_{M\times\{t\}}\langle\nabla^{h_k}(l_k-l_\infty-\epsilon_k),
\varphi\nabla^{h_k}  l_k\rangle_{h_k} d\vol(h_k)\le 0.$$
\end{claim}

\begin{proof}
Since $\varphi$ is a compactly supported, non-negative smooth
function, it follows from Theorem~\ref{weak} that we have the
following inequality of distributions:
$$\varphi\triangle^{h_k} l_k\le
\frac{\varphi}{2}\left(|\nabla^{h_k} l_k|_{h_k}^2-R_{h_k}-\frac{
l_k-n}{\tau}\right).$$ (Here $R_{h_k}$ is the scalar curvature of
$h_k$.)
 That is to say, for any non-negative
$C^\infty$-function $f$ we have
\begin{eqnarray*}
\lefteqn{\int_{M\times\{t\}}-\langle\nabla^{h_k}
l_k,\nabla^{h_k}(\varphi \cdot f)\rangle_{h_k} d\vol(h_k) \le}
\\
& & \int_{M\times\{t\}}\frac{\varphi f}{2}\left(|\nabla^{h_k}
l_k|_{h_k}^2-R_{h_k}-\frac{l_k-n}{\tau}\right)d\vol(h_k).
\end{eqnarray*}

We claim that the same inequality holds as long as $f$ is a
non-negative, locally  Lipschitz function. The point is that given
such a function $f$, we can find a sequence of non-negative
$C^\infty$-functions $f_k$ on the support of $\varphi$ (by say
mollifying $f$)  that converge to $f$ strongly in the $W^{1,2}$-norm
on the support of $\varphi$. The sought-after inequality holds for
every $f_k$. Since both sides of the inequality are continuous in
the $W^{1,2}$-norm of the function, the result holds for the limit
function $f$ as well.

Now we apply this with $f$ being the non-negative locally Lipschitz
function $l_\infty-l_k+\epsilon_k$. We conclude that
\begin{eqnarray*}
\lefteqn{\int_{M\times\{t\}}\langle \nabla^{h_k}(\varphi(
l_k-l_\infty-\epsilon_k)),\nabla^{h_k} l_k\rangle_{h_k}d\vol(h_k) \le}  \\
& & \int_{M\times\{t\}}\frac{\varphi(l_\infty-
l_k+\epsilon_k)}{2}\left(|\nabla^{h_k}
l_k|_{h_k}^2-R_{h_k}-\frac{l_k-n}{\tau}\right)d\vol(h_k).
\end{eqnarray*}
Now taking the limit as $k\rightarrow\infty$, we see that the
right-hand side of this inequality tends to zero since $(l_\infty-
l_k+\epsilon_k)$ tends uniformly to zero on the support of $\varphi$
and $|\nabla^{h_k}l_k|_{h_k}^2$, $R_k$ and $l_k$ are all uniformly
essentially bounded on the support of $\varphi$. Thus, the term
$$\int_{M\times\{t\}}\langle\nabla^{h_k}(\varphi(
l_k-l_\infty-\epsilon_k)),\nabla^{h_k} l_k\rangle_{h_k}d\vol(h_k)$$ has a limsup $\le 0$ as $k$ tends to $\infty$. Now we
expand
$$\nabla^{h_k}(\varphi(l_k-l_\infty-\epsilon_k))=\nabla^{h_k}(\varphi)(l_k-l_\infty-\epsilon_k)+
\varphi\nabla^{h_k}(l_k-l_\infty-\epsilon_k).$$ The first term on
the right-hand side converges to zero as $k\rightarrow\infty$ since
$l_k-l_\infty-\epsilon_k$ tends uniformly to zero on the support of
$\varphi$. This completes the proof of the claim.
\end{proof}
We have now established the  inequalities in both directions and
hence completed the proof of Lemma~\ref{lem:gradient-convergence-distributions}.
\end{proof}




\begin{lemma}[The candidate distribution is non-positively paired]\label{lem:distribution-d-nonpositive}
\uses{lem:gradient-convergence-distributions}
Consider the distribution
$${\mathcal D}=\frac{\partial
l_\infty}{\partial \tau}+|\nabla
l_\infty|^2-R+\frac{n}{2\tau}-\triangle l_\infty$$ on
$M_\infty\times (-\infty,0)$. Then ${\mathcal D}$ extends to a
continuous linear functional on the space of compactly supported
$W^{1,2}$-functions on $M_\infty\times (-\infty,0)$. Furthermore, if
$\psi$ is a non-negative Lipschitz function on $M_\infty\times
(-\infty,0)$ with compact support, then ${\mathcal D}(\psi)\le 0$.
\end{lemma}

\begin{proof}
\uses{lem:gradient-convergence-distributions}
Clearly, since the $l_k$ converge uniformly on compact subsets of
$M_{\infty}\times (-\infty,0)$ to $l_\infty$ and the metrics $h_k$
converge smoothly to $g_\infty$, uniformly on compact sets, it
follows that the $\triangle ^{h_k}l_k $ converge in the weak sense
to $\triangle l_\infty$ and similarly, the $\partial l_k/\partial
\tau$ converge in the weak sense to $\partial l_\infty/\partial
\tau$. Hence, by taking limits from Theorem~\ref{weak}, using
Lemma~\ref{lem:gradient-convergence-distributions}, we see that
\begin{equation}\label{1stlinftyeqn}
{\mathcal D}=\frac{\partial l_\infty}{\partial \tau}+|\nabla
l_\infty|^2-R+\frac{n}{2\tau}-\triangle l_\infty\ge 0\end{equation}
 in the weak sense on $M\times (-\infty,0)$.


Since $R$ and $\frac{n}{2\tau}$ are $C^\infty$-functions, it is
clear that the distributions given by these terms extend to
continuous linear functionals on the space of compactly supported
$W^{1,2}$-functions. Similarly, since $|\nabla l_\infty|^2$ is an
element of $L^\infty_\mathrm{loc}$, it also extends to a continuous
linear functional on compactly supported $W^{1,2}$-functions. Since
$|\partial l_\infty/\partial \tau|$ is an locally essentially
bounded function, $\partial l_\infty/\partial \tau$ extends to a
continuous functional on the space of compactly supported $W^{1,2}$
functions. Lastly, we consider $\triangle l_\infty$. As we have
seen, the value of the associated distribution on $\varphi$ is given
by
$$\int_{M\times (-\infty,,0)}-\langle \nabla \varphi,\nabla
l_\infty\rangle_{g_\infty} d\vol(g_\infty)d\tau.$$ Since
$|\nabla l_\infty|$ is a locally essentially bounded function, this
expression also extends to a continuous linear functional on
compactly supported $W^{1,2}$-functions.




Lastly, if $\psi$ is an element of $W^{1,2}$ with compact support
and hence can be approximated in the $W^{1,2}$-norm by non-negative
smooth functions. The last statement  is now immediate from
Equation~(\ref{1stlinftyeqn}).
\end{proof}


This leads immediately to:

\begin{corollary}[Weighted distribution is non-negative]\label{cor:weighted-distribution-nonnegative}
\label{etimesdis}
\uses{lem:distribution-d-nonpositive}
The functional
$$\varphi\mapsto{\mathcal D}(e^{-l_\infty}\varphi)$$ is a distribution and
its value on any non-negative, compactly supported $C^\infty$-function
$\varphi$ is $\ge 0$.
\end{corollary}


\begin{proof}
\uses{lem:distribution-d-nonpositive}
If $\varphi$ is a compactly supported non-negative
$C^\infty$-function, then $e^{-l_\infty}\varphi$ is a compactly
supported non-negative Lipschitz function. Hence, this result is an
immediate consequence of the previous corollary.
\end{proof}

\subsection{Extension to non-compactly supported functions}

The next step in this proof is to estimate the $l_{x_k}$ uniformly from below
in order to show that  the integrals involved in the distributions in
Proposition~\ref{prop:l-infinity-differential-equations} are absolutely convergent so that they extend to
continuous functionals on a certain  space of functions that includes
non-compactly supported functions.

\begin{lemma}[Lower bound for the reduced length]\label{lem:l-function-lower-bound}
\label{lowerbdd}
\uses{thm:asymptotic-gradient-shrinking-soliton}
There is a constant $c_1>0$ depending only on the dimension $n$ such
that for any $p,q\in M_k$ we have
$$l_{x_k}(p,\tau)\ge -l_{x_k}(q,\tau)-1+c_1\frac{d^2_{g(-\tau)}(p,q)}{\tau}.$$
\end{lemma}

\begin{proof}
Since both sides of this inequality and also Ricci flow are
invariant if the metric and time are simultaneously rescaled, it
suffices to consider the case when $\tau=1$. Also, since ${\mathcal
U}_x(1)$ is a dense subset, it suffices to assume that $p,q\in
{\mathcal U}_x(1)$. Also, by symmetry, we can suppose that
$l_{x_k}(q,1)\le l_{x_k}(p,1)$.

 Let
$\gamma_1$ and $\gamma_2$ be the minimizing ${\mathcal L}$-geodesics
from $x$ to $(p,1)$ and $(q,1)$ respectively. We define a function
$f\colon M_k\times M_k\times [0,\infty)\to \Ar$ by
$$f(a,b,\tau)=d_{g_k(-\tau)}(a,b).$$
Since $\gamma_1(0)=\gamma_2(0)$ we have
\begin{eqnarray}\nonumber
d_{g_k(-1)}(p,q) & = & f(p,q,1) \\ \nonumber & =&
\int_0^1\frac{d}{d\tau}f(\gamma_1(\tau),\gamma_2(\tau),\tau)d\tau \\
\nonumber & = & \int_0^1\bigl(\frac{\partial f}{\partial
\tau}(\gamma_1(\tau),\gamma_2(\tau),\tau)+\langle \nabla
f_a,\gamma_1'(\tau)\rangle\bigr. \\
\label{fform} & & \hskip.2in+\bigl.\langle \nabla
f_b,\gamma_2'(\tau)\rangle\bigr)d\tau,
\end{eqnarray}
where $\nabla_af$ and $\nabla_bf$ refer respectively to  the
gradient of $f$ with respect to the first copy of $M_k$ in the
domain and the second copy of $M_k$ in the domain. Of course,
$|\nabla f_a|=1$ and $|\nabla f_b|=1$.

By Corollary~\ref{DL}, we have $\gamma_1'(\tau)=\nabla
l_{x_k}(\gamma_1(\tau),\tau)$ and $\gamma_2'(\tau)=\nabla
l_{x_k}(\gamma_2(\tau),\tau)$. Since $R\ge 0$ we have
$$l_{x_k}(\gamma_1(\tau),\tau)=\frac{1}{2\sqrt{\tau}}{\mathcal
L}_{x_k}(\gamma_1|_{[0,\tau]})\le \frac{1}{2\sqrt{\tau}}{\mathcal
L}_{x_k}(\gamma_1)=\frac{1}{\sqrt{\tau}}l_{x_k}(p,1).$$
Symmetrically, we have
$$l_{x_k}(\gamma_2(\tau),\tau)\le \frac{1}{\sqrt{\tau}}l_{x_k}(q,1).$$
From this inequality, Corollary~\ref{cor:l-function-gradient-bound}, and the fact that
$R\ge 0$, we have
\begin{eqnarray}\nonumber \bigl|\langle \nabla
f_a(\gamma_1(\tau),\gamma_2(\tau),\tau),\gamma_1'(\tau)\rangle
\bigr|& \le &
|\gamma_1'(\tau)| =|\nabla l_{x_k}(\gamma_1(\tau),\tau) | \nonumber \\
& \le &
\frac{\sqrt{3}}{\tau^{3/4}}\sqrt{l_{x_k}(p,1)}  \nonumber \\
& \le &
\frac{\sqrt{3}}{\tau^{3/4}}\sqrt{l_{x_k}(p,1)+1}.\label{1eqn}\end{eqnarray}
Symmetrically, we have \begin{equation}\label{2eqn} \left|\langle
\nabla
f_b(\gamma_1(\tau),\gamma_2(\tau),\tau),\gamma'_2(\tau)\rangle\right|\le
\frac{\sqrt{3}}{\tau^{3/4}}\sqrt{l_{x_k}(q,1)}\le
\frac{\sqrt{3}}{\tau^{3/4}}\sqrt{l_{x_k}(q,1)+1}.\end{equation}


 It follows from
Corollary~\ref{cor:l-function-gradient-bound} that for any $p$
$$|\nabla (\sqrt{l_{x_k}(p,\tau)})|\le
\frac{\sqrt{3}}{2\sqrt{\tau}}.$$ Set
$r_0(\tau)=\tau^{3/4}(l_{x_k}(q,1)+1)^{-1/2}$. For any $p'\in
B_{g_k}(\gamma_1(\tau),\tau,r_0(\tau))$ integrating gives
$$l_{x_k}^{1/2}(p',\tau)\le
l_{x_k}^{1/2}(\gamma_1(\tau),\tau)+\frac{\sqrt{3}}{2\sqrt{\tau}}r_0(\tau)\le
\left(\tau^{-1/4}+\frac{\sqrt{3}}{2}\tau^{1/4}\right)\sqrt{l_{x_k}(p,1)+1},$$
where in the last inequality we have used the fact that $1\le
l_{x_k}(q,1)+1\le l_{x_k}(p,1)+1$. Again using
Corollary~\ref{cor:l-function-gradient-bound} we have
$$R(p',\tau)\le
\frac{3}{\tau}\left(\tau^{-1/4}+\frac{\sqrt{3}}{2}\tau^{1/4}\right)^2(l_{x_k}(p,1)+1).$$
Now consider $q'\in B_{g_k(\tau)}(\gamma_2(\tau),\tau,r_0(\tau))$.
Similarly to the above computations, we have
$$l^{1/2}_{x_k}(q',\tau)\le
l_{x_k}^{1/2}(q,1)+\frac{\sqrt{3}}{2\sqrt{\tau}}r_0(\tau),$$ so that
$$l_{x_k}^{1/2}(q',\tau)\le \left(\tau^{-1/4}
+\frac{\sqrt{3}}{2}\tau^{1/4}\right)\sqrt{l_{x_k}(q,1)+1},$$ and
$$|\Ric(q',\tau)\le R(q',\tau)\le
\frac{3}{\tau}\left(\tau^{-1/4}+\frac{\sqrt{3}}{2}\tau^{1/4}\right)^2(l_{x_k}(q,1)+1).$$
We set
$$K=\frac{3}{\tau}\left(\tau^{-1/4}+\frac{\sqrt{3}}{2}\tau^{1/4}\right)^2(l_{x_k}(q,1)+1).$$

Now, noting that $\partial/\partial \tau$ here is
$-\partial/\partial t$ of Proposition~\ref{I.8.3}, we apply
Proposition~\ref{I.8.3} to see that
\begin{eqnarray*}
\left|\frac{\partial}{\partial
\tau}f(\gamma_1(\tau),\gamma_2(\tau),\tau)\right|
& \le & 2(n-1)\left(\frac{2}{3(n-1)}Kr_0(\tau)+r_0(\tau)^{-1}\right) \\
& \le &
\left(C_1\tau^{-3/4}+C_2\tau^{-1/4}+C_3\tau^{1/4}\right)\sqrt{l_{x_k}(q,1)+1},
\end{eqnarray*}
where $C_1,C_2,C_3$ are constants depending only on the dimension
$n$.

Now plugging Equation~(\ref{1eqn}) and~(\ref{2eqn}) and  the above
inequality into Equation~(\ref{fform}) we see that
\begin{eqnarray*}
d_{g(-1)}(p,q) & \le &
\int_0^1\left(\left(C_1\tau^{-3/4}+C_2\tau^{-1/4}+C_3\tau^{1/4}\right)\sqrt{l_{x_k}(q,1)+1}\right.
\\ & & +
\left.\sqrt{3}\tau^{-3/4}\sqrt{l_{x_k}(q,1)+1}+\sqrt{3}\tau^{-3/4}\sqrt{l_{x_k}(p,1)+1}\right)d\tau.
\end{eqnarray*}
This implies that
$$d_{g(-1)}(p,q)\le
C\left(\sqrt{l_{x_k}(q,1)+1}+\sqrt{l_{x_k}(p,1)+1}\right),$$ for
some constant depending only on the dimension. Thus, since we are
assuming that $l_{x_k}(p,1)\ge l_{x_k}(q,1)$ we have
$$d^2_{g(-1)}(p,q)\le
C^2\left(3(l_{x_k}(p,1)+1)+(l_{x_k}(q,1)+1)\right)\le
4C^2(l_{x_k}(p,1)+1+l_{x_k}(q,1)),$$ for some constant $C<\infty$
depending only on the dimension. The result now follows immediately.
\end{proof}


\begin{corollary}[Refined lower bound for the reduced length]\label{cor:l-function-lower-bound-refined}
\uses{lem:l-function-lower-bound,claim:qk-length-bound}
For  any $q'\in M$ and any $0<\tau_0\le\tau'$ we have
$$l_{x_k}(q',\tau') \ge -\frac{n}{2(\tau')^2}-\frac{\tau'}{2}-1+c_1\frac{d_{g^2_k(-\tau_0)}(q_k,q')}{\tau'},$$
where  $c_1$ is the constant from Lemma~\ref{lem:l-function-lower-bound}.
\end{corollary}

\begin{proof}
\uses{claim:qk-length-bound,lem:l-function-lower-bound}
By Claim~\ref{claim:qk-length-bound}
$$l_{x_k}(q_k,\tau')\le \frac{n}{2(\tau')^2}+\frac{n\tau'}{2}.$$
Now applying Lemma~\ref{lem:l-function-lower-bound} we see that for any $0<\tau'$ and any $q'\in
M_k$ we have \begin{eqnarray*} l_{x_k}(q',\tau') & \ge &
-\frac{n}{2(\tau')^2}-\frac{n\tau'}{2}-1+c_1\frac{d^2_{g_k(-\tau')}(q_k,q')}{\tau'}
\\
& \ge &
-\frac{n}{2(\tau')^2}-\frac{n\tau'}{2}-1+c_1\frac{d^2_{g_k(-\tau_0)}(q_k,q')}{\tau'}.
\end{eqnarray*} In the last inequality, we use the fact that the
Ricci curvature is positive so that the metric is decreasing under
the Ricci flow.
\end{proof}

Since the time slices of all the flows in question have non-negative
curvature, by Theorem~\ref{BishopGromov} the volume of the ball of
radius $s$ is at most $\omega s^{n}$ where $\omega$ is the volume of
the ball of radius one in $\Ar^n$. Since the $l_k$ converge
uniformly to $l_\infty$ on compact sets and since the metrics $h_k$
converge uniformly in the $C^\infty$-topology on compact sets to
$g_\infty$, it follows that for any $\epsilon>0$, for any
$0<\tau_0\le \tau'<\infty$ there is a radius $r$ such that for every
$k$ and any $\tau\in [\tau_0,\tau']$ the integral
$$\int_{M_\infty\setminus B_{h_k(-\tau_0)}(q_k,r)}e^{-l_k(q,\tau)}dq<\epsilon.$$
It follows by Lebesgue dominated convergence that
$$\int_{M_\infty\setminus
B_{g_\infty(-\tau_0)}(q_\infty,r)}e^{-l_\infty(q,\tau)}dq\le
\epsilon.$$


\begin{claim}[Absolute convergence of the weighted distribution]\label{claim:absolute-convergence-weighted-distribution}
\label{absconv1}
\uses{cor:l-function-lower-bound-refined}
 Fix a compact interval $[-\tau,-\tau_0]\subset
(-\infty,0)$. Let $f$ be a locally Lipschitz function that is
defined on $M_\infty\times [-\tau,-\tau_0]$ and such that there is a
constant $C$ with the property that $f(q,\tau')$ by $C$ times $\max(l_\infty(q,\tau'),1)$ . Then the distribution ${\mathcal
D}_1=fe^{-l_\infty}$ is absolutely convergent in the following
sense. For any bounded smooth function $\varphi$ defined on all of
$M_\infty\times [-\tau,-\tau_0]$ and any sequence of compactly
supported, non-negative smooth functions $\psi_k$, bounded above by
$1$ everywhere that are eventually $1$ on every compact subset, the
following limit exists and is finite:
$$\lim_{k\rightarrow\infty}{\mathcal D}_1(\varphi\psi_k).$$
 Furthermore, the limit is independent of the choice of the
$\psi_k$ with the given properties.
\end{claim}

\begin{proof}
It follows from the above discussion that there are constants $c>0$
 and a ball $B\subset M_\infty$ centered at $q_\infty$ such that on
$M_\infty\times [-\tau,-\tau_0]\setminus B\times [\tau,-\tau_0]$ the
function $l_\infty$ is greater than
$cd^2_{g_\infty(-\tau_0)}(q_\infty,\cdot)-C'$. Thus,
$fe^{-l_\infty}$ has fixed exponential decay at infinity. Since the
Riemann curvature of $M_\infty\times\{\tau'\}$ is non-negative for
every $\tau'$, the flow is distance decreasing, and there is a fixed
polynomial upper bound to the growth rate of volume at infinity.
This leads to the claimed convergence property.
\end{proof}


\begin{corollary}[Absolute convergence of the gradient and time-derivative terms]\label{cor:absolute-convergence-gradient-terms}
\label{absconv}
\uses{claim:absolute-convergence-weighted-distribution,cor:l-function-gradient-bound}
The distributions $|\nabla l_\infty|^2e^{-l_\infty}$,
$Re^{-l_\infty}$, $|(\partial l_\infty/\partial \tau)|e^{-l_\infty}$
are absolutely convergent in the sense of the above claim. \end{corollary}


\begin{proof}
\uses{cor:l-function-gradient-bound,claim:absolute-convergence-weighted-distribution}
 By
Corollary~\ref{cor:l-function-gradient-bound}, each of the Lipschitz functions $|\nabla
l_\infty|^2$, $|\partial l_\infty/\partial \tau|$ and $R$
  is at most a constant multiple of $l_\infty$.
Hence, the corollary follows from the previous claim.
\end{proof}




There is a slightly weaker statement that is true for $\triangle
e^{-l_\infty}$.

\begin{claim}[Absolute convergence of the Laplacian term]\label{claim:absolute-convergence-laplacian-term}
\label{absconv2}
\uses{claim:absolute-convergence-weighted-distribution}
Suppose that $\varphi$ and $\psi_k$ are as in Claim~\ref{claim:absolute-convergence-weighted-distribution},
but in addition $\varphi$ and all the $\psi_k$ are uniformly
Lipschitz. Then
$$\lim_{k\rightarrow\infty}\int_{M_\infty}\varphi\psi_k\triangle
e^{-l_\infty}d\vol_{g_\infty}$$ converges absolutely.
\end{claim}

\begin{proof}
This time the value of the distribution on a compactly supported
smooth function $\rho$ is given by the integral of $$-\langle
\nabla\rho,\nabla e^{-l_\infty}\rangle=\langle \nabla \rho,\nabla
l_\infty\rangle e^{-l_\infty}.$$ Since $|\nabla l_\infty|$ is less
than or equal to the maximum of $1$ and $|\nabla l_\infty|^2$, it
follows immediately, that if $|\nabla \rho|$ is bounded, then the
integral is absolutely convergent. From this the claim follows
easily.
\end{proof}


\begin{corollary}[Non-negativity of the weak pairing]\label{cor:weak-inequality-nonnegative-integral}
\label{weakneg}
\uses{claim:absolute-convergence-weighted-distribution,claim:absolute-convergence-laplacian-term,cor:weighted-distribution-nonnegative}
Fix $0<\tau_0<\tau_1<\infty$. Let $f$ be a non-negative, smooth
bounded function on $M_\infty\times [\tau_0,\tau_1]$ with (spatial)
gradient of bounded norm. Then
$$\int_{\tau_0}^{\tau_1}\int_{M_\infty\times\{-\tau\}}\left(\frac{\partial l_\infty}{\partial \tau}+|\nabla
l_\infty|^2-R+\frac{n}{2\tau}-\triangle
l_\infty\right)f\tau^{-n/2}e^{-l_\infty}d\vol_{g_\infty}d\tau\ge 0.$$
\end{corollary}

\begin{proof}
\uses{claim:absolute-convergence-weighted-distribution,claim:absolute-convergence-laplacian-term,cor:weighted-distribution-nonnegative}
For the interval $[\tau_0,\tau']$  we construct a sequence of
uniformly Lipschitz functions $\psi_k$ on $M_\infty\times
[\tau_0,\tau']$ that are non-negative, bounded above by one and
eventually one on every compact set. Let $\rho(x)$ be a smooth bump
function which is one for $x$ less than $1/4$ and is zero from $x\ge
3/4$ and is everywhere between $0$ and $1$. For any $k$ sufficiently
large let $\psi_k$ be the composition of
$\rho(d_{g_\infty(-\tau_0)}(q_\infty,\cdot)-k)$. Being compositions
of $\rho$ with Lipschitz functions with Lipschitz constant $1$, the
$\psi_k$ are a uniformly Lipschitz family of functions on
$M_\infty\times\{-\tau_0\}$. Clearly then they form a uniformly
Lipschitz family on $M_\infty\times [\tau_0,\tau']$ as required.
This allows us to define any of the above distributions on Lipschitz
functions on $M_\infty\times [\tau_0,\tau']$.

Take a family $\psi_k$ of uniformly Lipschitz functions, each
bounded between $0$ and $1$ and eventually one of every compact
subset of $M_\infty\times [\tau_0,\tau_1]$. Then the family
$f\psi_k$ is a uniformly Lipschitz family of compactly supported
functions. Hence, we can apply Claims~\ref{claim:absolute-convergence-weighted-distribution}
and~\ref{claim:absolute-convergence-laplacian-term} to establish that the integral in question is the
limit of an absolutely convergence sequence. By
Corollary~\ref{cor:weighted-distribution-nonnegative} each term in the sequence is non-positive.
\end{proof}



 \subsection{Completion of the proof of
Proposition~{\ref{prop:l-infinity-differential-equations}}}\label{sect9.2.5}



Lebesgue dominated convergence implies that the following limit
exists
$$\lim_{k\rightarrow\infty}\widetilde V_k(\tau)\equiv\widetilde V_\infty(\tau)=\int_{M_\infty\times\{-\tau\}}\tau^{-n/2}e^{-
l_\infty(q,\tau)}d\vol_{g_\infty(\tau)}.$$ By Corollary~\ref{cor:reduced-volume-limit-constant}, the function $\tau\rightarrow
\widetilde V_\infty(\tau)$ is constant. On the other hand, note that
for any $0< \tau_0<\tau_1<\infty$, we have
\begin{eqnarray*}
\widetilde V_\infty(\tau_1)-\widetilde V_\infty(\tau_0) & = &
\int_{\tau_0}^{\tau_1}\frac{d\widetilde
V_\infty}{d\tau}d\tau \\
& = & \int_{\tau_0}^{\tau_1}\int_{M_\infty}\left(\frac{\partial
l_\infty}{\partial \tau}-R+\frac{n}{2\tau}\right)
\left(\tau^{-n/2}e^{-l_\infty(q,\tau)}d\vol_{g_\infty(\tau)}\right).\end{eqnarray*} According to
Corollary~\ref{cor:absolute-convergence-gradient-terms} this is an absolutely convergent integral,
and so  this integral is zero.

\begin{claim}[The pure Laplacian integral vanishes]\label{claim:laplacian-integral-vanishes}
\uses{cor:absolute-convergence-gradient-terms}
\begin{eqnarray*}
\int_{\tau_0}^{\tau_1}\int_{M_\infty\times\{-\tau\}}\triangle
e^{-l_\infty}d\vol_{g_\infty}d\tau & =
&\int_{\tau_0}^{\tau_1}\int_{M_\infty\times\{-\tau\}}\left(|\nabla
l_\infty|^2-\triangle l_\infty\right)e^{-l_\infty}d\vol_{g_\infty}d\tau \\
& = & 0.
\end{eqnarray*}
\end{claim}


\begin{proof}
\uses{cor:absolute-convergence-gradient-terms}
Since we are integrating against the constant function $1$, this
result is clear, given the convergence result,
Corollary~\ref{cor:absolute-convergence-gradient-terms}, necessary to show that this integral is
well defined.
\end{proof}

Adding these two results together gives us the following
\begin{equation}\label{tildeD}
\int_{\tau_0}^{\tau_1}\int_{M_\infty\times\{-\tau\}}\left(\frac{\partial
l_\infty}{\partial \tau}+|\nabla l_\infty|^2
-R+\frac{n}{2\tau}-\triangle
l_\infty\right)\tau^{-n/2}e^{-l_\infty}d\vol_{g_\infty}=0.\end{equation}


Now let $\varphi$ be any compactly supported, non-negative smooth
function. By scaling by a positive constant, we can assume that
$\varphi\le 1$ everywhere. Let $\widetilde{\mathcal D}$ denote the
distribution given by
$$\widetilde{\mathcal D}(\varphi)=\int_{\tau_0}^{\tau_1}\int_{M_\infty\times\{-\tau\}}
\varphi\left(\frac{\partial l_\infty}{\partial \tau}+|\nabla
l_\infty|^2 -R+\frac{n}{2\tau}-\triangle
l_\infty\right)\tau^{-n/2}e^{-l_\infty}d\vol_{g_\infty}.$$ Then
we have seen that $\widetilde{\mathcal D}$ extends to a functional
on bounded smooth functions of bounded gradient. Furthermore,
according to Equation~(\ref{tildeD}), we have $\widetilde {\mathcal
D}(1)=0$. Thus,
$$0={\mathcal D}(1)={\mathcal D}(\varphi)+{\mathcal D}(1-\varphi).$$
Since both $\varphi$ and $1-\varphi$ are non-negative, it follows from
Corollary~\ref{cor:weak-inequality-nonnegative-integral}, that ${\mathcal D}(\varphi)$ and ${\mathcal
D}(1-\varphi)$ are each $\ge 0$. Since their sum is zero, it must be the case
that each is individually zero.

This proves that the Inequality~(\ref{1stlinftyeqn}) is actually an
equality in the weak sense, i.e., an equality of distributions on
$M_\infty\times [\tau_0,\tau')$. Taking limits we see:
\begin{equation}\label{constcon}
\widetilde{\mathcal D}=\left(\frac{\partial l_\infty}{\partial
\tau}+|\nabla l_\infty|^2-R+\frac{n}{2\tau}-\triangle
l_\infty\right)\tau^{-n/2}e^{-l_\infty}=0,\end{equation} in the weak
sense on all of $M\times(-\infty,0)$. Of course, this implies that
$$\frac{\partial l_\infty}{\partial \tau}+|\nabla
l_\infty|^2-R+\frac{n}{2\tau}-\triangle l_\infty=0$$ in the weak
sense.


It now follows by parabolic regularity that $l_\infty$ is a smooth
function on $M_\infty\times(-\infty,0)$ and that
Equation~(\ref{constcon}) holds in the usual sense.


Now from the last two equations in Corollary~\ref{lformula} and the
convergence of the $l_{x_k}$ to $l_\infty$, we conclude that the
following equation also holds:
\begin{equation}\label{wk=}2\triangle l_\infty-|\nabla
l_\infty|^2+R+\frac{l_\infty-n}{\tau}=0.\end{equation}

This completes the proof of Proposition~\ref{prop:l-infinity-differential-equations}.

\subsection{The gradient shrinking soliton equation}

Now we return to the proof of Theorem~\ref{thm:asymptotic-gradient-shrinking-soliton}. We have shown that the
limiting Ricci flow   referred to in that result exists, and we have
established that the limit $l_\infty$ of the length functions $l_{x_k}$ is a
smooth function and satisfies the differential equalities given in
Proposition~\ref{prop:l-infinity-differential-equations}. We shall use these to establish the gradient
shrinking soliton equation,
Equation~(\ref{GSSequation}), for the limit for $f=l_\infty$.

\begin{proposition}[The limit satisfies the gradient shrinking soliton equation]\label{prop:limit-gradient-shrinking-soliton-equation}
\label{gradshrink}
\uses{prop:l-infinity-differential-equations,def:gradient-shrinking-soliton}
The following equation holds on $M_\infty\times (-\infty,0)$:
$$ \Ric_{g_{\infty}(t)} + \Hess^{g_\infty(t)}(l_\infty(\cdot,\tau))
-\frac{1}{2\tau}g_{\infty}(t) = 0,$$ where $\tau=-t$,
\end{proposition}

\begin{proof}  This result will follow immediately
from:


\begin{lemma}[Conjugate heat equation identity for gradient shrinking solitons]\label{lem:conjugate-heat-equation-soliton-identity}
\label{9.1}
\uses{def:ricci-flow}
Let $(M,g(t)),\ 0\le t\le T$, be an $n$-dimensional Ricci flow, and
let $f\colon M\times [0,T]\to \Ar$ be a smooth function. As usual
set $\tau=T-t$. Then the function
$$u=(4\pi\tau)^{-\frac{n}{2}}e^{-f}$$ satisfies the conjugate heat
equation
$$-\frac{\partial u}{\partial t}-\triangle u+Ru=0,$$ if and only if we have
$$\frac{\partial f}{\partial t}+\triangle f-|\nabla
f|^2+R-\frac{n}{2\tau}=0.$$ Assuming that $u$ satisfies the
conjugate heat equation, then setting
$$v=\left[\tau\left(2\triangle f-|\nabla
f|^2+R\right)+f-n\right]u,$$ we have
$$-\frac{\partial v}{\partial t}-\triangle v+Rv=-2\tau
\bigl|\Ric_g+\Hess^g(f)-\frac{1}{2\tau}g\bigr|^2u.$$
\end{lemma}

Let us assume the lemma for a moment and use it to complete the
proof of the proposition.

We apply the lemma to the limiting Ricci flow
$(M_\infty,g_\infty(t))$ with the function $f=l_\infty$. According
to Proposition~\ref{prop:l-infinity-differential-equations} and the first statement in
Lemma~\ref{lem:conjugate-heat-equation-soliton-identity}, the function $u$ satisfies the conjugate heat
equation. Thus, according to the second statement in
Lemma~\ref{lem:conjugate-heat-equation-soliton-identity}, setting
$$v=\left[\tau\left(2\triangle f-|\nabla
f|^2+R\right)+f-n\right]u,$$ we have
$$\frac{\partial v}{\partial\tau}-\triangle v+Rv=-2\tau
\bigl|\Ric_g+\Hess(f)-\frac{1}{2\tau}g\bigr|^2u.$$ On the
other hand, the second equality in Proposition~\ref{prop:l-infinity-differential-equations} shows
that $v=0$. Since $u$ is nowhere zero, this implies that
$$\Ric_{g_\infty}+\Hess^{g_\infty}(f)-\frac{1}{2\tau}g_\infty=0.$$
This completes the proof of the proposition assuming the lemma.
\end{proof}

Now we turn to the proof of the lemma.

\begin{proof}{\em (of Lemma~\ref{lem:conjugate-heat-equation-soliton-identity})}
Direct computation shows that $$-\frac{\partial u}{\partial
t}-\triangle u+Ru= \left(-\frac{n}{2\tau}+\frac{\partial f}{\partial
t}+\triangle f-|\nabla f|^2+R\right)u.$$ From this, the first
statement of the lemma is clear.
 Let
$$H=\bigl[\tau(2\triangle f-|\nabla f|^2+R)+f-n\bigr]$$ so that $v=Hu$.
Then, of course,
$$\frac{\partial v}{\partial t}=\frac{\partial H}{\partial t}u+H\frac{\partial u}{\partial t}\ \ \ \mathrm{and}$$
$$\triangle v=\triangle H\cdot u+2\langle\nabla H,\nabla u\rangle+H\triangle u.$$
Since $u$ satisfies the conjugate heat equation, we have
$$-\frac{\partial v}{\partial t}-\triangle v+Rv=\left(-\frac{\partial H}{\partial t}-\triangle
H\right)u-2\langle\nabla H,\nabla u\rangle.$$ Differentiating the
definition of $H$ yields \begin{equation}\label{dHeqn}
\frac{\partial H}{\partial t} = -(2\triangle f-|\nabla
f|^2+R)+\frac{\partial f}{\partial t}+ \tau \left(2\frac{\partial
}{\partial t}\triangle f -\frac{\partial }{\partial t}(|\nabla
f|^2)+\frac{\partial R}{\partial t}\right)
\end{equation}


\begin{claim}[Commuting the Laplacian and the time derivative of $f$]\label{claim:laplacian-time-derivative-commutation}
\label{11.20}
$$\frac{\partial }{\partial t}\triangle
f=\triangle(\frac{\partial f}{\partial t})+2\langle \Ric,\Hess(f)\rangle.$$
\end{claim}

\begin{proof}
\uses{claim:curvature-symmetries-bianchi}
We work in local coordinates. We have $$\triangle
f=g^{ij}\nabla_i\nabla_j
f=g^{ij}(\partial_i\partial_jf-\Gamma_{ij}^k\partial_kf),$$ so that
from the Ricci flow equation we have
\begin{eqnarray*}
\frac{\partial }{\partial t}\triangle f & = & 2\Ric^{ij}\Hess(f)_{ij}+g^{ij}\frac{\partial}{\partial t}\left(\Hess(f)_{ij}\right)
\\
& = & 2\Ric^{ij}\Hess(f)_{ij}+g^{ij}\Hess\left({\frac{\partial f}{\partial
t}}\right)_{ij}-g^{ij}\frac{\partial \Gamma_{ij}^k}{\partial
t}\partial_kf.
\end{eqnarray*}
Since the first term is $2\langle \Ric,\Hess(f)\rangle$
and the second is $\triangle(\frac{\partial f}{\partial t})$, to
complete the proof of the claim, we must show that the last term of
this equation vanishes. In order to simplify the computations, we
assume that the metric is standard to second order at the point and
time under consideration. Then, using the Ricci flow equation, the
definition of the Christoffel symbols in terms of the metric, and
the fact that $g_{ij}$ is the identity matrix at the given point and
time and that its covariant derivatives in all spatial directions
vanish at this point and time, we get
$$g^{ij}\frac{\partial \Gamma_{ij}^k}{\partial t}=g^{kl}g^{ij}\left(-(\nabla_j \Ric)_{li}-(\nabla_i
 \Ric)_{lj}+(\nabla_l\Ric)_{ij}\right).$$
 This expression vanishes by the second Bianchi identity (Claim~\ref{claim:curvature-symmetries-bianchi}).
 This completes the proof of the claim.
\end{proof}

We also have
$$\frac{\partial }{\partial t}(|\nabla f|^2)=2\Ric(\nabla f,\nabla f)+2\langle
\nabla\frac{\partial f}{\partial t},\nabla f\rangle.$$ (Here $\nabla
f$ is a one-form, which explains the positive sign in the Ricci
term.)



Plugging this and Claim~\ref{claim:laplacian-time-derivative-commutation} into Equation~(\ref{dHeqn})
yields
\begin{eqnarray*}
\frac{\partial H}{\partial t} & = & -2\triangle f+|\nabla f|^2-R+\frac{\partial f}{\partial t} \\
& & + \tau \left(4\langle \Ric,\Hess(f)\rangle +2\triangle
\frac{\partial f}{\partial t} -2\Ric(\nabla f,\nabla f)
-2\langle \nabla\frac{\partial f}{\partial t},\nabla
f\rangle+\frac{\partial R}{\partial t}\right).
\end{eqnarray*}
Also,
\begin{equation*} \triangle H  =  \triangle
f+\tau\left(2\triangle^2f-\triangle( |\nabla f|^2)+\triangle
R\right).
\end{equation*}

Since $u$ satisfies the conjugate heat equation, from the first part
of the lemma we have
\begin{equation}\label{feqn}\frac{\partial f}{\partial t}=-\triangle f+|\nabla
f|^2-R+\frac{n}{2\tau}.\end{equation}
 Putting all this together and
using the Equation~(\ref{Revol}) for $\partial R/\partial t$ yields
\begin{eqnarray*}
\frac{\partial H}{\partial t}+\triangle H & = & -\triangle f+|\nabla
f|^2+\frac{\partial f}{\partial t}-R \\ & & +\tau\bigl(4\langle \Ric,\Hess(f)\rangle+2\triangle \frac{\partial f}{\partial t}
+2\triangle^2f-2\Ric(\nabla f,\nabla f)\bigr. \\
& & \bigl.-\triangle(|\nabla f|^2) -2\langle \nabla\frac{\partial
f}{\partial t},\nabla f\rangle+2\triangle
R +2|\Ric|^2\bigr) \\
& = & -\triangle f+|\nabla f|^2+\frac{\partial f}{\partial t}-R
+\tau\bigl(4\langle \Ric,\Hess(f)\rangle+2\triangle(|\nabla f|^2-R)\bigr.
\\
& & \bigl.-2\Ric(\nabla f,\nabla f)-\triangle(|\nabla
f|^2)-2\langle\nabla\frac{\partial f}{\partial t},\nabla
f\rangle+2\triangle R+2|\Ric|^2\bigr)
\\
& = & -\triangle f+|\nabla f|^2+\frac{\partial f}{\partial
t}-R+\tau\Bigl[4\langle \Ric,\Hess(f)\rangle+\triangle(|\nabla f|^2)\Bigr.
\\
& & -2\Ric(\nabla f,\nabla f)  +2\langle\nabla(\triangle
f),\nabla f\rangle -2\langle \nabla(|\nabla f|^2),\nabla
f\rangle \\
& & \Bigl. +2\langle\nabla R,\nabla f\rangle+2|\Ric|^2\Bigr].
\end{eqnarray*}
Similarly, we have
\begin{eqnarray*}
\frac{2\langle\nabla u,\nabla H\rangle}{u} & = & -2\langle\nabla f,\nabla H\rangle \\
& = & -2|\nabla f|^2-2\tau\langle\nabla f,\left(\nabla(2\triangle
f)-|\nabla f|^2+R\right)\rangle \\
& = & -2|\nabla f|^2-\tau\left(4\langle\nabla f,\nabla(\triangle
f)\rangle-2\langle\nabla f,\nabla(|\nabla f|^2)\rangle+\langle\nabla
f,\nabla R\rangle\right).
\end{eqnarray*}

Thus,
\begin{eqnarray*}
\frac{\partial H}{\partial t} +\triangle H+\frac{2\langle\nabla
u,\nabla H\rangle}{u} & = & -\triangle f-|\nabla f|^2+\frac{\partial
f}{\partial t}-R + \tau\Bigl[4\langle \Ric,\Hess(f)\rangle
\Bigr.
\\
& & +\triangle(|\nabla f|^2)-2\Ric(\nabla f,\nabla f)+2|\Ric|^2
\\
& & \Bigl.-2\langle\nabla f,\nabla(\triangle f)\rangle\Bigr].
\end{eqnarray*}

\begin{claim}[Bochner formula for $|\nabla f|^2$]\label{claim:laplacian-gradient-squared-formula}
\uses{lem:laplacian-one-form}
\begin{eqnarray*} \triangle(|\nabla f|^2) & = &
2\langle\nabla(\triangle f),\nabla f\rangle+2\Ric(\nabla
f,\nabla f)+2|\Hess(f)|^2,
\end{eqnarray*}
\end{claim}

\begin{proof}
\uses{lem:laplacian-one-form}
We have
$$\triangle(|\nabla f|^2)=\triangle \langle \nabla f,\nabla f\rangle =\triangle \langle df,df\rangle=2\langle \triangle
df,df\rangle+2\langle \nabla df,\nabla df\rangle.$$ The last term is
$|\Hess(f)|^2$. According to Lemma~\ref{lem:laplacian-one-form} we have
$\triangle df=d(\triangle f)+\Ric(\nabla f,\cdot)$. Plugging
this in gives
$$\triangle(|\nabla f|^2)=2\langle d(\triangle f),df\rangle+2\langle
\Ric(\nabla f,\cdot),df\rangle+2|\Hess(f)|^2,$$ which is
clearly another way of writing the claimed result.
\end{proof}


Using this we can simplify the above to
\begin{eqnarray*}
\frac{\partial H}{\partial t} +\triangle H+\frac{2\langle\nabla
u,\nabla H\rangle}{u} & = &-\triangle f-|\nabla f|^2+\frac{\partial
f}{\partial t}-R
\\
& & +\tau\left(4\langle \Ric,\Hess(f)\rangle +2|\Hess(f)|^2+2|\Ric|^2\right).
\end{eqnarray*}
Now using Equation~(\ref{feqn}) we have
\begin{eqnarray*}
\frac{\partial H}{\partial t} +\triangle H+\frac{2\langle\nabla
u,\nabla H\rangle}{u} & = &
-2\triangle f-2R+\frac{n}{2\tau} \\
& & +\tau\left(4\langle \Ric,\Hess(f)\rangle+2|\Ric|^2+2|\Hess(f)|^2\right) \\
& = & 2\tau\bigl(2\langle \Ric,\Hess(f)\rangle+|\Ric|^2+|\Hess(f)|^2\bigr. \\
& & -\frac{\triangle f}{\tau}-\frac{R}{\tau}+\frac{n}{4\tau^2}\bigr) \\
& = & 2\tau\bigl|\Ric+
\Hess(f)-\frac{1}{2\tau}g_\infty\bigr|^2
\end{eqnarray*}
Since
$$-\frac{\partial v}{\partial t}-\triangle v+Rv=-u\left(\frac{\partial H}{\partial t}
 +\triangle H+\frac{2\langle\nabla u,\nabla H\rangle}{u}\right),$$
 this proves the lemma.
\end{proof}



At this point, setting $f=l_\infty$, we have established all the
results claimed in Theorem~\ref{thm:asymptotic-gradient-shrinking-soliton} except for the fact that
the limit is not flat. This we establish in the next chapter.


\subsection{Completion of the proof of
Theorem~{\ref{thm:asymptotic-gradient-shrinking-soliton}}}

 To complete the
proof of Theorem~\ref{thm:asymptotic-gradient-shrinking-soliton} it remains to show that for no $t\in
(-\infty,0)$ is $(M_\infty,g_\infty(t))$ flat.

\begin{claim}[The pullback of $l_\infty$ in the flat case]\label{claim:flat-limit-length-function-form}
\uses{prop:limit-gradient-shrinking-soliton-equation}
If, for some $t\in (-\infty,0)$, the Riemannian manifold
$(M_\infty,g_\infty(t))$ is flat, then there is an isometry from
$\Ar^n$ to $(M_\infty,g_\infty(t))$  and the pullback under this
isometry of the function $l_\infty(x,\tau)$ is the function
$|x|^2/4\tau+\langle x,a\rangle+b\cdot \tau$ for some $a\in \Ar^n$
and $b\in \Ar$.
\end{claim}

\begin{proof}
\uses{lem:conjugate-heat-equation-soliton-identity,prop:limit-gradient-shrinking-soliton-equation}
We know that $f=l_\infty(\cdot,\tau)$ solves the equation given in
Lemma~\ref{lem:conjugate-heat-equation-soliton-identity} and hence by the above argument, $f$ also satisfies
the equation given in  Proposition~\ref{prop:limit-gradient-shrinking-soliton-equation}. If the limit is
flat, then the equation becomes
$$\Hess(f)=\frac{1}{2\tau}g.$$ The universal covering of $(M_\infty,g_\infty(t))$
is isometric to $\Ar^n$. Choose an identification with $\Ar^n$, and lift $f$ to
the universal cover. Call the result $\tilde f$. Then $\tilde f$ satisfies
$\Hess(\tilde f)=\frac{1}{2\tau}\tilde g$, where $\tilde g$ is the usual
Euclidean metric on $\Ar^n$. This means that $\tilde f-|x|^2/4\tau$ is an
affine linear function. Clearly, then $\tilde f$ is not invariant under any
free action of a non-trivial group, so that the universal covering in question
is the trivial cover. This completes the proof of the claim.
\end{proof}

If $(M_\infty,g_\infty(t))$ is flat for some $t\in (-\infty,0)$, then by the
above $(M_\infty,g_\infty(t))$ is isometric to $\Ar^n$. According to
Theorem~\ref{4picase} this implies that $\widetilde
V_\infty(\tau)=(4\pi)^{n/2}$. This contradicts Corollary~\ref{cor:reduced-volume-limit-constant}, and
the contradiction establishes that $(M_\infty,g_\infty(t))$ is not flat for any
$t<0$.  Together with Proposition~\ref{prop:limit-gradient-shrinking-soliton-equation}, this completes the proof of
Theorem~\ref{thm:asymptotic-gradient-shrinking-soliton}, namely of the fact that $(M_\infty,g_\infty(t)),\
-\infty<t<0$, is a non-flat, $\kappa$-non-collapsed Ricci flow with
non-negative curvature operator that satisfies the gradient shrinking soliton
equation, Equation~(\ref{GSSequation}).

To emphasize once again, we do not claim that $(M_\infty,g_\infty(t))$ is a
$\kappa$-solution, since we do not claim that each time-slice has bounded
curvature operator.

\section{Splitting results at infinity}


\subsection{Point-picking}

There is a very simple, general result about Riemannian manifolds
that we shall use in various contexts to prove that certain types of
Ricci flows split at infinity as a product with $\Ar$.


\begin{lemma}[Point-picking lemma]\label{lem:point-picking-general}
\label{4times}
Let $(M,g)$ be a Riemannian manifold  and let $p\in M$ and $r>0$ be
given. Suppose that $B(p,2r)$ has compact closure in $M$ and suppose
that $f\colon B(p,2r)\times (-2r,0]\to \Ar$ is a continuous, bounded
function with $f(p,0)>0$. Then there is a point $(q,t)\in
B(p,2r)\times (-2r,0]$ with the following properties:
\begin{enumerate}
\item[(1)] $f(q,t)\ge f(p,0)$.
\item[(2)] Setting $\alpha=f(p,0)/f(q,t)$ we have $d(p,q)\le 2r(1-\alpha)$
and $t\ge -2r(1-\alpha)$.
\item[(3)] $f(q',t')<2f(q,t)$ for all $(q',t')\in B(q,\alpha r)\times (t-\alpha r,t]$.
\end{enumerate}
\end{lemma}


\begin{proof}
Consider sequences of points
$x_0=(p,0),x_1=(p_1,t_1),\ldots,x_j=(p_j,t_j)$ in $B(p,2r)\times
(-2r,0]$ with the following properties:
\begin{enumerate}
\item[(1)] $f(x_i)\ge 2f(x_{i-1})$;
\item[(2)] Setting $r_i=rf(x_0)/f(x_{i-1})$, then $r_i\le 2^{i-1}r$, and
 we have that $$x_i\in
B(p_{i-1},r_i)\times (t_{i-1}-r_i,t_{i-1}].$$
\end{enumerate}

Of course, there is exactly one such sequence with  $j=0$: it has
$x_0=(p,0)$. Suppose we have such a sequence defined for some $j\ge
0$. If follows immediately from the properties of the sequence that
$f(p_j,t_j)\ge 2^jf(p,0)$, that
$$t_j\ge -r(1+2^{-1}+\cdots+2^{1-j}),$$
and that $r_{j+1}\le 2^{-j}r$. It also follows immediately from the
triangle inequality that $d(p,p_j)\le r(1+2^{-1}+\cdots +2^{1-j})$.
 This means that
 $$B(p_j,r_{j+1})\times(t_j-r_{j+1},t_j]\subset B(p,2r)\times (-2r,0].$$
  Either the point
$x_j$ satisfies the conclusion of the lemma, or we can find
$x_{j+1}\in B(p_j,r_{j+1})\times (t_j-r_{j+1},t_j]$ with
$f(x_{j+1})\ge 2f(x_j)$. In the latter case we extend our sequence
by one term. This shows that either the process terminates at some
$j$, in which case $x_j$ satisfies the conclusion of the lemma, or
it continues indefinitely. But it cannot continue indefinitely since
$f$ is bounded on $B(p,2r)\times (-2r,0]$.
\end{proof}

One special case worth stating separately is when $f$ is independent
of $t$.


\begin{corollary}[Time-independent point-picking]\label{cor:point-picking-time-independent}
\label{pp}
\uses{lem:point-picking-general}
Let $(M,g)$ be a Riemannian manifold  and let $p\in M$ and $r>0$ be given.
Suppose that $B(p,2r)$ has compact closure in $M$ and suppose that $f\colon
B(p,2r)\to \Ar$ is a continuous, bounded function with $f(p)>0$. Then there is
a point $q\in B(x,2r)$ with the following properties:
\begin{enumerate}
\item[(1)] $f(q)\ge f(p)$.
\item[(2)] Setting $\alpha=f(p)/f(q)$ we have $d(p,q)\le 2r(1-\alpha)$
and $f(q')<2f(q)$ for all $q'\in B(q,\alpha r)$.
\end{enumerate}
\end{corollary}

\begin{proof}
\uses{lem:point-picking-general}
Apply the previous lemma to $\widehat f\colon B(p,2r)\times
(-2r,0]\to \Ar$ defined by $\widehat f(p,t)=f(p)$.
\end{proof}


\subsection{Splitting results}

Here we prove a splitting result for ancient solutions
of non-negative curvature. They are both based on Theorem~\ref{topsplit}.



\begin{proposition}[Splitting at infinity]\label{prop:splitting-at-infinity}
\label{prodatinf}
\uses{def:kappa-noncollapsed}
Suppose that $(M,g(t)),\ -\infty <t< 0$, is a $\kappa$-non-collapsed
 Ricci flow of dimension\footnote{This result in fact
 holds in all dimensions.} $n\le 3$. Suppose that $(M,g(t))$ is a
 complete, non-compact, non-flat Riemannian manifold
with non-negative curvature operator for each $t$. Suppose  that
$\partial R(q,t)/\partial t\ge 0$ for all $q\in M$ and all $t<0$.
Fix $p\in M$. Suppose that there is a sequence of points $p_i\in M$
going to infinity with the property that
$$\lim_{i\rightarrow \infty}R(p_i,-1)d^2_{g(-1)}(p,p_i)=\infty.$$
  Then
there is a sequence of points $q_i\in M$ tending to infinity such
that, setting $Q_i=R(q_i,-1)$, we have $\lim_{i\rightarrow\infty}d^2(p,q_i)Q_i=\infty$. Furthermore, setting
$g_i(t)=Q_ig(Q_i^{-1}(t+1)-1)$,
 the
sequence of  based flows $(M,g_i(t),(q_i,-1)),\ -\infty<t\le -1$,
converges smoothly to $(N^{n-1},h(t))\times (\Ar,ds^2)$, a product
Ricci flow defined for $-\infty<t\le -1$ with $(N^{n-1},h(-1))$
being non-flat and of bounded, non-negative curvature.
\end{proposition}


\begin{corollary}[No two-dimensional example]\label{cor:no-two-dimensional-splitting-example}
\label{compact2D}
\uses{prop:splitting-at-infinity}
There is no two-dimensional flow satisfying the hypotheses of
Proposition~\ref{prop:splitting-at-infinity}.
\end{corollary}

\begin{proof} (of Proposition~\ref{prop:splitting-at-infinity})
\uses{cor:point-picking-time-independent}
Take a sequence $p_i\in M$ such that
$$d_{g(-1)}^2(p,p_i)R(p_i,-1)\rightarrow\infty$$ as $i\rightarrow \infty$.
  We set $d_i=d_{g(-1)}(p,p_i)$ and we set
$B_i=B(p_i,-1,d_i/2)$, and we let $f\colon B_i\to \Ar$ be the square
root of the scalar curvature. Since $(M,g(-1))$ is complete, $B_i$
has compact closure in $M$, and consequently $f$ is a bounded
continuous function on $B_i$. Applying Corollary~\ref{cor:point-picking-time-independent} to
$(B_i,g(-1))$ and $f$, we conclude that there is a point $q_i\in
B_i$ with the following properties:
\begin{enumerate}
\item[(1)] $R(q_i,-1)\ge R(p_i,-1)$
\item[(2)] $B'_i=B(q_i,-1,(d_iR(p_i,-1)^{1/2})/(4R(q_i,t_i)^{1/2})\subset B(p_i,-1,d_i/2)$.
\item[(3)] $R(q',-1)\le 4R(q_i,-1)$ for
all $(q',-1)\in B'_i$.
\end{enumerate}
Since $d_{g(-1)}(p,q_i)\ge d_i/2$, it is also the case that
$d^2_{g(-1)}(p,q_i)R(q_i,-1)$ tends to infinity as $i$ tends to infinity.
Because of our assumption on the time derivative of $R$, it follows that
$R(q',t)\le 4R(q_i,-1)$ for all $q'\in B'_i$ and for all $t\le -1$.




 Set
$Q_i=R(q_i,-1)$. Let $M_i=M$, and set $x_i=(q_i,-1)$. Lastly, set
$g_i(t)=Q_ig(Q_i^{-1}(t+1)-1)$. We consider the based Ricci flows
$(M_i,g_i(t),x_i),\ -\infty<t\le -1$. We see that $R_{g_i}(q',t)\le
4$ for all $(q',t)\in B_{g_i}(q_i,-1,d_iR(p_i,-1)^{1/2}/4)\times
(-\infty,-1]$. Since the original Ricci flows are
$\kappa$-non-collapsed, the same is true for the rescaled flows.
Since $d_iR(p_i,-1)^{1/2}/4\rightarrow \infty$, by
Theorem~\ref{flowlimit} there is a geometric limit flow
$(M_\infty,g_\infty(t),(q_\infty,-1))$ defined for $t\in
(-\infty,-1]$. Of course, by taking limits we see that
$(M_\infty,g_\infty(t))$ is $\kappa$-non-collapsed, its scalar
curvature is bounded above by $4$, and its curvature operator is
non-negative. It follows that $(M_\infty,g_\infty(t))$ has bounded
curvature.


To complete the proof we show that the Ricci flow $(M_\infty,g_\infty(t))$
splits as a product of a line with a Ricci flow of one lower dimension. By
construction $(M_\infty,g_\infty(-1))$ is the geometric limit constructed from
$(M,g(-1))$ in the following manner. We have a sequence of points $q_i$ tending
to infinity in $M$ and constants $\lambda_i=R(q_i,-1)$ with the property that
$\lambda_id^2_{g(-1)}(p,q_i)$ tending to infinity such that
$(M_\infty,g_\infty(-1))$ is the geometric limit of $(M,\lambda_ig(-1),q_i)$.
 Thus, according
to Theorem~\ref{topsplit}, the limit $(M_\infty,g_\infty(-1))$ splits as a
Riemannian product with a line. If the dimension of $M_\infty$ is two, then
this is a contradiction: We have that $(M_\infty,g_\infty(-1))$ splits as the
Riemannian product of a line and a one-manifold and hence is flat, but
$R(q_\infty,-1)=1$. Suppose that the dimension of $M_\infty$ is three. Once we
know that $(M_\infty,g_\infty(-1))$ splits as a product with a line, it follows
from the maximum principle (Corollary~\ref{localprod}) that the entire flow
splits as a product with a line, and the Ricci flow on the surface has strictly
positive curvature.
\end{proof}


\section{Classification of gradient shrinking solitons in dimensions $2$ and $ 3$}

In this section we fix $\kappa>0$ and we classify all
$\kappa$-solutions $(M,g_\infty(t)),\ -\infty<t<0$, that satisfy the
gradient shrinking soliton equation at the time-slice $t=-1$ in the
sense that there is a function $f\colon M \to \Ar$ such that
\begin{eqnarray}\label{GSS-1}
\Ric_{g_\infty(-1)}+\Hess^{g_\infty(-1)}(f)-\frac{1}{2}
g_\infty(-1)=0.\end{eqnarray} This will give a classification of the
two- and three-dimensional asymptotic gradient shrinking solitons
constructed in Theorem~\ref{thm:asymptotic-gradient-shrinking-soliton}.


 Let us give some examples in
dimensions two and three of ancient solutions that have such
functions. It turns out, as we shall see below, that in dimensions
two and three the only such are compact manifolds of constant
positive curvature -- i.e., Riemannian manifolds finitely covered by
the round sphere. We can create another, non-flat gradient shrinking
soliton in dimension three by taking $(M,g_{-1})$ equal to the
product of $(S^2,h_{-1})$, the round sphere of Gaussian curvature
$1/2$, with the real line (with the metric on the real line denoted
$ds^{ 2}$) and setting $g(t)=|t|h_{-1}+ ds^2$ for all $t<0$. We
define $f\colon M\times (-\infty,0)\to\Ar$ by $f(p,t)=s^2/4|t|$
where $s\colon M\to \Ar$ is the projection onto the second factor.
Then it is easy to see that
$$\Ric_{g(t)}+\Hess^{g(t)}(f)-\frac{1}{2|t|}g(t)=0,$$
so that this example is a gradient shrinking soliton. There is a
free, orientation-preserving involution on this Ricci flow: the
product of the sign change on $\Ar$ with the antipodal map on $S^2$.
This preserves the family of metrics and hence there is an induced
Ricci flow on the quotient. Since this involution also preserves the
function $f$, the quotient  is also a gradient shrinking soliton.
These are the basic $3$-dimensional examples. As the following
theorem shows, they are all the $\kappa$-non-collapsed gradient
shrinking solitons in dimension three.

First we need a definition for a single Riemannian manifold analogous to a
definition we have already made for Ricci flows.

\begin{definition}[$\kappa$-non-collapsed Riemannian manifold]\label{def:kappa-noncollapsed-riemannian-manifold}
\uses{def:kappa-noncollapsed}
Let $(M,g)$ be an $n$-dimensional complete Riemannian manifold and
fix $\kappa>0$. We say that $(M,g)$ is {\em $\kappa$-non-collapsed}
if for every $p\in M$ and any $r>0$, if $|\Rm_g|\le r^{-2}$ on
$B(p,r)$ then $\Vol\, B(p,r)\ge \kappa r^n$.
\end{definition}


Here is the theorem that we shall prove:

\begin{theorem}[Classification of gradient shrinking solitons in dimension 2 and 3]\label{thm:gradient-shrinking-soliton-classification}
\label{GSSclass}
\uses{def:kappa-noncollapsed-riemannian-manifold,def:gradient-shrinking-soliton}
Let $(M,g)$ be a  complete, non-flat Riemannian manifold of bounded
non-negative curvature of dimension $2$ or $3$. Suppose that the
Riemannian manifold $(M,g)$ is $\kappa$-non-collapsed. Lastly,
suppose that there is a $C^2$-function $f\colon M \to \Ar$ such that
$$\Ric_{g}+\Hess^{g}(f)=\frac{1}{2}g.$$
 Then there is a Ricci flow $(M,G(t)),\ -\infty<t<0$, with $G(-1)=g$ and with $(M,G(t))$ isometric
 to $(M,|t|g)$ for every $t<0$. In addition, $(M,G(t))$ is of one
 of the following three types:
\begin{enumerate}
\item[(1)] The flow $(M,G(t)),\ -\infty<t<0$, is a shrinking family of
compact, round (constant positive curvature) manifolds.
\item[(2)] The flow $(M,G(t)),\ -\infty<t<0$, is a product of a shrinking family
of round $2$-spheres with the real line. \item[(3)] $(M,G(t))$ is
isomorphic to the quotient family of metrics of the product of a
shrinking family of round $2$-spheres and the real line under the
action of an isometric involution.
\end{enumerate}
\end{theorem}



Now let us begin the proof of Theorem~\ref{thm:gradient-shrinking-soliton-classification}


\subsection{Integrating $\nabla f$}

Since the curvature of $(M,g)$ is bounded, it follows immediately from the
gradient shrinking soliton equation that $\Hess^g(f)$ is bounded. Fix a
point $p\in M$. For any $q\in M$ let $\gamma(s)$ be a minimal geodesic from $p$
to $q$ parameterized at unit length. Since
$$\frac{d}{ds}\left(|\nabla f(\gamma(s))|\right)^2=2\langle \Hess(f)(\gamma'(s),\nabla f(\gamma(s))\rangle,$$ it follows that
$$\frac{d}{ds}\left(|\nabla f(\gamma(s))|\right)\le C,$$
where $C$ is an upper bound for $|\Hess(f)|$.
 By integrating, it
follows  that
$$|\nabla f(q)|\le Cd_g(p,q)+|\nabla f(p)|.$$
This means that any flow line $\lambda(t)$ for $\nabla f$ satisfies
$$\frac{d}{dt}d_g(p,\lambda(t))\le Cd_g(p,\lambda (t))+|\nabla f(p)|,$$
and hence these flow lines do not escape to infinity in finite time.
It follows that there is a flow $\Phi_t\colon M\to M$ defined for
all time with $\Phi_0=\Id$ and $\partial\Phi_t/\partial
t=\nabla f$. We consider the one-parameter family of diffeomorphisms
$\Phi_{-\log(|t|)}\colon M\to M$ and define
\begin{equation}\label{Gdefn} G(t)=|t|\Phi_{-\log(|t|)}^*g.\ \
\ -\infty<t<0.\end{equation} We compute
$$\frac{\partial G}{\partial t}=-\Phi_{h(t)}^*g+2\Phi_{h(t)}^*\Hess^g(f)=-2\Phi_{h(t)}^*\Ric(g)=-2\Ric(G(t)),$$ so that
$G(t)$ is a Ricci flow. Clearly, every time-slice is a complete,
non-flat manifold of non-negative bounded curvature. It is clear
from the construction that $G(-1)=g$ and that $(M,G(t))$ is
isometric to $(M,|t|g)$. This shows that $(M,g)$ is the $-1$
time-slice of a Ricci flow $(M,G(t))$ defined for all $t<0$, and
that, furthermore, all the manifolds $(M,G(t))$ are equivalent up to
diffeomorphism and scaling by $|t|$.


\subsection{Case 1: $M$ is compact and the curvature is strictly
positive}\label{compactcs}


\begin{claim}[Compact positive-curvature case is round]\label{claim:compact-positive-curvature-round}
\uses{thm:gradient-shrinking-soliton-classification}
Suppose that $(M,g)$ and $f\colon M\to \Ar$ satisfies the hypotheses of
Theorem~\ref{thm:gradient-shrinking-soliton-classification} and that $M$ is compact and of positive curvature. Then
the Ricci flow $(M,G(t))$ with $G(-1)=g$ given in Equation~(\ref{Gdefn}) is a
shrinking family of compact round manifolds.
\end{claim}

\begin{proof} The
manifold $(M,G(t))$ given in Equation~(\ref{Gdefn}) is equivalent up to
diffeomorphism and scaling by $|t|$ to $(M,g)$. If the dimension of $M$ is
three, then according to Hamilton's pinching toward positive curvature result
(Theorem~\ref{flowtoround}), the Ricci flow becomes singular in finite time and
as it becomes singular the metric approaches constant curvature in the sense
that the ratio of the largest sectional curvature to the smallest goes to one.
But this ratio is invariant under scaling and diffeomorphism, so that it must
be the case that for each $t$, all the sectional curvatures of the metric
$G(t)$ are equal; i.e., for each $t$ the metric $G(t)$ is round. If the
dimension of $M$ is two, then the results go back to Hamilton in
\cite{Hamiltonsurface}. According to Proposition 5.21 on p.. 118 of
\cite{ChowKnopf}, $M$ is a shrinking family of constant positive curvature
surfaces, which must be either $S^2$ or $\Ar P^2$. This completes the analysis
in the compact case.
\end{proof}


From this result, we can easily deduce a complete classification of
$\kappa$-solutions with compact asymptotic gradient shrinking
soliton.



\begin{corollary}[Compact asymptotic soliton forces time-shifted equivalence]\label{cor:compact-asymptotic-soliton-classification}
\label{compactcase}
\uses{claim:compact-positive-curvature-round}
Suppose that $(M,g(t))$ is a $\kappa$-solution of dimension $3$ with a compact
asymptotic gradient shrinking soliton. Then the Ricci flow $(M,g(t))$ is
isomorphic to a time-shifted version of its asymptotic gradient shrinking
soliton.
\end{corollary}

\begin{proof}
\uses{claim:compact-positive-curvature-round}
We suppose that the compact asymptotic gradient shrinking soliton is
the limit of the $(M,g_{\tau_n}(t),(q_n,-1))$ for some sequence of
$\tau_n\rightarrow\infty$. Since by the discussion in the compact
case, this limit is of constant positive curvature. It follows that
for all $n$ sufficiently large, $M$ is diffeomorphic to the limit
manifold and the metric $g_{\tau_n}(-1)$ is close to a metric of
constant positive curvature.  In particular, for all $n$
sufficiently large, $(M,g_{\tau_n}(-1))$ is compact and of strictly
positive curvature. Furthermore, as $n\rightarrow\infty$
$\tau_n\rightarrow\infty$ and Riemannian manifolds
$(M,g_{\tau_n}(-1))$ become closer and closer to round in the sense
that the ratio of its largest sectional curvature to its smallest
sectional curvature goes to one. Since this is a scale invariant
ratio, the same is true for the sequence of Riemannian manifolds
$(M,g(-\tau_n))$. In the case when the dimension of $M$ is three, by
Hamilton's pinching toward round result or Ivey's theorem (see
Theorem~\ref{flowtoround}), this implies that the $(M,g(t))$ are all
exactly round.



This proves that $(M,g(t))$ is a shrinking family of round metrics.
The only invariants of such a family are the diffeomorphism type of
$M$ and the time $\Omega$ at which the flow becomes singular. Of
course, $M$ is diffeomorphic to its asymptotic soliton. Hence, the
only remaining invariant is the singular time, and hence $(M,g(t))$
is equivalent to a time-shifted version of its asymptotic soliton.
\end{proof}

\subsection{Case 2: Non-strictly positively curved.}


\begin{claim}[Non-strictly positive curvature splits off a line]\label{claim:non-strictly-positive-curvature-splitting}
\label{flatssplit}
\uses{thm:gradient-shrinking-soliton-classification}
Suppose that $(M,g)$ and $f\colon M\to \Ar$ are as in the statement
of Theorem~\ref{thm:gradient-shrinking-soliton-classification} and that $(M,g)$ does not have strictly
positive curvature. Then $n=3$ and the Ricci flow $(M,G(t))$ with
$G(-1)=g$ given in Equation~(\ref{Gdefn}) has a one- or two-sheeted
covering that is a product of a two-dimensional
$\kappa$-non-collapsed Ricci flow of positive curvature and a
constant flat copy of $\Ar$. The curvature is bounded on each
time-slice.
\end{claim}


\begin{proof}
According to Hamilton's strong maximum principle
(Corollary~\ref{nullspace}), the Ricci flow $(M,G(t))$ has a one- or
two-sheeted covering that splits as a product of an evolving family
of manifolds of one dimension less of positive curvature and a
constant one-manifold.  It follows immediately that $n=3$.
 Let
$\tilde f$ be the lifting of $f$ to this one- or two-sheeted
covering. Let $Y$ be a unit tangent vector in the direction of the
one-manifold. Then it follows from Equation~(\ref{GSS-1}) that the
value of the Hessian of $\tilde f$ of $(Y,Y)$ is one. If the flat
one-manifold factor is a circle then there can be no such function
$\tilde f$. Hence, it follows that the one- or two-sheeted covering
is a product of an evolving surface with a constant copy of $\Ar$.
 Since
$(M,g)$ is $\kappa$-non-collapsed and of bounded curvature,
$(M,G(t))$ is $\kappa$-non-collapsed and each time-slice has
positive bounded curvature. These statements are also true for the
flow of surfaces.
\end{proof}




\subsection{Case 3: $M$ is non-compact and
strictly positively curved}\label{noncomp}

Here the main result is that this case does not occur.

\begin{proposition}[No non-compact positively curved examples]\label{prop:no-noncompact-positive-curvature-soliton}
\label{noncompsol}
\uses{thm:gradient-shrinking-soliton-classification}
There is no  two- or three-dimensional  Ricci flow satisfying the hypotheses of
Theorem~\ref{thm:gradient-shrinking-soliton-classification} with $(M,g)$ non-compact and of positive curvature.
\end{proposition}

We suppose that we have $(M,g)$ as in Theorem~\ref{thm:gradient-shrinking-soliton-classification} with
$(M,g)$ being non-compact and of positive curvature. Let $n$ be the
dimension of $M$, so that $n$ is either $2$ or $3$. Taking the trace
of the gradient shrinking soliton equation yields
$$R+\triangle f-\frac{n}{2}=0,$$
and consequently that $$ dR+d(\triangle f)=0.$$ Using
Lemma~\ref{lem:laplacian-one-form} we rewrite this equation as
\begin{equation}\label{Rfeqn}
dR+\triangle(df)-\Ric(\nabla f,\cdot)=0.\end{equation}
 On the other
hand, taking the divergence of the gradient shrinking soliton
equation and using the fact that $\nabla^* g=0$ gives
$$\nabla^*\Ric+\nabla^*\Hess(f)=0.$$
Of course, $$\nabla^*\Hess(f)=\nabla^*(\nabla\nabla
f)=(\nabla^*\nabla)\nabla f=\triangle (df),$$ so that
$$\triangle(df)=-\nabla^*\Ric.$$ Plugging this into Equation~\ref{Rfeqn}
 gives
$$dR-\nabla^*\Ric-\Ric(\nabla f,\cdot)=0.$$
Now invoking Lemma~\ref{lem:div-ricci-scalar-curvature} we have  \begin{equation}\label{dR}
dR=2\Ric(\nabla f,\cdot).\end{equation}

 Fix a point $p\in M$. Let
$\gamma(s);\ 0\le s\le \bar s$, be a shortest geodesic (with respect to the
metric $g$), parameterized at unit speed, emanating from $p$, and set
$X(s)=\gamma'(s)$.

\begin{claim}[Bounded Ricci integral along minimizing rays]\label{claim:ricci-integral-bound}
There is a constant $C$ independent of the choice of
$\gamma$ and of $\bar s$ such that
$$\int_0^{\bar s}\Ric(X,X)ds\le C.$$
\end{claim}

\begin{proof}
Since the curvature is bounded, clearly it suffices to assume that
$\bar s>>1$. Since $\gamma$ is length-minimizing and parameterized
at unit speed, it follows that it is a local minimum for the energy
functional $E(\gamma)=\frac{1}{2}\int_0^{\bar s}|\gamma'(s)|^2ds$
among all paths with the same end points. Thus, letting
$\gamma_u(s)=\gamma(s,u)$ be a one-parameter family of variations
(fixed at the endpoints) with $\gamma_0=\gamma$ and with
$d\gamma/du|_{u=0}=Y$, we see
$$0\le \delta_{Y}^2E(\gamma_u)=\int_0^{\bar
s}|\nabla_XY|^2+\langle {\mathcal R}(Y,X)Y,X\rangle ds.$$ We
conclude that
\begin{equation}\label{Yeqn}
\int_0^{\bar s}\langle -{\mathcal R}(Y,X)Y,X\rangle ds\le
\int_0^{\bar s}|\nabla_XY|^2ds.
\end{equation} Fix an orthonormal basis $\{E_i\}_{i=1}^n$  at $p$ with $E_n=X$, and let
$\widetilde E_i$ denote the parallel translation of $E_i$ along $\gamma$. (Of
course, $\widetilde E_n=X$.) Then, for $i\le n-1$, we define
$$Y_i=\left\{\begin{array}{ll} s\widetilde E_i & \mbox{if $0\le s\le
1$} \\ \widetilde E_i & \mbox{if $1\le s\le \bar s-1$} \\
(\bar s-s)E_i & \mbox{if $\bar s-1\le s\le \bar s$.} \end{array}
\right. $$ Adding up Equation~(\ref{Yeqn}) for each $i$ gives
$$-\sum_{i=1}^{n-1}\int_0^{\bar s}\langle{\mathcal R}(Y_i,X)Y_i,X\rangle ds\le
\sum_{i=1}^{n-1}\int_0^{\bar s}|\nabla_XY_i|^2ds.$$ Of course, since the
$\widetilde E_i$ are parallel along $\gamma$, we have
$$|\nabla_XY_i|^2=\left\{\begin{array}{ll} 1 & \mbox{if $0\le s\le 1$}
\\ 0 & \mbox{if $1\le s\le \bar s-1$} \\ 1 & \mbox{if $\bar s-1\le s\le
\bar s$}\end{array}\right. ,$$ so that
$$\sum_{i=1}^{n-1}\int_0^{\bar s}|\nabla_XY_i|^2=2(n-1).$$
On the other hand,
$$-\sum_{i=1}^{n-1}\langle{\mathcal R}(Y_i,X)(Y_i),X\rangle =\left\{
\begin{array}{ll} s^2\Ric(X,X) & \mbox{if $0\le s\le 1$} \\
\Ric(X,X) & \mbox{if $1\le s\le \bar s-1$} \\ (\bar s-s)^2
\Ric(X,X) & \mbox{if $\bar s-1\le s\le \bar s$.} \end{array} \right.
$$ Since the curvature is bounded and $|X|=1$, we see that
$\int_0^{\bar s}(1-s^2)\Ric(X,X)ds+\int_{\bar s-1}^{\bar
s}(\bar s-s)^2\Ric(X,X)$ is bounded independent of $\gamma$ and
of $\bar s$. This concludes the proof of the claim.
\end{proof}

\begin{claim}[Cauchy--Schwarz bound for the Ricci tensor]\label{claim:ricci-cauchy-schwarz-bound}
$|\Ric(X,\cdot)|^2\le R \cdot \Ric(X,X)$.
\end{claim}


\begin{proof}
This is obvious if $n=2$, so we may as well assume that $n=3$. We
diagonalize $\Ric$ in an orthonormal basis $\{e_i\}$. Let
$\lambda_i\ge 0$
 be the eigenvalues. Write $X=X^ie_i$ with $\sum_i(X^i)^2=1$. Then
$$\Ric(X,\cdot)=X^i\lambda_i(e_i)^*,$$ so that
$|\Ric(X,\cdot)|^2=\sum_i(X^i)^2\lambda_i^2$. Of course, since
the $\lambda_i\ge 0$, this gives
$$R\cdot \Ric(X,X)=(\sum_i\lambda_i)\sum_i\lambda_i(X^i)^2\ge
\sum_i\lambda_i^2(X^i)^2,$$ establishing the claim.
\end{proof}


Now we compute, using Cauchy-Schwarz,
\begin{eqnarray*}
\left(\int_0^{\bar s}|\Ric(X,\widetilde E_i)|ds\right)^2 & \le
&\bar s\int_0^{\bar s}|\Ric(X,\widetilde E_i)|^2ds\le \bar
s\int_0^{\bar s}|\Ric(x,\cdot)|^2ds\\ & \le &  \bar
s\int_0^{\bar s}R\cdot \Ric(X,X)ds.\end{eqnarray*} Since $R$ is
bounded, it follows from the first claim that there is a constant
$C'$ independent of $\gamma$ and $\bar s$ with
\begin{equation}\label{sqrts}
\int_0^{\bar s}|\Ric(X,\widetilde E_i)|ds\le C'\sqrt{\bar
s}.\end{equation} Since $\gamma$ is a geodesic in the metric $g$, we have
$\nabla_XX=0$. Hence,
$$\frac{d^2f(\gamma(s))}{ds^2}=X(X(f))=\Hess(f)(X,X).$$
Applying  the gradient shrinking soliton equation to the pair
$(X,X)$ gives
$$\frac{d^2f(\gamma(s))}{ds^2}=\frac{1}{2}-\Ric_{g}(X,X).$$
Integrating we see
$$\frac{df(\gamma(s))}{ds}|_{s=\bar
s}=\frac{df(\gamma(s))}{ds}|_{s=0}+\frac{\bar s}{2}-\int_0^{\bar
s}\Ric(X,X)ds.$$ It follows that
\begin{equation}\label{X}
X(f)(\gamma(\bar s))\ge \frac{\bar s}{2}-C'',\end{equation} for some
constant $C''$ depending only on $(M,g)$ and $f$. Similarly,
applying the gradient shrinking soliton equation to the pair
$(X,\widetilde E_i)$, using Equation~(\ref{sqrts}) and the fact that
$\nabla_X\widetilde E_i=0$ gives
\begin{equation}\label{sqrts2}
|\widetilde E_i(f)(\gamma(\bar s))|\le C''(\sqrt{\bar
s}+1).\end{equation}

These two inequalities imply that for $\bar s$ sufficiently large,
$f$ has no critical points  and that $\nabla f$ makes a small angle
with the gradient of the distance function from $p$, and $|\nabla
f|$ goes to infinity as the distance from $p$ increases. In
particular, $f$ is a proper function going off to $+\infty$ as we
approach infinity in $M$.

Now apply Equation~(\ref{dR}) to see that $R$ is increasing along the gradient
curves of $f$. Hence, there is a sequence $p_k$ tending to infinity in $M$ with
$\lim_kR(p_k)=\limsup_{q\in M} R_g(q)>0$.


The Ricci flow $(M,G(t)),\ -\infty<t<0$, given in
Equation~(\ref{Gdefn}) has the property that $G(-1)=g$ and that
$(M,G(t))$ is isometric to $(M,|t|g)$. Since the original Riemannian
manifold $(M,g)$ given in the statement of Theorem~\ref{thm:gradient-shrinking-soliton-classification} is
$\kappa$-non-collapsed, it follows that, for every $t<0$, the
Riemannian manifold $(M,G(t))$ is $\kappa$-non-collapsed.
Consequently, the Ricci flow $(M,G(t))$ is $\kappa$-non-collapsed.
It clearly has bounded non-negative curvature on each time-slice and
is non-flat.  Fix a point $p\in M$. There is a sequence of points
$p_i$ tending to infinity with $R(p_i,-1)$ bounded away from zero.
It follows  that $\lim_{i\rightarrow
\infty}R(p_i,-1)d_{g(-1)}^2(p,p_i)=\infty$. Thus, this flow
satisfies all the hypotheses of Proposition~\ref{prop:splitting-at-infinity}. Hence,
by Corollary~\ref{cor:no-two-dimensional-splitting-example} we see that $n$ cannot be equal to two.
Furthermore, by Proposition~\ref{prop:splitting-at-infinity}, when $n=3$ there is
another subsequence $q_i$ tending to infinity in $M$ such that there
is a geometric limit $(M_\infty,g_\infty(t),(q_\infty,-1)), \
-\infty<t\le -1$, of the flows $(M,G(t),(q_i,-1))$ defined for all
$t<0$ and this limit splits as a product of a surface flow
$(\Sigma^2,h(t))$ times the real line where the surfaces
$(\Sigma^2,h(t))$ are all of positive, bounded curvature and the
surface flow is $\kappa$-non-collapsed. Since there is a constant
$C<\infty$ such that the curvature of $(M,G(t)), -\infty<t\le
t_0<0$, is bounded by $C/|t_0|$, this limit actually exists for
$-\infty<t<0$ with the same properties.


Let us summarize our progress to date.

\begin{corollary}[Summary: non-compact soliton splitting]\label{cor:noncompact-soliton-splitting-summary}
\label{compact2Dsol}
\uses{claim:non-strictly-positive-curvature-splitting,prop:splitting-at-infinity}
There is no non-compact, two-dimensional Riemannian manifold $(M,g)$
satisfying the hypotheses of Theorem~\ref{thm:gradient-shrinking-soliton-classification}. For any
non-compact three-manifold $(M,g)$ of positive curvature satisfying
the hypotheses of Theorem~\ref{thm:gradient-shrinking-soliton-classification}, there is a sequence of
points $q_i\in M$ tending to infinity such that $\lim_{i\rightarrow \infty}R_g(q_i)=\sup_{p\in M}$ such that the
based Ricci flows $(M,G(t),(q_i,-1))$ converge to a Ricci flow
$(M_\infty,G_\infty(t),(q_\infty,-1))$ defined for $-\infty<t<0$
that splits as a product of a line and a family of surfaces, each of
positive, bounded curvature $(\Sigma^2,h(t))$. Furthermore, the flow
of surfaces is $\kappa$-non-collapsed.
\end{corollary}

\begin{proof}
In Claim~\ref{claim:non-strictly-positive-curvature-splitting} we saw that every two-dimensional $(M,g)$ satisfying
the hypotheses of Theorem~\ref{thm:gradient-shrinking-soliton-classification} has strictly positive curvature. The
argument that we just completed shows that there is no non-compact
two-dimensional example of strictly positive curvature.

The final statement  is exactly what we just established.
\end{proof}


\begin{corollary}[Classification of two-dimensional gradient shrinking solitons]\label{cor:two-dimensional-gradient-shrinking-soliton}
\label{2DGSS}
\uses{cor:no-two-dimensional-splitting-example,cor:noncompact-soliton-splitting-summary}
\begin{enumerate}
\item[(1)] Let $(M,g(t))$ be a two-dimensional Ricci flow satisfying all
the hypotheses of Proposition~\ref{prop:splitting-at-infinity} except possible the
non-compactness hypothesis. Then $M$ is compact and for any $a>0$
the restriction of the flow to any interval of the form
$(-\infty,-a]$ followed by a shift of time by $+a$ is a
$\kappa$-solution.
\item[(2)] Any asymptotic gradient
shrinking soliton for a two-dimensional $\kappa$-solution is a shrinking family
of round surfaces.
\item[(3)] Let $(M,g(t)), -\infty<t\le 0$, be a
two-dimensional $\kappa$-solution. Then $(M,g(t))$ is a shrinking family of
compact, round surfaces.
\end{enumerate}
\end{corollary}

\begin{proof}
\uses{cor:no-two-dimensional-splitting-example,cor:noncompact-soliton-splitting-summary,thm:asymptotic-gradient-shrinking-soliton}
Let $(M,g(t))$ be a two-dimensional Ricci flow satisfying all the hypotheses of
Proposition~\ref{prop:splitting-at-infinity} except possibly non-compactness. It then follows
from Corollary~\ref{cor:no-two-dimensional-splitting-example} that $M$ is compact. This proves the first item.


Now suppose that $(M,g(t))$ is an asymptotic soliton for a
$\kappa$-solution of dimension two. If $(M,g(-1))$ does not have
bounded curvature, then there is a sequence $p_i\rightarrow \infty$
so that $\lim_{i\rightarrow \infty}R(p_i,-1)=\infty$. By this
and
 Theorem~\ref{thm:asymptotic-gradient-shrinking-soliton} the Ricci flow $(M,g(t))$ satisfies all the hypotheses
 of Proposition ~\ref{prop:splitting-at-infinity}. But this contradicts Corollary~\ref{cor:no-two-dimensional-splitting-example}.
 We conclude that $(M,g(-1))$ has bounded curvature. According to
 Corollary~\ref{cor:noncompact-soliton-splitting-summary} this means that
 $(M,g(t))$ is compact.
Results going back to Hamilton in \cite{Hamiltonsurface} imply that
this compact asymptotic shrinking soliton is a shrinking family of
compact, round surfaces. For example, this result is contained in
Proposition 5.21 on p. 118 of \cite{ChowKnopf}.
  This proves the second item.


 Now suppose that $(M,g(t))$ is a two-dimensional $\kappa$-solution.
 By the second item any asymptotic gradient shrinking soliton for
 this $\kappa$-solution is compact.
It follows that $M$ is compact. We know that as $t$ goes to $-\infty$ the
Riemannian surfaces $(M,g(t))$ are converging to compact, round surfaces.
Extend the flow forward from $0$ to a maximal time $\Omega<\infty$. By Theorem
5.64 on p. 149 of \cite{ChowKnopf} the surfaces $(M,g(t))$ are also becoming
round as $t$ approaches $\Omega$ from below. Also, according to Proposition
5.39 on p. 134 of \cite{ChowKnopf} the entropy of the flow is weakly monotone
decreasing and is strictly decreasing unless the flow is a gradient shrinking
soliton. But we have seen that the limits at both $-\infty$ and $\Omega$ are
round manifolds, and hence of the same entropy. It follows that the
$\kappa$-solution is a shrinking family of compact, round surfaces.
\end{proof}



Now that we have shown that every two-dimensional $\kappa$-solution
is a shrinking family of round surfaces, we can complete the proof
of Proposition~\ref{prop:no-noncompact-positive-curvature-soliton}. Let $(M,g)$ be a non-compact
manifold of positive curvature satisfying the hypotheses of
Theorem~\ref{thm:gradient-shrinking-soliton-classification}. According to Corollary~\ref{cor:two-dimensional-gradient-shrinking-soliton} the
limiting Ricci flow $(M_\infty,G_\infty(t))$ referred to in
Corollary~\ref{cor:noncompact-soliton-splitting-summary} is the product of a line and a
shrinking family of round surfaces. Since $(M,g)$ is non-compact and
has positive curvature, it is diffeomorphic to $\Ar ^3$ and hence
does not contain an embedded copy of a projective plane. It follows
that the round surfaces are in fact  round two-spheres. Thus,
$(M_\infty,G_\infty(t)),\ -\infty<t<0$,  splits as the product of a
shrinking family $(S^2,h(t)),\ -\infty<t<0$, of round two-spheres
and the real line.

\begin{claim}[Normalization of the limiting cylinder's curvature]\label{claim:cylinder-limit-scalar-curvature-normalization}
The scalar curvature of $(S^2,h(-1))$ is equal to $1$.
\end{claim}

\begin{proof}
Since the shrinking family of round two-spheres $(S^2,h(t))$ exists for all
$-\infty<t<0$, it follows that the scalar curvature of $(S^2,h(-1))$ is at most
$1$. On the other hand, since the scalar curvature is increasing along the
gradient flow lines of $f$, the infimum of the scalar curvature of $(M,g)$,
$R_\mathrm{inf}$, is positive. Thus, the infimum of the scalar curvature of
$(M,G(t))$ is $R_\mathrm{inf}/|t|$ and goes to infinity as $|t|$ approaches $0$.
Thus, the infimum of the scalar curvature of  $(S^2,h(t))$ goes to infinity as
$t$ approaches zero. This means that the shrinking family of two-spheres
becomes singular as $t$ approaches zero, and consequently the scalar curvature
of $(S^2,h(-1))$ is equal to $1$.
\end{proof}


 It follows that for any $p$ in a
neighborhood of infinity of $(M,g)$, we have
$$R_g(p)< 1.$$
For any unit vector $Y$ at any point of $M\setminus K$ we have
$$\Hess(f)(Y,Y)=\frac{1}{2}-\Ric(Y,Y)\ge \frac{1}{2}-\frac{R}{2}>0.$$
(On a manifold with non-negative curvature $\Ric(Y,Y)\le R/2$ for any unit
tangent vector $Y$.) This means that for $u$ sufficiently large the level
surfaces of $N_u=f^{-1}(u)$ are convex and hence have increasing area as $u$
increases.

According to Equations~(\ref{X}) and~(\ref{sqrts2}) the angle between $\nabla
f$ and the gradient of the distance function from $p$ goes to zero as we go to
infinity. According to Theorem~\ref{topsplit} the gradient of the distance
function from $p$ converges to the unit vector field in the $\Ar$-direction of
the product structure. It follows that the unit vector in the $\nabla
f$-direction converges to the unit vector in the $\Ar$-direction. Hence, as $u$
tends to $\infty$ the level surfaces $f^{-1}(u)$ converge in the $C^1$-sense to
$\Sigma\times \{0\}$. Thus, the areas of these level surfaces converge to the
area of $(\Sigma,h(-1))$ which is $8\pi$ since the scalar curvature of this
limiting surface is $\limsup_{p\in M} R(p,-1)= 1$. It follows that the
area of $f^{-1}(u)$ is less than $8\pi$ for all $u$ sufficiently large.

Now let us estimate the intrinsic curvature of $N=N_u=f^{-1}(u)$.
Let $K_N$ denote the sectional curvature of the induced metric on
$N$, whereas $K_M$ is the sectional curvature of $M$. We also denote
by $R_N$ the scalar curvature of the induced metric on $N$. Fix an
orthonormal basis $\{e_1,e_2,e_3\}$ at a point of $N$, where
$e_3=\nabla f/|\nabla f|$. Then by the Gauss-Codazzi formula we have
$$R_N=2K_N(e_1,e_2)=2(K_M(e_1,e_2)+\det\, S)$$
where $S$ is the shape operator
$$S=\frac{\Hess(f|TN)}{|\nabla f|}.$$
Clearly, we have $R-2\Ric(e_3,e_3)=2K_M(e_1,e_2)$, so that
$$R_N=R-2\Ric(e_3,e_3)+2\det\, S.$$
We can assume that the basis is chosen so that $\Ric|_{TN}$ is
diagonal; i.e., in the given basis we have
$$\Ric=\begin{pmatrix} r_1 & 0  & c_1 \\ 0 & r_2 & c_2 \\ c_1 & c_2 &
r_3\end{pmatrix}.$$ From the gradient shrinking soliton equation we
have $\Hess(f)=(1/2)g-\Ric$ so that \begin{eqnarray*}
\det(\Hess(f|TN)) & = & \left(\frac{1}{2}-r_1\right)\left(\frac{1}{2}- r_2\right) \\
& \le &
\frac{1}{4}(1-r_1-r_2)^2 \\
& = & \frac{1}{4} (1-R+\Ric(e_3,e_3))^2.
\end{eqnarray*}
Thus, it follows that \begin{equation}\label{R_N} R_N\le R-2
\Ric(e_3,e_3)+\frac{(1-R+\Ric(e_3,e_3))^2}{2|\nabla f|^2}.
\end{equation}
 It follows from Equation~(\ref{X}) that $|\nabla
f(x)|\rightarrow\infty$ as $x$ goes to infinity in $M$. Thus, since
the curvature of $(M,g(-1))$ is bounded, provided that $u$ is
sufficiently large, we have $1-R+\Ric(e_3,e_3)<2|\nabla f|^2$.
Since the left-hand side of this inequality is positive (since $R<
1$), it follows that
\begin{equation*}
(1-R+\Ric(e_3,e_3))^2  <  2(1-R+\Ric(e_3,e_3))|\nabla
f|^2.
\end{equation*}
Plugging this into Equation~(\ref{R_N}) gives that
$$R_N<1-\Ric(e_3,e_3)\le 1,$$
assuming that  $u$ is sufficiently large.

This contradicts the Gauss-Bonnet theorem for the surface $N$: Its
area is less than $8\pi$, and the  scalar curvature of the induced
metric is less than $1$, meaning that its Gaussian curvature is less
than $1/2$; yet $N$ is diffeomorphic to a $2$-sphere. This completes
the proof of Proposition~\ref{prop:no-noncompact-positive-curvature-soliton}, that is to say this shows
that there are no non-compact positive curved examples satisfying
the hypotheses of Theorem~\ref{thm:gradient-shrinking-soliton-classification}.

\subsection{Case of non-positive curvature revisited}



We return now to the second case of Theorem~\ref{thm:gradient-shrinking-soliton-classification}.  We
extend $(M,g)$ to a Ricci flow $(M,G(t))$ defined for $-\infty<t<0$
as given in Equation~(\ref{Gdefn}). By Claim~\ref{claim:non-strictly-positive-curvature-splitting} $M$
has either a one- or $2$-sheeted covering $\widetilde M$ such that
$(\widetilde M,\tilde G(t))$ is a metric product of a surface and a
one-manifold for all $t<0$. The evolving metric on the surface is
itself a $\kappa$-solution and hence by Corollary~\ref{cor:two-dimensional-gradient-shrinking-soliton} the
surfaces are compact and the metrics are all round. Thus, in this
case, for any $t<0$, the manifold $(\widetilde M,\tilde G(t))$ is a
metric product of a round $S^2$ or $\Ar P^2$ and a flat copy of
$\Ar$.  The conclusion in this case is that the one- or two-sheeted
covering $(\widetilde M,\tilde G(t))$ is a product of a round $S^2$
or $\Ar P^2$ and the line for all $t<0$.

\subsection{Completion of the proof of Theorem~{\ref{thm:gradient-shrinking-soliton-classification}}}



\begin{corollary}[The splitting-at-infinity limit is a family of round surfaces]\label{cor:three-dimensional-splitting-compact-surfaces}
\label{3Dinfty}
\uses{prop:splitting-at-infinity,cor:two-dimensional-gradient-shrinking-soliton}
Let $(M,g(t))$ be a three-dimensional Ricci flow satisfying the hypotheses of
Proposition~\ref{prop:splitting-at-infinity}. Then the limit constructed in that proposition
splits as a product of a shrinking family of compact round surfaces with a
line. In particular, for any non-compact gradient shrinking soliton of a
three-dimensional $\kappa$-solution the limit constructed in
Proposition~\ref{prop:splitting-at-infinity} is the product of a shrinking family of round
surfaces and the real line.
\end{corollary}

\begin{proof}
\uses{prop:splitting-at-infinity,cor:two-dimensional-gradient-shrinking-soliton}
Let $(M,g(t))$ be a three-dimensional Ricci flow satisfying the
hypotheses of Proposition~\ref{prop:splitting-at-infinity} and let $(N^2,h(t))\times
(\Ar,ds^2)$ be the limit constructed in that proposition. Since this
limit is $\kappa$-non-collapsed, $(N,h(t))$ is
$\kappa'$-non-collapsed for some $\kappa'>0$ depending only on
$\kappa$. Since the limit is not flat and has non-negative
curvature, the same is true for $(N,h(t))$. Since $\partial
R/\partial t\ge 0$ for the limit, the same is true for $(N,h(t))$.
That is to say $(N,h(t))$ satisfies all the hypotheses of
Proposition~\ref{prop:splitting-at-infinity} except possibly non-compactness. It now
follows from Corollary~\ref{cor:two-dimensional-gradient-shrinking-soliton} that $(N,h(t))$ is a shrinking
family of compact, round surfaces.
\end{proof}


\begin{corollary}[Asymptotic soliton of a 3D $\kappa$-solution has bounded curvature]\label{cor:three-dimensional-asymptotic-soliton-bounded-curvature}
\label{3DGSSkappa}
\uses{cor:three-dimensional-splitting-compact-surfaces,cor:two-dimensional-gradient-shrinking-soliton}
Let $(M,g(t)),\ -\infty<t<0$, be an asymptotic gradient shrinking soliton for a
three-dimensional $\kappa$-solution. Then for each $t<0$, the Riemannian
manifold $(M,g(t))$ has bounded curvature. In particular, for any $a>0$ the
flow $(M,g(t)),\ -\infty<t\le -a$, followed by a shift of time by $+a$ is a
$\kappa$-solution.
\end{corollary}

\begin{proof}
\uses{cor:three-dimensional-splitting-compact-surfaces,cor:two-dimensional-gradient-shrinking-soliton}
If an asymptotic gradient shrinking soliton $(M,g(t))$ of a three-dimensional
$\kappa$-solution does not have strictly positive curvature, then according to
Corollary~\ref{nullspace}, $(M,g(t))$ has a covering that splits as a product
of a a two-dimensional Ricci flow and a line. The two-dimensional Ricci flow
satisfies all the hypotheses of Proposition~\ref{prop:splitting-at-infinity} except possibly
compactness, and hence by Corollary~\ref{cor:two-dimensional-gradient-shrinking-soliton} it is a shrinking family of
round surfaces. In this case, it is clear that each time-slice of $(M,g(t))$
has bounded curvature.

Now we consider the remaining case when $(M,g(t))$ has strictly positive
curvature.  Assume that $(M,g(t))$ has unbounded curvature. Then there is a
sequence of points $p_i$ tending to infinity in $M$ such that $R(p_i,t)$ tends
to infinity. By Corollary~\ref{cor:three-dimensional-splitting-compact-surfaces} we can replace the points $p_i$ by
points $q_i$ with $Q_i=R(q_i,t)\ge R(p_i,t)$ so that the based Riemannian
manifolds $(M,Q_ig(t),q_i)$ converge to a product of a round surface $(N,h(t))$
with $\Ar$. The surface $N$ is either diffeomorphic to $S^2$ or $\Ar P^2$.
Since $(M,g(t))$ has positive curvature, by Theorem~\ref{thm:soul-theorem}, it is
diffeomorphic to $\Ar ^3$, and hence it contains no embedded $\Ar P^2$. It
follows that $(N,h(t))$ is a round two-sphere.

Fix $\epsilon>0$ sufficiently small as in Proposition~\ref{narrows}. Then the
limiting statement means that, for every $i$ sufficiently large, there is an
$\epsilon$-neck in $(M,g(t))$ centered at $q_i$ with scale $Q_i^{-1/2}$. This
contradicts Proposition~\ref{narrows}, establishing that for each $t<0$ the
curvature of $(M,g(t))$ is bounded.
\end{proof}

\begin{corollary}[The asymptotic soliton falls under the classification theorem]\label{cor:asymptotic-soliton-classification-application}
\uses{cor:three-dimensional-asymptotic-soliton-bounded-curvature,thm:gradient-shrinking-soliton-classification,thm:asymptotic-gradient-shrinking-soliton}
Let $(M,g(t)),\ -\infty<t\le 0$, be a $\kappa$-solution of dimension three.
Then any asymptotic gradient shrinking soliton $(M_\infty,g_\infty(t))$ for
this $\kappa$-solution, as constructed in Theorem~\ref{thm:asymptotic-gradient-shrinking-soliton}, is of one of
the three types listed in Theorem~\ref{thm:gradient-shrinking-soliton-classification}.
\end{corollary}


\begin{proof}
\uses{cor:three-dimensional-asymptotic-soliton-bounded-curvature,thm:gradient-shrinking-soliton-classification,thm:asymptotic-gradient-shrinking-soliton}
Let $(M_\infty,g_\infty(t)),\ -\infty<t<0$, be an asymptotic
gradient shrinking soliton for $(M,g(t))$. According to
Corollary~\ref{cor:three-dimensional-asymptotic-soliton-bounded-curvature}, this soliton is a $\kappa$-solution,
implying that $(M_\infty,g_\infty(-1))$ is a complete Riemannian
manifold of bounded, non-negative curvature. Suppose that
$B(p,-1,r)\subset M_\infty$ is a metric ball and $|\Rm_{g_\infty}|(x,-1)\le r^{-2}$ for all $x\in B(p,-1,r)$. Since
$\partial R_{g_\infty}(x,t)/\partial t\ge 0$, it follows that
$R(x,t)\le 3r^{-2}$ on $B(p,-1,r)\times (-1-r^2,-1]$, and hence that
$|\Rm_{g_\infty}|\le 3r^{-2}$ on this same region. Since the
Ricci flow $(M_\infty,g_\infty(t))$ is $\kappa$-non-collapsed, it
follows that $\Vol\, B(p,-1,r/\sqrt{3})\ge
\kappa(r/\sqrt{3})^3$. Hence, $\Vol\, B(p,-1,r)\ge
(\kappa/3\sqrt{3})r^3$. This proves that the manifold
$(M_\infty,g_\infty(-1))$ is $\kappa'$-non-collapsed for some
$\kappa'>0$ depending only on $\kappa$. On the other hand, according
to Theorem~\ref{thm:asymptotic-gradient-shrinking-soliton} there is a function $f(\cdot, -1)$ from
$M_\infty$ to $\Ar$ satisfying the gradient shrinking soliton
equation at the time-slice $-1$.
 Thus,
Theorem~\ref{thm:gradient-shrinking-soliton-classification} applies to $(M_\infty,g_\infty(-1))$ to
produce a Ricci flow $G(t),\ -\infty<t<0$, of one of the three types
listed in that theorem and with $G(-1)=g_\infty(-1)$.

Now we must show that $G(t)=g_\infty(t)$ for all $t<0$. In the first
case when $M_\infty$ is compact, this is clear by uniqueness of the
Ricci flow in the compact case. Suppose that $(M_\infty,G(t))$ is of
the second type listed in Theorem~\ref{thm:gradient-shrinking-soliton-classification}. Then
$(M_\infty,g_\infty(-1))$ is a product of a round two-sphere and the
real line. By Corollary~\ref{nullspace} this implies that the entire
flow $(M_\infty,g_\infty(t))$ splits as the product of a flow of
compact two-spheres and the real line. Again by uniqueness in the
compact case, this family of two-spheres must be a shrinking family
of round two-spheres. In the third case, one passes to a finite
sheeted covering of the second type, and applies the second case.
\end{proof}






\subsection{Asymptotic curvature}

There is one elementary result that will be needed in what follows.

\begin{definition}[Asymptotic scalar curvature]\label{def:asymptotic-scalar-curvature}
Let $(M,g)$ be a complete, connected, non-compact Riemannian manifold of
non-negative curvature. Fix a point $p\in M$. We define the {\em asymptotic
scalar curvature}
$${\mathcal R}(M,g)=\limsup_{x\rightarrow\infty}R(x)d^2(x,p).$$ Clearly, this limit is
independent of $p$.
\end{definition}

\begin{proposition}[Non-compact $\kappa$-solutions have infinite asymptotic curvature]\label{prop:asymptotic-scalar-curvature-infinite}
\label{asympscal}
\uses{def:asymptotic-scalar-curvature,def:ancient-solution,def:kappa-noncollapsed}
Suppose that $(M,g(t)),\ -\infty<t<0$, is a connected, non-compact
$\kappa$-solution of dimension at most three\footnote{This result, in fact,
holds in all dimensions.}.  Then ${\mathcal R}(M,g(t))=+\infty$ for every
$t<0$.
\end{proposition}


\begin{proof}
\uses{cor:two-dimensional-gradient-shrinking-soliton,claim:non-strictly-positive-curvature-splitting}
By Corollary~\ref{cor:two-dimensional-gradient-shrinking-soliton} the only two-dimensional $\kappa$-solutions are
compact, so that the result is vacuously true in this case. Suppose that
$(M,g(t))$ is three-dimensional
 If $(M,g(t))$ does
not have strictly positive curvature, then, since it is not flat, by
Corollary~\ref{nullspace} it must be three-dimensional and it has a
finite-sheeted covering space that splits as a product $(Q,h(0))\times
(\Ar,ds^2)$ with $(Q,h(0))$ being a surface of strictly positive curvature and
$T$ being a flat one-manifold. Clearly, in this case the asymptotic curvature
is infinite.

Thus, without loss of generality we can assume that $(M,g(t))$ has
strictly positive curvature. Let us first consider the case when
${\mathcal R}(M,g(t))$ has a finite, nonzero value. Fix a point
$p\in M$. Take a sequence of points $x_n$ tending to infinity and
set $\lambda_n=d^2_0(x_n,p)$ and $Q_n=R(x_n,t)$. We choose this
sequence such that
$$\lim_{n\rightarrow\infty}Q_n\lambda_n={\mathcal R}(M,g(t)).$$
We consider the sequence of Ricci flows $(M,h_n(t),(x_n,0))$, where
$$h_n(t)=Q_ng(Q_n^{-1}t).$$

 Fix $0<a<\sqrt{{\mathcal R}(M,g(t))}<b<\infty$. Consider  the annuli
$$A_n=\{y\in M\left|\right. a<d_{h_n(0)}(y,p)<b\}.$$
Because of the choice of sequence, for all $n$ sufficiently large,
the scalar curvature of the restriction of $h_n(0)$ to $A_n$ is
bounded independent of $n$. Furthermore, since $d_{h_n}(p,x_n)$
converges to $\sqrt{{\mathcal R}(M,g(t))}$, there is $\alpha>0$ such
that for all $n$ sufficiently large, the annulus $A_n$ contains
$B_{h_n}(x_n,0,\alpha)$. Consequently, we have a bound, independent
of $n$, for the scalar curvature of $h_n(0)$ on these balls. By the
hypothesis that $\partial R/\partial t\ge 0$, there is a bound,
independent of $n$, for the scalar curvature of $h_n$ on
$B_{h_n}(x_n,0,\alpha)\times (-\infty,0]$. Using the fact that the
flows have non-negative curvature, this means that there is a bound,
independent of $n$, for $|\Rm_{h_n}(y,0)|$ on
$B_{h_n}(x_n,0,\alpha)\times (-\infty,0]$.
 This means
that by Shi's theorem (Theorem~\ref{thm:shi-derivative-estimates}), there are bounds,
independent of $n$, for every covariant derivative of the curvature
on $B_{h_n}(x_n,0,\alpha/2)\times (-\infty,0]$.

Since the original flow is $\kappa$-non-collapsed on all scales, it
follows that the rescaled flows are also $\kappa$ non-collapsed on
all scales. Since the curvature is bounded, independent of $n$, on
$B_{h_n}(x_n,0,\alpha)$, this implies that there is $\delta>0$,
independent of $n$, such that for all $n$ sufficiently large, every
ball of radius $\delta$ centered at any point of
$B_{h_n}(x_n,0,\alpha/2)$ has volume at least $\kappa \delta^3$, Now
applying Theorem~\ref{basicconv} we see that  a subsequence
converges geometrically to a limit which will automatically be a
metric ball $B_{g_\infty}(x_\infty,0,\alpha/2)$. In fact, by
Hamilton's result (Proposition~\ref{partialflowlimit}) there is a
limiting flow on $B_{g_\infty}(x_\infty,0,\alpha/4)\times
(-\infty,0]$. Notice that the limiting flow is not flat since
$R(x_\infty,0)=1$.



On the other hand, according to Lemma~\ref{conelimit} the
Gromov-Hausdorff limit of a
subsequence $(M,\lambda_n^{-1}g_n(0),x_n)$ is the Tits
cone, i.e., the cone over $S_\infty(M,p)$. Since
$Q_n={\mathcal R}(M,g(t))\lambda_n^{-1}$, the rescalings
$(M,Q_ng_n(0),x_n)$  also converge to  a cone, say $(C,h,y_\infty)$,
which is in fact simply a rescaling of the Tits cone by a factor
${\mathcal R}(M,g(t))$.  Pass to a subsequence so that both the
geometric limit on the ball of radius $\alpha/2$ and the
Gromov-Hausdorff limit exist. Then the geometric limit
$B_{g_\infty}(x_\infty,0,\alpha/2)$ is isometric to an open ball in
the cone. Since we have a limiting Ricci flow
$$(B_{g_\infty}(x_\infty,0,\alpha/2),g_\infty(t)),\ -\infty<t\le 0,$$
 this contradicts Proposition~\ref{nocones}. This completes the proof that
 it is not possible for the asymptotic curvature to be finite
and nonzero.





Lastly, we consider the possibility that the asymptotic curvature is zero.
Again we fix $p\in M$. Take any sequence of points $x_n$ tending to infinity
and let $\lambda_n=d^2_0(p,x_n)$. Form the sequence of based Ricci flows
$(M,h_n(t),(x_n,0))$ where $h_n(t)=\lambda_n^{-1}g(\lambda_nt)$. On the one
hand, the Gromov-Hausdorff limit (of a subsequence) is the Tits cone. On the
other hand, the curvature condition tells us the following: For any $0<a<1<b$
on the regions
$$\{y\in M\left|\right. a< d_{h_n(0)}(y,p)<b\},$$
the curvature tends uniformly to zero as $n$ tends to infinity. Arguing as in
the previous case, Shi's theorem, Hamilton's result,
Theorem~\ref{partialflowlimit}, and the fact that the original flow is $\kappa$
non-collapsed on all scales tells us that we can pass to a subsequence so that
these annuli centered at $x_n$ converge geometrically to a limit. Of course,
the limit is flat. Since this holds for all $0<a<1<b$, this implies that the
Tits cone is smooth and flat except possibly at its cone point. In particular,
the sphere at infinity, $S_\infty(M,p)$, is a smooth surface of constant
curvature $+1$.

\begin{claim}[Sphere at infinity is round]\label{claim:sphere-at-infinity-round}
 $S_\infty(M,p)$ is isometric to a round $2$-sphere.
 \end{claim}

\begin{proof}
Since $M$ is orientable the complement of the cone point in the Tits
cone is an orientable manifold and hence $S_\infty(M,p)$ is an
orientable surface. Since we have already established that it has a
metric of constant positive curvature, it must be diffeomorphic to
$S^2$, and hence isometric to a round sphere. (In higher dimensions
one can prove that $S_\infty(M,p)$ is simply connected, and hence
isometric to a round sphere.)
\end{proof}

It follows that the Tits cone is a smooth flat manifold even at the origin, and
hence is isometric to Euclidean $3$-space. This means that in the limit, for
any $r>0$ the volume of the ball of radius $r$ centered at the cone point
  is exactly $\omega_3 r^3$, where
$\omega_3$ is the volume of the unit ball in $\Ar^3$. Consequently,
$$\lim_{n\rightarrow\infty}
\Vol\left(B_{g}(p,0,\sqrt{\lambda_n}r)\setminus
B_g(p,0,1)\right)\rightarrow \omega_3\lambda_n^{3/2}r^3.$$ By
Theorem~\ref{BishopGromov} and the fact that the Ricci curvature is
non-negative, this implies that
$$\Vol\, B_g(p,0,R)=\omega_3R^3$$
for all $R<\infty$. Since the Ricci curvature is non-negative, this
means that $(M,g(t))$ is Ricci-flat, and hence flat. But this
contradicts the fact that $(M,g(t))$ is a $\kappa$-solution and
hence is not flat.

Having ruled out the other two cases, we are left with only one possibility:
${\mathcal R}(M,g(t))=\infty$.
\end{proof}


\section{Universal $\kappa$}

The first consequence of the existence of an asymptotic gradient
shrinking soliton is that there is a universal $\kappa$ for all
$3$-dimensional $\kappa$-solutions, except those of constant
positive curvature.


\begin{proposition}[Universal $\kappa_0$ for non-round $3$-dimensional $\kappa$-solutions]\label{prop:universal-kappa-non-round}
\label{kappa0}
\uses{cor:compact-asymptotic-soliton-classification,thm:gradient-shrinking-soliton-classification}
There is a $\kappa_0>0$ such that any non-round $3$-dimensional
$\kappa$-solution is a $\kappa_0$-solution.
\end{proposition}

\begin{proof}
\uses{cor:compact-asymptotic-soliton-classification,thm:gradient-shrinking-soliton-classification}
Let $(M,g(t))$ be a non-round $3$-dimensional $\kappa$-solution. By
Corollary~\ref{cor:compact-asymptotic-soliton-classification} since $(M,g(t))$ is not a family of
round manifolds, the asymptotic soliton for the $\kappa$-solution
cannot be compact. Thus, according to Corollary~\ref{thm:gradient-shrinking-soliton-classification} there
are only two possibilities for the asymptotic soliton
$(M_\infty,g_\infty(t))$ -- either $(M_\infty,g_\infty(t))$ is the
product of a round $2$-sphere of Gaussian curvature $1/2|t|$ with a
line or has a two-sheeted covering by such a product. In fact, there
are three possibilities: $S^2\times \Ar$, $\Ar P^2\times \Ar$ or the
twisted $\Ar$-bundle over $\Ar P^2$ whose total space is
diffeomorphic to the complement of a point in $\Ar P^3$.


Fix a point $x=(p,0)\in M\times\{0\}$. Let $\bar\tau_k$ be a
sequence converging to $\infty$, and $q_k\in M$ a point with
$l_{x}(q_k,\bar\tau_k)\le 3/2$. The existence of an asymptotic
soliton means that, possibly after passing to a subsequence, there
is a gradient shrinking soliton $(M_\infty,g_\infty(t))$ and a ball
$B$ of radius $1$ in $(M_\infty,g_\infty(-1))$ centered at a point
$q_\infty\in M_\infty$ and a sequence of embeddings $\psi_k\colon
B\to M$ such that $\psi_k(q_\infty)=q_k$ and such that the map
$$B\times [-2,-1]\to M\times [-2\bar\tau_k,-\bar\tau_k]$$
given by $(b,t)\mapsto (\psi_k(b),\bar\tau_kt)$ has the property
that the pullback of $\bar\tau^{-1}_kg(\bar\tau_k t)$ converges
smoothly and uniformly as $k\rightarrow\infty$ to the restriction of
$g_\infty(t),\ -2\le t\le -1$, to $B$. Let $(M_k,g_k(t))$ be this
rescaling of the the $\kappa$-solution by $\tau_k$. Then the
embeddings $\psi_k\times \id\colon B\times (-2,-1]\to
(M_k\times [-2,-1]$ converge as $k\rightarrow \infty$ to a
one-parameter family of isometries. That is to say, the image
$\psi_k(B\times [-2,-1])\subset M_k\times [-2,-1]$ is an almost
isometric embedding. Since the reduced length function from $x$ to
$(\psi_k(a),-1)$ is at most $3/2$ (from the invariance of reduced
length under rescalings, see Corollary~\ref{COR}), it follows easily
that the reduced length function on $\psi_k(B\times \{-2\})$ is
bounded independent of $k$. Similarly, the volume of
$\psi_k(B\times\{-2\})$ is bounded independent of $k$. This means
the reduced volume of $\psi_k(B\times\{-2\})$ in $(M_k,g_k(t))$ is
bounded independent of $k$. Now by Theorem~\ref{THM} this implies
that $(M_k,g_k(t))$ is $\kappa_0$-non-collapsed at $(p,0)$ on scales
$\le \sqrt{2}$ for some $\kappa_0$ depending only on the geometry of
the three possibilities for $(M_\infty,g_\infty(t)),\ -2\le t\le
-1$. Being $\kappa_0$-non-collapsed is invariant under rescalings,
so that it follows immediately that $(M,g(t))$ is
$\kappa_0$-non-collapsed on scales $\le \sqrt{2\bar\tau_k}$. Since
this is true for all $k$, it follows that $(M,g(t))$ is
$\kappa_0$-non-collapsed on all scales at $(p,0)$.

This result holds of course for every $p\in M$, showing that at
$t=0$ the flow is $\kappa_0$-non-collapsed. To prove this result at
points of the form $(p,t)\in M\times(-\infty,0]$ we simply shift the
original $\kappa$-solution upward by $|t|$ and remove the part of
the flow at positive time. This produces a new $\kappa$-solution and
the point in question has been shifted to the time-zero slice, so
that we can apply the previous results.
\end{proof}



\section{Asymptotic volume}

Let $(M,g(t))$ be an $n$-dimensional $\kappa$-solution. For any
$t\le 0$ and any point $p\in M$ we consider $(\Vol\,B_{g(t)}(p,r))/r^n$. According to the Bishop-Gromov Theorem
(Theorem~\ref{thm:bishop-gromov}), this is a non-increasing function of
$r$. We define the {\em asymptotic volume} ${\mathcal V}_\infty(M,g(t))$, or ${\mathcal
V}_\infty(t)$ if the flow is clear from the context, to be the limit
as $r\rightarrow\infty$ of this function. Clearly, this limit is
independent of $p\in M$.

\begin{theorem}[Asymptotic volume of a $\kappa$-solution is zero]\label{thm:asymptotic-volume-zero}
\label{asympvol}
\uses{def:ancient-solution,def:kappa-noncollapsed}
For\footnote{This theorem and all the other results of this section are valid
in all dimensions. Our proofs use Theorem~\ref{prop:asymptotic-scalar-curvature-infinite} and
Proposition~\ref{prop:splitting-at-infinity} which are also valid in all dimensions but which we
proved only in dimensions $2$ and $3$. Thus, while we state the results of this
section for all dimensions, strictly speaking we give proofs only for
dimensions $2$ and $3$. These are the only cases we need in what follows.} any
$\kappa>0$ and any $\kappa$-solution $(M,g(t))$ the asymptotic volume
${\mathcal V}_\infty(M,g(t))$ is identically zero.
\end{theorem}


\begin{proof}
\uses{cor:two-dimensional-gradient-shrinking-soliton,prop:splitting-at-infinity}
The proof is by induction on the dimension $n$ of the solution. For $n=2$ by
Corollary~\ref{cor:two-dimensional-gradient-shrinking-soliton} there are only compact $\kappa$-solutions, which clearly
have zero asymptotic volume. Suppose that we have established the result for
$n-1\ge 2$ and let us prove it for $n$.

According to Proposition~\ref{prop:splitting-at-infinity} there is a sequence of points $p_n\in
M$ tending to infinity such that setting $Q_n=R(p_n,0)$ the sequence of Ricci
flows
$$(M,Q_ng(Q_n^{-1}t),(q_n,0))$$ converges geometrically to a limit
$(M_\infty,g_\infty(t),(q_\infty,0))$, and this limit splits off a
line: $(M_\infty,g_\infty(t))= (N,h(t))\times \Ar$. Since the ball
of radius $R$ about a point $(x,t)\in N\times \Ar$ is contained in
the product of the ball of radius $R$ about $x$ in $N$ and an
interval of length $2R$, it follows that $(N,h(t))$ is a
$\kappa/2$-ancient solution. Hence, by induction, for every $t$, the
asymptotic volume of $(N,h(t))$ is zero, and hence so is that of
$(M,g(t))$.
\end{proof}



\subsection{Volume comparison}

One important consequence of the asymptotic volume result is a
volume comparison result.

\begin{proposition}[Volume comparison for $\kappa$-non-collapsed balls]\label{prop:volume-comparison-bound}
\label{volcomp}
\uses{thm:asymptotic-volume-zero}
Fix the dimension $n$.
For every $\nu>0$ there is $A<\infty$ such that the following holds.
Suppose that $(M_k,g_k(t)),\ -t_k\le t\le 0$, is a sequence of (not
necessarily complete) $n$-dimensional Ricci flows of non-negative
curvature operator. Suppose in addition we have points $p_k\in M_k$
and radii $r_k>0$ with the property that for each $k$ the ball
$B(p_k,0,r_k)$ has compact closure in $M_k$. Let $Q_k=R(p_k,0)$ and
suppose that $R(q,t)\le 4Q_k$ for all $q\in B(p_k,0,r_k)$ and for
all $t\in [-t_k,0]$, and suppose that $t_kQ_k\rightarrow\infty$ and
$r_k^2Q_k\rightarrow\infty$ as $k\rightarrow\infty$. Then $\Vol\,B(p_k,0,A/\sqrt{Q_k})< \nu(A/\sqrt{Q_k})^n$ for all $k$
sufficiently large.
\end{proposition}


\begin{proof}
\uses{thm:asymptotic-volume-zero}
Suppose that the result fails for some $\nu>0$. Then there is a
sequence $(M_k,g_k(t)),\ -t_k\le t\le 0$, of $n$-dimensional Ricci
flows, points $p_k\in M_k$, and radii $r_k$ as in the statement of
the lemma such that for every $A<\infty$ there is an arbitrarily
large $k$ with $\Vol\,B(p_k,0,A/\sqrt{Q_k})\ge
\nu(A/\sqrt{Q_k})^n$. Pass to a subsequence so that for each
$A<\infty$ we have
$$\Vol\,B(p_k,0,A/\sqrt{Q_k})\ge\nu(A/\sqrt{Q_k})^n$$ for all $k$
sufficiently large. Consider now the flows
$h_k(t)=Q_kg_k(Q_k^{-1}t)$, defined for $-Q_kt_k\le t\le 0$. Then
for every $A<\infty$ for all $k$ sufficiently large we have
$R_{h_k}(q,t)\le 4$ for all $q\in B_{h_k}(p_k,0,A)$ and all $t\in
(-t_kQ_k,0]$. Also, for every $A<\infty$ for all $k$ sufficiently
large we have $\Vol\,B(p_k,0,A)\ge\nu A^n$. According to
Theorem~\ref{flowlimit} we can then pass to a subsequence that has a
geometric limit which is an ancient flow of complete Riemannian
manifolds. Clearly, the time-slices of the limit have non-negative
curvature operator, and the scalar curvature is bounded (by $4$) and
is equal to $1$ at the base point of the limit. Also, the asymptotic
volume ${\mathcal V}(0)\ge \nu$.

\begin{claim}[Monotonicity of the asymptotic volume]\label{claim:asymptotic-volume-monotone}
 Suppose that $(M,g(t))$ is an ancient Ricci flow such that for each
 $t\le 0$ the Riemannian manifold $(M,g(t))$ is complete and has
bounded, non-negative  curvature operator. Let ${\mathcal V}(t)$ be
the asymptotic volume of the manifold $(M,g(t))$.
\begin{enumerate}
\item[(1)] The asymptotic volume ${\mathcal V}(t)$ is a non-increasing function
of $t$. \item[(2)] If ${\mathcal V}(t)=V>0$ then every metric ball
$B(x,t,r)$ has volume at least $Vr^n$.
\end{enumerate}
\end{claim}


\begin{proof}
We begin with the proof of the first item. Fix $a<b\le 0$. By
hypothesis there is a constant $K<\infty$ such that the scalar
curvature of $(M,g(0))$ is bounded by $(n-1)K$. By the Harnack
inequality (Corollary~\ref{posderiv}) the scalar curvature of
$(M,g(t))$ is bounded by $(n-1)K$ for all $t\le 0$. Hence, since the
$(M,g(t))$ have non-negative curvature, we have $\Ric(p,t)\le
(n-1)K$ for all $p$ and $t$. Set $A=4(n-1)\sqrt{\frac{2K}{3}}$. Then
by Corollary~\ref{corI.8.3} we have
$$d_a(p_0,p_1)\le d_b(p_0,p_1)+A(b-a).$$
 This means that for any $r>0$ we have
$$B(p_0,b,r)\subset B(p_0,a,r+A(b-a)).$$

On the other hand, since $d\Vol/dt=-Rd\Vol$, it follows
that in the case of non-negative curvature that the volume of any
open set is non-increasing in time. Consequently,
$$\Vol_{g(b)}B(p_0,b,r)\le
\Vol_{g(a)}B(p_0,a,r+A(b-a)),$$ and hence
$$\frac{\Vol_{g(b)}B(p_0,b,r)}{r^n}\le\frac{\Vol_{g(a)}B(p_0,a,r+A(b-a))}{(r+A(b-a))^n}\frac{(r+A(b-a))^n}{r^n}.$$
Taking the limit as $r\rightarrow \infty$ gives
$${\mathcal V}(b)\le {\mathcal V}(a).$$

The second item of the claim is immediate from the Bishop-Gromov
inequality (Theorem~\ref{thm:bishop-gromov}).
\end{proof}

Now we return to the proof of the proposition. Under the assumption
that there is a counterexample to the proposition for some $\nu>0$,
we have constructed a limit that is an ancient Ricci flow with
bounded, non-negative curvature with ${\mathcal V}(0)\ge \nu$. Since
${\mathcal V}(0)\ge \nu$, it follows from the claim that ${\mathcal
V}(t)\ge \nu$ for all $t\le 0$ and hence, also by the claim,  we see
that $(M,g(t))$ is $\nu$-non-collapsed for all $t$. This completes
the proof that the limit is a $\nu$-solution. This contradicts
Theorem~\ref{thm:asymptotic-volume-zero} applied with $\kappa=\nu$, and proves the
proposition.
\end{proof}

This proposition has two useful corollaries about balls in
$\kappa$-solutions with volumes bounded away from zero. The first
says that the normalized curvature is bounded on such balls.

\begin{corollary}[Volume bounded below implies curvature bounded above]\label{cor:volume-bound-implies-curvature-bound}
\label{bdvolbdcurv}
\uses{prop:volume-comparison-bound}
For any $\nu>0$ there is a $C=C(\nu)<\infty$ depending only on the
dimension $n$ such that the following holds. Suppose that
$(M,g(t)),\ -\infty<t\le 0$, is an $n$-dimensional Ricci flow with
each $(M,g(t))$ being complete and with bounded, non-negative
curvature operator. Suppose $p\in M$, and $r>0$ are such that
$\Vol\,B(p,0,r)\ge \nu r^n$. Then $r^2R(q,0)\le C$ for all $q\in
B(p,0,r)$.
\end{corollary}

\begin{proof}
\uses{prop:volume-comparison-bound,lem:point-picking-general,thm:asymptotic-volume-zero}
Suppose that the result fails for some $\nu>0$. Then there is a
sequence $(M_k,g_k(t))$ of $n$-dimensional Ricci flows, complete,
with bounded non-negative curvature operator and  points $p_k\in
M_k$, constants $r_k>0$, and points $q_k\in B(p_k,0,r_k)$ such that:
\begin{enumerate}
\item[(1)] $\Vol\,B(p_k,0,r_k)\ge \nu r_k^n$, and
\item[(2)] setting $Q_k=R(q_k,0)$ we have $r_k^2Q_k\rightarrow\infty$ as $k\rightarrow\infty$.
\end{enumerate}
Using Lemma~\ref{lem:point-picking-general} we can find points $q_k'\in B(p_k,0,2r_k)$
and constants $s_k\le r_k$, such that setting $Q_k'=R(q_k',0)$ we
have $Q_k's_k^2=Q_kr_k^2$ and  $R(q,0)<4Q_k'$ for all $q\in
B(q_k',0,s_k)$.  Of course, $Q'_ks_k^2\rightarrow\infty$ as
$k\rightarrow\infty$. Since $d_0(p_k,q'_k)<2r_k$, we have
$B(p_k,0,r_k)\subset B(q_k',0,3r_k)$ so that
$$\Vol\,B(q_k',0,3r_k)\ge \Vol\,B(p_k,0,r_k)\ge \nu
r_k^n=(\nu/3^n)(3r_k)^n.$$ Since the sectional curvatures of
$(M,g_k(0))$ are non-negative, it follows from the Bishop-Gromov
inequality (Theorem~\ref{thm:bishop-gromov}) that $\Vol\,B(q_k',0,s)\ge (\nu/3^n)s^n$ for any $s\le s_k$.

Of course, by Corollary~\ref{posderiv}, we have $R(q,t)<4Q_k'$ for
all $t\le 0$ and all $q\in B(q'_k,0,s_k)$. Now consider the sequence
of based, rescaled flows $$(M_k,Q'_kg(Q_k'^{-1}t),(q_k',0)).$$ In
these manifolds all balls centered at $(q_k',0)$ of radii at most
$\sqrt{Q_k}s_k$ are
 $(\nu/3^n)$ non-collapsed. Also, the curvatures of these manifolds
 are non-negative and the scalar curvature is bounded by $4$.
 It follows that by passing to a subsequence we can extract a
geometric limit.
 Since $Q_k's_k^2\rightarrow\infty$
as $k\rightarrow\infty$ the asymptotic volume of this limit is at
least $\nu/3^n$. But this geometric limit is a
$\nu/3^n$-non-collapsed ancient solution with non-negative curvature
operator with scalar curvature bounded by $4$.  This contradicts
Theorem~\ref{thm:asymptotic-volume-zero}.
\end{proof}

The second corollary gives curvature bounds at all points in terms
of the distance to the center of the ball.


\begin{corollary}[Global curvature bound from a volume lower bound]\label{cor:global-curvature-bound-from-volume}
\label{global}
\uses{cor:volume-bound-implies-curvature-bound}
Fix the dimension $n$. Given $\nu>0$,  there is a function
$K(A)<\infty$, defined for $A\in (0,\infty)$, such that if
$(M,g(t)),\ -\infty<t\le 0$, is an $n$-dimensional Ricci flow,
complete of bounded, non-negative curvature operator, $p\in M$ is a
point and $0<r<\infty$ is such that $\Vol\,B(p,0,r)\ge \nu r^n$
then for all $q\in M$ we have
$$(r+d_0(p,q))^2R(q,0)\le K(d_0(p,q)/r).$$
\end{corollary}


\begin{proof}
\uses{cor:volume-bound-implies-curvature-bound}
Fix $q\in M$ and let $d=d_0(p,q)$. We have
$$\Vol\,B(q,0,r+d)\ge
\Vol\,B(p,0,r)\ge \nu r^n=\frac{\nu}{(1+(d/r))^n}(r+d)^n.$$ Let
$K(A)=C(\nu/^n)$, where $C$ is the constant provided by the previous
corollary. The result is immediate from the previous corollary.
\end{proof}

\section{Compactness of the space of $3$-dimensional
$\kappa$-solutions}

This section is devoted to proving the following result.

\begin{theorem}[Compactness of based $3$-dimensional $\kappa$-solutions]\label{thm:kappa-solution-compactness}
\label{compactkappa}
\uses{def:ancient-solution,def:kappa-noncollapsed}
Let $(M_k,g_k(t),(p_k,0))$ be a sequence of based $3$-dimensional
$\kappa$-solutions satisfying $R(p_k,0)=1$. Then there is a
subsequence converging smoothly to a based $\kappa$-solution.
\end{theorem}

The main point in proving this theorem is to establish the uniform
curvature bounds given in the next lemma.




\begin{lemma}[Uniform curvature bound near the base point]\label{lem:uniform-curvature-bound-kappa-solution}
\label{proposition}
\uses{thm:kappa-solution-compactness}
For each $r<\infty$ there is a constant $C(r)<\infty$, such that the
following holds. Let $(M,g(t),(p,0))$ be a  based $3$-dimensional
$\kappa$-solution satisfying $R(p,0)=1$. Then $R(q,0)\le C(r)$ for
all $q\in B(p,0,r)$.
\end{lemma}


\begin{proof}
\uses{prop:asymptotic-scalar-curvature-infinite}
Fix a based $3$-dimensional $\kappa$-solution $(M,g(t),(p,0))$. By
Theorem~\ref{prop:asymptotic-scalar-curvature-infinite} we have
$$\sup_{q\in M}d_{0}(p,q)^2R(q,0)=\infty.$$
Let $q$ be a closest point to $p$ satisfying
$$d_0(p,q)^2R(q,0)=1.$$

We set $d=d_{0}(p,q)$, and we set $Q=R(q,0)$. Of course, $d^2Q=1$.
We carry this notation and these assumptions through the next five
claims. The goal of these claims is to show that $R(q',0)$ is
uniformly bounded for $q'$ near $(p,0)$ so that in fact the distance
$d$ from the point $q$ to $p$ is uniformly bounded from below by a
positive constant (see Claim~\ref{claim:basepoint-neighborhood-curvature-bound} for a more precise
statement). Once we have this the lemma will follow easily. To
establish this uniform bound requires a sequence of claims.





\begin{claim}[Universal curvature ratio bound near $q$]\label{claim:curvature-ratio-bound-near-basepoint}
\label{2Rbound}
\uses{cor:volume-bound-implies-curvature-bound,cor:three-dimensional-asymptotic-soliton-bounded-curvature}
There is a universal (i.e., independent of the $3$-dimensional
$\kappa$-solution) upper bound $C$ for $R(q',0)/R(q,0)$ for all
$q'\in B(q,0,2d)$.
\end{claim}

\begin{proof}
Suppose not. Then there is a sequence $(M_k,g_k(t),(p_k,0))$ of
$3$-dimensional $\kappa$-solutions with $R(p_k,0)=1$, points $q_k$
in $(M_k,g_k(0))$ closest to $p_k$ satisfying $d_k^2R(q_k,0)=1$,
where $d_k=d_0(p_k,q_k)$, and points $q_k'\in B(q_k,0,2d_k)$ with
$$\lim_{k\rightarrow\infty}(2d_k)^2R(q_k',0)=\infty.$$ Then according
to Corollary~\ref{cor:volume-bound-implies-curvature-bound} for every $\nu>0$ for all $k$
sufficiently large, we have \begin{equation}\label{V3}
\Vol\,B(q_k,0,2d_k)<\nu (2d_k)^3.\end{equation} Therefore, by
passing to a subsequence, we can assume that for each $\nu>0$
\begin{equation}\label{lowerbd}
\Vol\,B(q_k,0,2d_k)< \nu (2d_k)^3 \end{equation} for all $k$
sufficiently large. Let $\omega_3$ be the volume of the unit ball in
$\Ar^3$. Then for all $k$ sufficiently large, $\Vol\,B(q_k,0,2d_k)<[\omega_3/2](2d_k)^3$. Since the sectional
curvatures of $(M_k,g_k(0))$ are non-negative, by the Bishop-Gromov
inequality (Theorem~\ref{thm:bishop-gromov}),  it follows that for every
$k$ sufficiently large there is $r_k<2d_k$ such that
\begin{equation}\label{V/2}
\Vol\,B(q_k,0,r_k)=[\omega_3/2]r_k^3.\end{equation} Of course,
because of Equation~(\ref{lowerbd}) we see that $\lim_{k\rightarrow\infty}r_k/d_k=0$. Then, according to
Corollary~\ref{cor:global-curvature-bound-from-volume}, we have for all $q\in M_k$

$$(r_k+d_{g_k(0)}(q_k,q))^2R(q,0)\le K(d_{g_k(0)}(q_k,q)/r_k),$$
where $K$ is as given in Corollary~\ref{cor:global-curvature-bound-from-volume}. Form the sequence
$(M_k,g'_k(t),(q_k,0))$, where $g'_k(t)=r_k^{-2}g_k(r_k^2t)$. This
is a sequence  of based Ricci flows. For each $A<\infty$ we have
$$(1+A)^2R_{g'_k}(q,0)\le K(A)$$ for all
$q\in B_{g_k'(0)}(q_k,0,A)$. Hence, by the consequence of Hamilton's
Harnack inequality (Corollary~\ref{posderiv})
$$R_{g'_k}(q,t)\le K(A),$$
for all  $(q,t)\in B_{g_k'(0)}(q_k,0,A)\times(-\infty,0]$. Using
this and the fact that all the flows are $\kappa$-non-collapsed,
Theorem~\ref{flowlimit} implies that, after passing to a
subsequence, the sequence $(M_k,g_k'(t),(q_k,0))$ converges
geometrically to a limiting Ricci flow
$(M_\infty,g_\infty(t),(q_\infty,0))$ consisting of non-negatively
curved, complete manifolds $\kappa$-non-collapsed on all scales
(though possibly with unbounded curvature).

Furthermore, Equation~(\ref{V/2}) passes to the limit to give
\begin{equation}\label{V3/2}
\Vol\,B_{g_\infty}(q_\infty,0,1)=\omega_3/2. \end{equation}
 Since $r_k/d_k\rightarrow 0$ as $k\rightarrow\infty$ and since $R_{g_k}(q_k,0)=d_k^{-2}$, we see that
$R_{g_\infty}(q_\infty,0)=0$. By the strong maximum principle for
scalar curvature (Theorem~\ref{firststrongmax}), this implies that
the limit $(M_\infty,g_\infty(0))$ is flat. But
Equation~(\ref{V3/2}) tells us that this limit is not $\Ar^3$. Since
it is complete and flat, it must be a quotient of $\Ar^3$ by an
action of a non-trivial group of isometries acting freely and
properly discontinuously. But the quotient of $\Ar^3$ by any
non-trivial group of isometries acting freely and properly
discontinuously has zero asymptotic volume. [Proof: It suffices to
prove the claim in the special case when the group is infinite
cyclic. The generator of this group has an axis $\alpha$ on which it
acts by translation and on the orthogonal subspace its acts by an
isometry. Consider the circle in the quotient that is the image of
$\alpha$, and let $L_\alpha$ be its length. The volume of the ball
of radius $r$ about $L_\alpha$ is $\pi r^{2}L_\alpha$. Clearly then,
for any $p\in \alpha$, the volume of the ball of radius $r$ about
$p$ is at most $\pi L_\alpha r^2$. This proves that the asymptotic
volume of the quotient is zero.]

We have now shown that $(M_\infty,g_\infty(0))$ has zero curvature and zero
asymptotic volume. But this implies that it is not $\kappa$-non-collapsed on
all scales, which is a contradiction. This contradiction completes the proof of
Claim~\ref{claim:curvature-ratio-bound-near-basepoint}.
\end{proof}

This claim establishes the existence of a universal constant $C<\infty$
(universal in the sense that it is independent of the $3$-dimensional
$\kappa$-solution) such that $R(q',0)\le CQ$ for all $q'\in B(q,0,2d)$. Since
the curvature of $(M,g(t))$ is non-negative and bounded, we know from the
Harnack inequality (Corollary~\ref{posderiv}) that  $R(q',t)\le CQ$ for all
$q'\in B(q,0,2d)$ and all $t\le 0$. Hence, the Ricci curvature $\Ric(q',t)\le CQ$ for all $q'\in B(q,0,2d)$ and all $t\le 0$.

\begin{claim}[Distance bound under rescaling]\label{claim:distance-bound-under-rescaling}
Given any constant $c>0$ there is a constant $\widetilde
C=\widetilde C(c)$, depending only on $c$ and not on the
$3$-dimensional $\kappa$-solution, so that
$$d_{g(-cQ^{-1})}(p,q)\le \tilde CQ^{-1/2}.$$
\end{claim}

\begin{proof}
Let $\gamma\colon[0,d]\to M$ be a $g(0)$-geodesic from $p$ to $q$,
parameterized at unit speed. Denote by $\ell_t(\gamma)$ the length
of $\gamma$ under the metric $g(t)$. We have $d_t(p,q)\le
\ell_t(\gamma)$. We estimate $\ell_t(\gamma)$ using the fact that
$|\Ric|\le CQ$ on the image of $\gamma$ at all times.
\begin{eqnarray*}
\frac{d}{dt}\ell_{t_0}(\gamma) & = &
\frac{d}{dt}\left(\int_0^d\sqrt{\langle\gamma'(s),\gamma'(s)\rangle_{g(t)}}ds\right)
\Bigl|_{t=t_0}\Bigr.
\\
& = & \int_0^{d}\frac{-\Ric_{g(t_0)}(\gamma'(s),\gamma'(s))}
{\sqrt{\langle\gamma'(s),\gamma'(s)\rangle_{g(t_0)}}}ds \\
& \ge & -CQ\int_0^d|\gamma'(s)|_{g(t_0)}ds=-CQ\ell_{t_0}(\gamma).
\end{eqnarray*}
Integrating yields
$$\ell_{-t}(\gamma)\le e^{CQt}\ell_0(\gamma)=e^{CQt}Q^{-1/2}.$$
(Recall $d^2Q=1$.) Plugging in $t=cQ^{-1}$ gives us
$$d_{-cQ^{-1}}(p,q)\le \ell_{-cQ^{-1}}(\gamma)\le e^{cC}Q^{-1/2}.$$
Setting $\widetilde C=e^{cC}$ completes the proof of the claim.
\end{proof}

The integrated form of  Hamilton's Harnack inequality,
Theorem~\ref{Harnack2}, tells us that
$$\log\left(\frac{R(p,0)}{R(q,-cQ^{-1})}\right)\ge
-\frac{d_{-cQ^{-1}}^2(p,q)}{2cQ^{-1}}.$$ According to the above
claim, this in turn tells us
$$\log\left(\frac{R(p,0)}{R(q,-cQ^{-1})}\right)\ge -\widetilde
C^2/2c.$$ Since $R(p,0)=1$, it immediately follows that
$R(q,-cQ^{-1})\le \exp(\widetilde C^2/(2c))$.

\begin{claim}[Universal upper bound for $Q$]\label{claim:curvature-upper-bound-universal}
\label{zkbound}
\uses{claim:distance-bound-under-rescaling}
There is a universal (i.e., independent of the $3$-dimensional
$\kappa$-solution) upper bound  for $Q=R(q,0)$.
\end{claim}



\begin{proof}
\uses{claim:distance-bound-under-rescaling}
Let $G'=QG$ and $t'=Qt$. Then $R_{G'}(q',0)\le C$ for all $q'\in
B_{G'}(q,0,2)$. Consequently, $R_{G'}(q',t')\le C$ for all $q'\in
B_{G'}(q,0,2)$ and all $t'\le 0$.  Thus, by Shi's derivative
estimates (Theorem~\ref{thm:shi-derivative-estimates})applied  with $T=2$ and $r=2$, there is
a universal constant $C_1$ such that for all $-1\le t'\le 0$
$$|\triangle \Rm_{G'}(q,t')|_{G'}\le C_1,$$
(where the Laplacian is taken with respect to the metric $G'$).
Rescaling by $Q^{-1}$ we see that for all $-Q^{-1}\le t\le 0$ we
have
$$|\triangle \Rm_{G}(q,t)|\le C_1Q^2,$$
where the Laplacian is taken with respect to the metric $G$. Since
the metric is non-negatively curved, by Corollary~\ref{posderiv} we
have  $2|\Ric(q,t)|^2\le 2Q^2$ for all $t\le 0$.
 From these two facts we conclude from
the flow equation~(\ref{Revol}) that there is a constant $1<C''<\infty$ with
the property that $\partial R/\partial t(q,t)\le C''Q^2$ for all $-Q^{-1}<t\le
0$.
 Thus for any $0<c<1$, we have $Q=R(q,0)\le
cC''Q+R(q,-cQ^{-1})\le cC''Q+e^{(\widetilde C^2(c)/2c)}$. Now we
take $c=(2C'')^{-1}$ and $\widetilde C=\widetilde C(c)$. Plugging
these values into the previous inequality yields
$$Q\le 2e^{(\tilde C^2C'')}.$$
\end{proof}


This leads immediately to:

\begin{claim}[Universal neighborhood curvature bound]\label{claim:basepoint-neighborhood-curvature-bound}
\label{bdnearxk}
\uses{claim:curvature-ratio-bound-near-basepoint,claim:curvature-upper-bound-universal}
There are universal constants $\delta>0$ and  $C_1<\infty$
(independent of the based $3$-dimensional $\kappa$-solution
$(M,g(t),(p,0))$ with $R(p,0)=1$) such that $d(p,q)\ge \delta$. In
addition,  $R(q',t)\le C_1$ for all $q'\in B(p,0,d)$ and all $t\le
0$.
\end{claim}


\begin{proof}
Since, according to the previous claim, $Q$ is universally bounded
above and $d^2Q=1$, the existence of $\delta>0$ as required is
clear. Since $B(p,0,d)\subset B(q,0,2d)$, since $R(q',0)/R(q,0)$ is
universally bounded on $B(q,0,2d)$ by Claim~\ref{claim:curvature-ratio-bound-near-basepoint}, and since
$R(q,0)$ is universally bounded by Lemma~\ref{claim:curvature-upper-bound-universal}, the second
statement is clear for all $(q',0)\in B(p,0,d)\subset B(q,0,2d)$.
Given this, the fact that the second statement holds for all
$(q',t)\in B(p,0,d)\times(-\infty,0]$ then follows immediately from
the derivative inequality for $\partial R(q,t)/\partial t$,
Corollary~\ref{posderiv}.
\end{proof}

This, in turn, leads immediately to:

\begin{corollary}[Curvature bound on a fixed universal ball]\label{cor:curvature-bound-fixed-ball}
\label{rho}
\uses{claim:basepoint-neighborhood-curvature-bound}
Fix $\delta>0$ the universal constant of the last
claim. Then $R(q',t)\le \delta^{-2}$ for all $q'\in B(p,0,\delta)$
and all $t\le 0$.
\end{corollary}

Now we return to the proof of Lemma~\ref{lem:uniform-curvature-bound-kappa-solution}. Since
$(M,g(t))$ is $\kappa$-non-collapsed, it follows from the previous
corollary that $\Vol\,B(p,0,\delta)\ge \kappa \delta^3$. Hence,
according to Corollary~\ref{cor:global-curvature-bound-from-volume} for each $A<\infty$ there is a
constant $K(A)$ such that $R(q',0)\le K(A/\delta)/(\delta+A)^2$ for
all $q'\in B(p,0,A)$. Since $\delta$ is a universal positive
constant, this completes the proof of Lemma~\ref{lem:uniform-curvature-bound-kappa-solution}.
\end{proof}

Now let us turn to the proof of Theorem~\ref{thm:kappa-solution-compactness}, the
compactness result for $\kappa$-solutions.

\begin{proof}
\uses{lem:uniform-curvature-bound-kappa-solution,cor:three-dimensional-asymptotic-soliton-bounded-curvature}
 Let
$(M_k,g_k(t),(p_k,0))$ be a sequence of based $3$-dimensional
$\kappa$-solutions with $R(p_k,0)=1$ for all $k$.
 The immediate
consequence of Lemma~\ref{lem:uniform-curvature-bound-kappa-solution} and Corollary~\ref{posderiv}
is the following. For every $r<\infty$ there is a constant
$C(r)<\infty$ such that $R(q,t)\le C(r)$ for all $q\in B(p_k,0,r)$
and for all $t\le 0$. Of course, since, in addition, the elements in
the sequence are $\kappa$-non-collapsed, by Theorem~\ref{flowlimit}
this implies that there is a subsequence of the
$(M_k,g_k(t),(p_k,0))$ that converges geometrically to an ancient
flow $(M_\infty,g_\infty(t),(p_\infty,0))$. Being a geometric limit
of $\kappa$-solutions, this limit is complete and
$\kappa$-non-collapsed, and each time-slice is of non-negative
curvature. Also, it is not flat since, by construction,
$R(p_\infty,0)=1$. Of course, it also follows from the limiting
procedure that $\partial R(q,t)/\partial t\ge 0$ for every $(q,t)\in
M_\infty\times(-\infty,0]$. Thus, according to
Corollary~\ref{cor:three-dimensional-asymptotic-soliton-bounded-curvature} the limit
 $(M_\infty,g_\infty(t))$ has bounded curvature for each $t\le 0$.
Hence, the limit is a $\kappa$-solution. This completes the proof of
Theorem~\ref{thm:kappa-solution-compactness}.
\end{proof}

\begin{corollary}[Universal scale-invariant gradient and time-derivative bounds]\label{cor:curvature-derivative-bounds-universal}
\label{derivcor}
\uses{thm:kappa-solution-compactness}
Given $\kappa>0$, there is $C<\infty$ such that for any
$3$-dimensional $\kappa$-solution $(M,g(t)),\ -\infty<t\le 0$, we
have
\begin{eqnarray}\label{gradest} \sup_{(x,t)}\frac{|\nabla R(x,t)|}{R(x,t)^{3/2}}< C\\
\label{timeest} \sup_{(x,t)}\frac{\left|\frac{d}{dt}R(x,t)\right|}{R(x,t)^2}<C.
\end{eqnarray}

\end{corollary}



\begin{proof}
\uses{thm:kappa-solution-compactness}
Notice that the two inequalities are scale invariant. Thus, this
result is immediate from the compactness theorem,
Theorem~\ref{thm:kappa-solution-compactness}.
\end{proof}

Because of Proposition~\ref{prop:universal-kappa-non-round}, and the fact that the previous
corollary obviously holds for any shrinking family of round metrics,
we can take the constant $C$ in the above corollary to be
independent of $\kappa>0$.

Notice that, using Equation~(\ref{Revol}), we can rewrite the second
inequality in the above corollary as
\begin{eqnarray*}
\sup_{(x,t)}\frac{|\triangle R+2|\Ric|^2|}{R(x,t)^2}<C.
\end{eqnarray*}



\section{Qualitative description of
$\kappa$-solutions}

In Chapter~\ref{nonnegcurv} we defined the notion of an
$\epsilon$-neck. In this section we define a stronger version of
these, called strong $\epsilon$-necks. We also introduce other types
of canonical neighborhoods -- $\epsilon$-caps, $\epsilon$-round
components and $C$-components. These definitions pave the way for a
qualitative description of $\kappa$-solutions.



\subsection{Strong canonical neighborhoods}

The next manifold we introduce is one with controlled topology
(diffeomorphic either to the $3$-disk or a punctured $\Ar P^3$) with
the property that the complement of a compact submanifold is an
$\epsilon$-neck.

\begin{definition}[$(C,\epsilon)$-cap]\label{def:c-epsilon-cap}
\label{epscap}
\uses{def:epsilon-neck}
Fix constants $0<\epsilon<1/2$ and $C<\infty$. Let $(M,g)$ be a
Riemannian $3$-manifold. A $(C,\epsilon)$-cap in $(M,g)$ is an open
submanifold $({\mathcal C},g|_{\mathcal C})$ together with an open
submanifold $N\subset {\mathcal C}$ with the following properties:
\begin{enumerate}
\item[(1)] ${\mathcal C}$ is diffeomorphic either to an open $3$-ball or to a
punctured $\Ar P^3$.
\item[(2)] $N$ is an $\epsilon$-neck with compact complement in
${\mathcal C}$.
\item[(3)] $\overline Y={\mathcal C}\setminus N$ is a compact submanifold with boundary.
Its  interior, $Y$, is called
 the {\em core} of ${\mathcal C}$. The frontier of $Y$, which is $\partial \overline Y$,
  is a central
 $2$-sphere of an $\epsilon$-neck contained in ${\mathcal C}$.
\item[(4)] The scalar curvature $R(y)>0$ for every $y\in {\mathcal C}$ and
$$ \diam({\mathcal C},g|_{\mathcal C})< C \left(\sup_{y\in
{\mathcal C}}R(y)\right)^{-1/2}.$$
\item[(5)] $\sup_{x,y\in {\mathcal C}}\left[ R(y)/R(x)\right]< C$.
\item[(6)]  $\Vol\,{\mathcal C} < C(\sup_{y\in {\mathcal C}}R(y))^{-3/2}$.
\item[(7)] For any $y\in Y$  let $r_y$ be defined so
that $\sup_{y'\in B(y,r_y)}R(y')=r_y^{-2}$. Then for each $y\in
Y$, the ball $B(y,r_y)$ lies in ${\mathcal C}$ and indeed has
compact closure in ${\mathcal C}$. Furthermore,
$$C^{-1}<\inf_{y\in
Y}\frac{\Vol\,B(y,r_y)}{r_y^3}.$$
\item[(8)] Lastly,
$$\sup_{y\in {\mathcal C}}\frac{|\nabla R(y)|}{R(y)^{3/2}}<C$$
and
$$\sup_{y\in {\mathcal C}}\frac{\left|\triangle R(y)+2|\Ric|^2\right|}{R(y)^2}<C.$$
\end{enumerate}
\end{definition}

\begin{remark}[The ball $B(y,r_y)$ has compact closure in the cap]\label{rem:cap-core-ball-remark}
\uses{def:c-epsilon-cap}
If the ball $B(y,r_y)$ meets the complement of the core of
${\mathcal C}$ then it contains a point whose scalar curvature is
close to $R(x)$, and hence $r_y$ is bounded above by, say
$2R(x)^{-1}$. Since $\epsilon<1/2$, using the fact that $y$ is
contained in the core of ${\mathcal C}$ it follows that $B(y,r_y)$
is contained in ${\mathcal C}$ and has compact closure in ${\mathcal
C}$.

Implicitly, we always orient the $\epsilon$-neck structure on $N$ so
that the closure of its negative end meets the core of ${\mathcal
C}$.
\end{remark}

Condition (8) in the above definition may seem unnatural, but here
is the reason for it.

\begin{claim}[Condition (8) is a time-derivative bound]\label{claim:cap-condition-time-derivative-equivalence}
\label{partialtR}
\uses{def:c-epsilon-cap,def:ricci-flow}
Suppose that $(M,g(t))$ is a Ricci flow and that $\left({\mathcal
C},g(t)|_{\mathcal C}\right)$ is a subset of a $t$ time-slice. Then
Condition (8) above is equivalent to
$$\sup_{(x,t)\in {\mathcal C}}\frac{\left|\frac{\partial R(x,t)}{\partial
t}\right|}{R^2(x,t)}<C.$$
\end{claim}

\begin{proof}
This is immediate from Equation~(\ref{Revol}).
\end{proof}


Notice that the definition of a $(C,\epsilon)$-cap is a scale
invariant notion.

\begin{definition}[$C$-component]\label{def:c-component}
\label{Ccomponent}
Fix a positive constant $C$. A compact connected Riemannian manifold
$(M,g)$ is called {\em a $C$-component} if
\begin{enumerate}
\item[(1)] $M$ is diffeomorphic to either $S^3$ or $\Ar P^3$.
\item[(2)] $(M,g)$ has positive sectional curvature.
\item[(3)] $$C^{-1}<\frac{\inf_{P}K(P)}{\sup_{y\in M}R(y)}$$
where $P$ varies over all $2$-planes in $TX$ (and $K(P)$ denotes the
sectional curvature in the $P$-direction).
\item[(4)] $$C^{-1}\sup_{y\in M}\left(R(y)^{-1/2}\right)<\diam(M)<C\inf_{y\in M}\left(R(y)^{-1/2}\right).$$
\end{enumerate}
\end{definition}


\begin{definition}[$\epsilon$-round component]\label{def:epsilon-round-component}
\label{defnepsround}
Fix $\epsilon>0$. Let $(M,g)$ be a compact, connected $3$-manifold.
Then $(M,g)$ is {\em within $\epsilon$ of round in the
$C^{[1/\epsilon]}$-topology} if there exist a constant $R>0$, a
compact manifold $(Z,g_0)$ of constant curvature $+1$, and a
diffeomorphism $\varphi\colon Z\to M$ with the property that the
pull back under $\varphi$ of $Rg$ is within $\epsilon$ in the
$C^{[1/\epsilon]}$-topology of $g_0$.
\end{definition}


Notice that both of these notions are scale invariant notions.

\begin{definition}[$(C,\epsilon)$-canonical neighborhood]\label{def:canonical-neighborhood}
\uses{def:epsilon-neck,def:c-epsilon-cap,def:c-component,def:epsilon-round-component}
Fix $C<\infty$ and $\epsilon>0$. For any Riemannian manifold
$(M,g)$,  an open neighborhood $U$ of a point $x\in M$ is a {\em
$(C,\epsilon)$-canonical neighborhood} if one of the following holds:
\begin{enumerate}
\item[(1)] $U$ is an $\epsilon$-neck in $(M,g)$ centered at $x$.
\item[(2)] $U$ a $(C,\epsilon)$-cap in $(M,g)$ whose core contains $x$.
\item[(3)] $U$ is a $C$-component of $(M,g)$.
\item[(4)] $U$ is  an $\epsilon$-round component of $(M,g)$.
\end{enumerate}
\end{definition}

Whether or not a point $x\in M$ has a $(C,\epsilon)$-canonical
neighborhood in $M$ is a scale invariant notion.

The notion of $(C,\epsilon)$-canonical neighborhoods is sufficient
for some purposes, but often we need a stronger notion.

\begin{definition}[Evolving and strong $\epsilon$-neck; strong canonical neighborhood]\label{def:strong-epsilon-neck-canonical-neighborhood}
\label{defncanonnbhd}
\uses{def:generalized-ricci-flow,def:epsilon-neck,def:c-epsilon-cap,def:c-component,def:epsilon-round-component,def:epsilon-close-topology}
Fix constants $C<\infty$ and $\epsilon>0$. Let $({\mathcal M},G)$ be
a generalized Ricci flow. An {\em evolving  $\epsilon$-neck defined
for an interval of normalized time of length
$t'>0$} centered at a point $x\in
{\mathcal M}$ with ${\bf t}(x)=t$ is an embedding $\psi\colon
S^2\times (-\epsilon^{-1},\epsilon^{-1})\stackrel{\cong}{\to}
N\subset M_t$ with $x\in \psi(S^2\times\{0\})$ satisfying the
following properties:
\begin{enumerate}
\item[(1)] There is an embedding $N\times (t-R(x)^{-1}t',t]\to {\mathcal
M}$ compatible with time and the vector field.
\item[(2)] The pullback under $\psi$ of the one-parameter family of
metrics on $N$ determined by restricting $R(x)G$ to the image of
this embedding is within $\epsilon$ in the
$C^{[1/\epsilon]}$-topology of the standard family $(h(t),ds^2),
-t'<t\le 0$, where $h(t)$ is the round metric of scalar curvature
$1/(1-t)$ on $S^2$ and $ds^2$ is the usual Euclidean metric on the
interval (see Definition~\ref{epsclose} for the notion of two
families of metrics being close).
\end{enumerate}
A {\em strong $\epsilon$-neck} is
the image of an evolving $\epsilon$-neck which is defined for an
interval of normalized time of length  $1$.

Both of these notions are scale invariant notions.

Let $({\mathcal M},G)$ be a generalized Ricci flow. Let $x\in
{\mathcal M}$ be a point with ${\bf t}(x)=t$. We say that an open
neighborhood $U$ of $x$ in $M_t$  is a {\em strong
$(C,\epsilon)$-canonical neighborhood} of $x$ if one of the following holds
\begin{enumerate}
\item[(1)]   $U$ is a strong $\epsilon$-neck in $({\mathcal M},G)$ centered at $x$.
\item[(2)] $U$ is a $(C,\epsilon)$-cap in $M_t$ whose  core contains $x$.
\item[(3)] $U$ is a $C$-component of $M_t$.
\item[(4)] $U$ is an $\epsilon$-round component of $M_t$.
\end{enumerate}
\end{definition}

Whether or not a point $x$ in a generalized Ricci flow has a strong
$(C,\epsilon)$-canonical neighborhood is a scale invariant notion.


\begin{proposition}[Continuity of canonical neighborhoods]\label{prop:canonical-neighborhood-stability}
\label{canonvary}
\uses{def:strong-epsilon-neck-canonical-neighborhood,def:canonical-neighborhood}
The following holds for any  $\epsilon<1/4$ and any $C<\infty$. Let
$({\mathcal M},G)$ be a generalized Ricci flow and let $x\in
{\mathcal M}$ be a point with ${\bf t}(x)=t$.

(1) Suppose that $U\subset M_t$ is a  $(C,\epsilon)$-canonical
neighborhood for $x$. Then for any horizontal metric $G'$
sufficiently close to $G|_U$ in the $C^{[1/\epsilon]}$-topology,
$(U,G'|_U)$ is a  $(C,\epsilon)$ neighborhood for any $x'\in U$
sufficiently close to $x$.



(2) Suppose that in $({\mathcal M},G)$ there is an evolving
$\epsilon$-neck $U$ centered at $(x,t)$ defined for an interval of
normalized time of length  $a>1$. Then any Ricci flow on $U\times
(t-aR(x,t)^{-1},t]$ sufficiently close in the
$C^{[1/\epsilon]}$-topology to the pullback of $G$ contains a strong
$\epsilon$-neck centered at $(x,t)$.

(3) Given $(C,\epsilon)$ and $(C',\epsilon')$ with $C'>C$ and
$\epsilon'>\epsilon$ there is $\delta>0$ such that the following
holds. Suppose that $R(x)\le 2$. If $(U,g)$ is a
$(C,\epsilon)$-canonical neighborhood of $x$ then for any metric
$g'$ within $\delta$ of $g$ in the $C^{[1/\epsilon]}$-topology
$(U,g')$ contains a $(C',\epsilon')$-neighborhood of $x$.

(4) Suppose that $g(t),\ -1<t\le 0$, is a one-parameter family of
metrics on $(U,g)$ that is a strong $\epsilon$-neck centered at
$(x,0)$ and $R_g(x,0)=1$. Then any one-parameter family $g'(t)$
within $\delta$ in the $C^{[1/\epsilon]}$-topology of $g$ with
$R_{g'}(x,0)=1$ is a strong $\epsilon'$-neck.
\end{proposition}


\begin{proof}
Since $\epsilon<1/4$,  the diameter of $(U,g)$, the volume of
$(U,g)$, the supremum over $x\in U$ of $R(x)$, the supremum over $x$
and $y$ in $U$ of $R(y)/R(x)$, and the infimum over all $2$-planes
$P$ in ${\mathcal H}TU$ of $K(P)$ are all continuous functions of
the horizontal metric $G$ in the $C^{[1/\epsilon]}$-topology.


Let us consider the first statement. Suppose $(U,G|_U)$ is a
$C$-component or an $\epsilon$-round component. Since the defining
inequalities are strict, and, as we just remarked, the quantities in
these inequalities vary continuously with the metric in the
$C^{[1/\epsilon]}$-topology the result is clear in this case.

Let us consider the case when $U\subset M_t$ is an $\epsilon$-neck
centered at $x$. Let $\psi\colon S^2\times
(-\epsilon^{-1},\epsilon^{-1})\to U$ be the map giving the
$\epsilon$-neck structure. Then for all horizontal metrics $G'$
sufficiently close to $G$ in the $C^{[1/\epsilon]}$-topology the
same map $\psi$ is determines an $\epsilon$-neck centered at $x$ for
the structure $(U,G'|_U)$. Now let us consider moving $x$ to a
nearby point $x'$, say $x'$ is the image of $(a,s)\in S^2\times
(-\epsilon^{-1},\epsilon^{-1})$. We pre-compose $\psi$ by a map
which is the product of the identity in the $S^2$-factor with a
diffeomorphism $\alpha$ on $(-\epsilon^{-1},\epsilon^{-1})$ that is
the identity near the ends and moves $0$ to $s$. As $x'$ approaches
$x$, $s$ tends to zero, and hence we can choose $\alpha$ so that it
tends to the identity in the $C^\infty$-topology. Thus, for $x'$
sufficiently close to $x$, this composition will determine an
$\epsilon$-neck structure centered at $x'$. Lastly, let us consider
the case when $(U,G|_U)$ is a $(C,\epsilon)$-cap whose core $Y$
contains $x$. Let $G'$ be a horizontal metric sufficiently close to
$G|_U$ in the $C^{[1/\epsilon]}$-topology. Let $N\subset U$ be the
$\epsilon$-neck $U\setminus \bar Y$. We have just seen that
$(N,G'|_N)$ is an $\epsilon$-neck. Similarly, if $N'\subset U$ is an
$\epsilon$-neck with central $2$-sphere $\partial \overline Y$, then
$(N',G'|_{N'})$ is an $\epsilon$-neck if $G'$ is sufficiently close
to $G$ in the $C^{[1/\epsilon]}$-topology.


Thus, Conditions (1),(2), and (3) in the definition of a
$(C,\epsilon)$-cap hold for $(U,G'_U)$.
 Since the
curvature, volume and diameter inequalities in Conditions (4), (5),
and (6) are strict, they also hold for $g'$. To verify that
Condition (7) holds for $G'$, we need only remark that $r_y$ is a
continuous function of the metric. Lastly, since the derivative
inequalities for the curvature in Condition (8) are strict
inequalities and $\epsilon^{-1}>4$, if these inequalities hold for
all horizontal metrics $G'$ sufficiently close to $G$ in the
$C^{[1/\epsilon]}$-topology. This completes the examination of all
cases and proves the first statement.

The second statement is proved in the same way using the fact that
if $g'(t)$ is sufficiently close to $g$ in the
$C^{[1/\epsilon]}$-topology and if $x'$ is sufficiently close to $x$
then $R_{g'(x')}^{-1}<aR_g(x)^{-1}$.


Now let us turn to the third statement. The result is clear for
$\epsilon$-necks. Also, since $R(x)\le 2$ the result is clear for
$\epsilon$-round components and $C$-components as well. Lastly, we
consider a $(C,\epsilon)$-cap $U$ whose core $Y$ contains $x$.
 Clearly,
since $R(x)$ is bounded above by  $2$, for $\delta>0$ sufficiently
small, any metric $g'$ within $\delta$ of $g$ will satisfy the
diameter, volume and curvature and the derivative of the curvature
inequalities with $C'$ replacing $C$. Let $N$ be the $\epsilon$-neck
in $(U,g)$ containing the end of $U$. Assuming that $\delta$ is
sufficiently small, let $N'$ be the image of $S^2\times
\left(-\epsilon^{-1},2(\epsilon')^{-1}-\epsilon^{-1}\right)$. Then
$(N',g')$ becomes an $\epsilon'$-neck structure once we shift the
parameter in the $s$-direction by $\epsilon^{-1}-(\epsilon')^{-1}$.
We let $U'=\overline Y\cup N'$. Clearly, the $\epsilon$-neck with
central $2$-sphere $\partial \overline Y$ will also determine an
$\epsilon'$-neck with the same central $2$-sphere provided that
$\delta>0$ is sufficiently small. Thus, for $\delta>0$ sufficiently
small, for any $(C,\epsilon)$ the result of this operation is a
$(C',\epsilon')$-cap with the same core.


The fourth statement is immediate.
\end{proof}

\begin{corollary}[Centers of strong necks form an open set]\label{cor:strong-neck-centers-open}
\label{strepsopen}
\uses{prop:canonical-neighborhood-stability}
In an ancient solution $(M,g(t))$ the set of points that are centers
of strong $\epsilon$-necks is an open subset
\end{corollary}

\begin{proof}
\uses{prop:canonical-neighborhood-stability}
Let $T$ be the final time of the flow. Suppose that $(x,t)$ is the
center of a strong $\epsilon$-neck $U\times (t-R(x,t)^{-1},t]\subset
M\times (-\infty,0]$. This neck extends backwards for all time and
forwards until the final time $T$ giving an embedding $U\times
(-\infty,T]\to M\times (-\infty,T]$. There is $a>1$ such that for
all $t'$ sufficiently close to $t$ the restriction of this embedding
determines an evolving $\epsilon$-neck centered at $(x,t')$ defined
for an interval of normalized time of length  $a$. Composing this
neck structure with a self-diffeomorphism of $U$ moving $x'$ to $x$,
as described above, shows that all $(x',t')$ sufficiently close to
$(x,t)$ are centers of strong $\epsilon$-necks.
\end{proof}

\begin{definition}[$\epsilon$-tube and capped $\epsilon$-tube]\label{def:epsilon-tube}
\label{epstube}
\uses{def:epsilon-neck,def:c-epsilon-cap}
{\em An $\epsilon$-tube ${\mathcal T}$} in
a Riemannian $3$-manifold $M$ is a submanifold diffeomorphic to the
product of $S^2$ with a non-degenerate interval with the following
properties:
\begin{enumerate}
\item[(1)] Each boundary component $S$ of ${\mathcal T}$ is the central
$2$-sphere of an $\epsilon$-neck $N(S)$ in $M$.
\item[(2)] ${\mathcal T}$
is a union of $\epsilon$-necks and the closed half $\epsilon$-necks
whose boundary sphere is a component of $\partial{\mathcal T}$.
Furthermore, the central $2$-sphere of each of the $\epsilon$-necks
is isotopic in ${\mathcal T}$ to the $S^2$-factors of the product
structure.
\end{enumerate}

An {\em open $\epsilon$-tube} is one without boundary. It is a union
of $\epsilon$-necks with the central spheres that are isotopic to
the $2$-spheres of the product structure.


A {\em $C$-capped $\epsilon$-tube}
in $M$ is a connected submanifold that is the union of a
$(C,\epsilon)$-cap ${\mathcal C}$
 and an open $\epsilon$-tube where the intersection of ${\mathcal C}$ with the
$\epsilon$-tube   is diffeomorphic to $S^2\times (0,1)$ and contains
an end of the $\epsilon$-tube and an end of the cap.
 A {\em doubly $C$-capped
$\epsilon$-tube} in $M$ is a
closed, connected submanifold of $M$ that is the union of two
$(C,\epsilon)$-caps ${\mathcal C}_1$ and ${\mathcal C}_2$  and  an
open $\epsilon$-tube. Furthermore, we require (i) that the cores
$Y_1$ and $Y_2$ of ${\mathcal C}_1$ and ${\mathcal C}_2$ have
disjoint closures, (ii) that the union of either ${\mathcal C}_i$
with the $\epsilon$-tube is  a capped $\epsilon$-tube and ${\mathcal
C}_1$ and ${\mathcal C}_2$ contain the opposite ends of the
$\epsilon$-tube. There is one further closely related notion, that
of an {\em $\epsilon$-fibration}. By
definition an $\epsilon$-fibration is a closed, connected manifold
that fibers over the circle with fibers $S^2$ that is also a union
of $\epsilon$-necks with the property that the central $2$-sphere of
each neck is isotopic to a fiber of the fibration structure. We
shall not see this notion again until the appendix, but because it
is clearly closely related to the notion of an $\epsilon$-tube, we
introduce it here.
\end{definition}

\begin{definition}[Strong $\epsilon$-tube]\label{def:strong-epsilon-tube}
\uses{def:epsilon-tube,def:strong-epsilon-neck-canonical-neighborhood}
A {\em strong $\epsilon$-tube} in a
generalized Ricci flow is an $\epsilon$-tube with the property that
each point of the tube is the center of a strong $\epsilon$-neck in
the generalized flow.
\end{definition}

\subsection{Canonical neighborhoods for $\kappa$-solutions}


\begin{proposition}[Trichotomy for $3$-dimensional $\kappa$-solutions]\label{prop:kappa-solution-trichotomy}
\label{infprod2}
\uses{def:ancient-solution,def:kappa-noncollapsed}
Let $(M,g(t))$ be a  $3$-dimensional $\kappa$-solution.  Then one of
the following hold:
\begin{enumerate}
\item[(1)] For every $t\le 0$ the manifold $(M,g(t))$ has positive
curvature.
\item[(2)] $(M,g(t))$ is the product of an evolving family of
round $S^2$'s with a line.
\item[(3)] $M$ is diffeomorphic to a line  bundle over $\Ar P^2$, and
there is a finite covering of $(M,g(t))$ that is  a flow as in (2).
\end{enumerate}
\end{proposition}

\begin{proof}
Suppose that $(M,g(t))$ does not have positive curvature for some
$t$. Then, by the application of the strong maximum principle given
in Corollary~\ref{nullspace}, there is a covering $\widetilde M$ of
$M$, with either one or $2$-sheets, such that $(\widetilde M,\tilde
g(t))$ is the product of an evolving family of round surfaces with a
flat one-manifold (either a circle or the real line).  Of course,
the covering must be $\kappa$-solution. In the case in which
$(\widetilde M,\tilde g(t))$ is isometric to the product of an
evolving family of round surfaces  and a circle, that circle has a
fixed length, say $L<\infty$. Since the curvature of the surface in
the $t$ time-slice goes to zero as $t\rightarrow -\infty$, we see
that the flow is not $\kappa$-non-collapsed on all scales for any
$\kappa>0$. Thus, $(M,g(t))$ has either a trivial cover or a double
cover isometric to the product of a shrinking family of round
surfaces with $\Ar$. If the round surface is $S^2$, then we have
established the result. If the round surface is $\Ar P^2$ a further
double covering is a product of round two-spheres with $\Ar$.
 This proves the proposition.
\end{proof}

\begin{lemma}[Neck structure outside a bounded ball]\label{lem:noncompact-kappa-solution-neck-structure}
\uses{prop:kappa-solution-trichotomy}
Let $(M,g(t))$ be a non-compact $3$-dimensional $\kappa$-solution of
positive curvature and let $p\in M$. Then there is $D'<\infty$,
possibly depending on $(M,g(0))$ and $p$, such that every point of
$M\times\{0\}\setminus B(p,0,D'R(p,0)^{-1/2})$ is the center of an
evolving $\epsilon$-neck in $(M,g(t))$ defined for an interval of
normalized time of length  $2$. Furthermore, there is $D'_1<\infty$
such that for any point $x\in B(p,0,D'R(p,0)^{-1/2})$ and any
$2$-plane $P_x$ in $T_xM$
 we have $(D'_1)^{-1}< K(P_x)/R(p,0)< D'_1$ where $K(P_x)$ denotes the sectional
 curvature in the direction of the $2$-plane $P_x$.
\end{lemma}

\begin{proof}
\uses{thm:kappa-solution-compactness,cor:two-dimensional-gradient-shrinking-soliton,lem:uniform-curvature-bound-kappa-solution}
 Given $(M,g(t))$ and $p$, suppose that no such $D'<\infty$ exists.
Because the statement is scale invariant, we can arrange that
$R(p,0)=1$.
 Then we
can find a sequence of points $p_k\in M$ with $d_0(p,p_k)\rightarrow
\infty$ as $k\rightarrow \infty$ such that no $p_k$ is the center of
an evolving $\epsilon$-neck in $(M,g(0))$ defined for an interval of
normalized time of length  $2$. By passing to a subsequence we can
assume that one of two possibilities holds: either
$d_0^2(p,p_k)R(p_k,0)\rightarrow\infty$ as $k\rightarrow\infty$ or
$\lim_{k\rightarrow\infty}d_0^2(p,p_k)R(p_k,0)=\ell<\infty$. In
the first case,  set $\lambda_k=R(p_k,0)$ and consider the based
flows $(M,\lambda_kg(\lambda_k^{-1}t),(p_k,0))$. According to
Theorem~\ref{thm:kappa-solution-compactness}, after passing to a subsequence there is
a geometric limit. Thus, by Theorem~\ref{topsplit} and
Corollary~\ref{localprod} the limit splits as a product of a
$2$-dimensional $\kappa$-solution and $\Ar$. By
Corollary~\ref{cor:two-dimensional-gradient-shrinking-soliton} it follows that the limit is the standard
round evolving cylinder. This implies that for all $k$ sufficiently
large $(p_k,0)$ is the center of an evolving $\epsilon$-neck in
$(M,g(t))$ defined for an interval of normalized time of length
$2$. This contradiction establishes the existence of $D'$ as
required in this case.

Now suppose that $\lim_{k\rightarrow\infty}d_0^2(p,p_k)R(p_k,0)=\ell<\infty$. Of
course, since $d_0(p,p_k)\rightarrow\infty$, it must be the case
that $R(p_k,0)\rightarrow 0$ as $k\rightarrow \infty$. Set
$Q_k=R(p_k,0)$. By passing to a subsequence we can arrange that
$d_0^2(p,p_k)Q_k< \ell+1$ for all $k$. Consider the
$\kappa$-solutions $(M_k,g_k(t))=(M,Q_kg(Q_k^{-1}t))$. For each $k$
we have $p\in B_{g_k}(p_k,0,\ell+1)$, and
$R_{g_k}(p,0)=Q_k^{-1}\rightarrow\infty$ as $k\rightarrow\infty$.
This contradicts Lemma~\ref{lem:uniform-curvature-bound-kappa-solution}, and completes the proof of
the existence of $D$ as required in this case as well.



 The existence of $D'_1$ is immediate since the
closure of the ball is compact and the manifold has positive
curvature.
\end{proof}

In fact a much stronger result is true. The constants $D'$ and
$D_1'$ in the above lemma can be chosen independent of the
non-compact $\kappa$-solutions.


\begin{proposition}[Uniform neck-structure constants]\label{prop:uniform-neck-structure-constants}
\label{uniformD}
\uses{lem:noncompact-kappa-solution-neck-structure}
For any $0<\epsilon$ sufficiently small there are constants
 $D=D(\epsilon)<\infty$ and $D_1=D_1(\epsilon)<\infty$
 such that the following holds for  any non-compact $3$-dimensional
 $\kappa$-solution $(M,g(t))$ of positive
curvature. Let $p\in M$ be a soul of $(M,g(0))$. Then:
\begin{enumerate}
\item[(1)] Every point in $M\setminus B(p,0,DR(p,0)^{-1/2})$ is the center
of a strong $\epsilon$-neck in $(M,g(t))$. Furthermore, for any
$x\in B(p,0,DR(p,0)^{-1/2})$ and any $2$-plane $P_x$ in $T_xM$ we
have
$$D_1^{-1}<K(P_x)/R(p,0)<D_1.$$
Also,
$$D_1^{-3/2}R(p,0)^{-3/2}< \Vol(B(p,0,DR(p,0)^{-1/2})<D_1^{3/2}R(p,0)^{-3/2}.$$
 \item[(2)] Let $f$
denote the distance function from $p$. For any $\epsilon$-neck
$N\subset (M,g(0))$, the middle two-thirds of $N$ is disjoint from
$p$, and the central $2$-sphere $S_N$ of $N$ is (topologically)
isotopic in $M\setminus \{p\}$ to $f^{-1}(a)$ for any $a>0$. In
particular, given two disjoint central $2$-spheres of
 $\epsilon$-necks in $(M,g(0))$ the region of $M$ bounded by these $2$-spheres is
 diffeomorphic
 to $S^2\times [0,1]$.
 \end{enumerate}
\end{proposition}

\begin{remark}[Basepoint may be replaced]\label{rem:uniform-neck-structure-basepoint-remark}
\uses{prop:uniform-neck-structure-constants}
In Part 1 of this theorem one can  replace $p$ by any point $p'\in
M$ that is not the center of a strong $\epsilon$-neck.
\end{remark}


\begin{proof}
\uses{thm:kappa-solution-compactness,lem:noncompact-kappa-solution-neck-structure,prop:canonical-neighborhood-stability,prop:kappa-solution-trichotomy,cor:strong-neck-centers-open}
First suppose that no $D$ exists so that the first statement holds.
Then there is a sequence of such solutions $(M_k,g_k(t))$, with
$p_k\in M_k$ being a soul of $(M_k,g_k(0))$ and points $q_k\in M_k$
with $d^2_0(p_k,q_k)R(p_k,0)\rightarrow\infty$ as
$k\rightarrow\infty$ such that $q_k$ is not the center of a strong
$\epsilon$-neck in $(M_k,g_k(0))$. By rescaling we can assume that
$R(p_k,0)=1$ for all $k$, and hence that
$d_0(p_k,q_k)\rightarrow\infty$. Then, according to
Theorem~\ref{thm:kappa-solution-compactness}, by passing to a subsequence we can
assume that there is a geometric limit
$(M_\infty,g_\infty(t),(p_\infty,0))$ with $R(p_\infty,0)=1$.
 By
Lemma~\ref{topiso}, provided that $\epsilon$ is sufficiently small
for all $k$ the soul $(p_k,0)$ is not the center of a strong
$2\epsilon$-neck in $(M_k,g_k(t))$.  Hence, invoking Part 4 of
Proposition~\ref{prop:canonical-neighborhood-stability} and using the fact that $R(p_k,0)=1$ for
all $k$ and hence $R(p_\infty,0)=1$, we see that $(p_\infty,0)$ is
not the center of a strong $\epsilon$-neck in
$(M_\infty,g_\infty(t))$. Since the manifolds $M_k$ are non-compact
and have metrics of positive curvature they are diffeomorphic to
$\Ar^3$ and in particular, do not contain embedded copies of $\Ar
P^2$. Thus, the limit $(M_\infty,g_\infty(t))$ is a non-compact
$\kappa$-solution containing no embedded copy of $\Ar P^2$. Thus, by
Proposition~\ref{prop:kappa-solution-trichotomy} either it is positively curved or it is a
Riemannian product $S^2$ times $\Ar$. In the second case every point
is the center of a strong $\epsilon$-neck. Since we have seen that
the point $(p_\infty,0)$ is not the center of a strong
$\epsilon$-neck, it follows that the limit is a positively curved
$\kappa$-solution.

 Then according to the previous lemma there is $D'$, depending on
$(M_\infty,g_\infty(0))$ and $p_\infty$, such that every point
outside $B(p_\infty,0,D')$ is the center of an evolving
$\epsilon/2$-neck defined for an interval of normalized time of
length $2$.

Now since $(M_k,g_k(t),(p_k,0))$ converge geometrically to
$(M_\infty,g_\infty(t),(p_\infty,0))$, by Part 2 of
Proposition~\ref{prop:canonical-neighborhood-stability} for any $L<\infty$, for all $k$
sufficiently large, all points of $B(p_k,0,L)\setminus B(p_k,0,2D')$
are centers of strong $\epsilon$-necks in $(M_k,g_k(t))$. In
particular, for all $k$ sufficiently large, $d_0(p_k,q_k)>L$. Let
$L_k$ be a sequence tending to infinity as $k\rightarrow \infty$.
Passing to a subsequence, we can suppose that every point of
$\left(B(p_k,0,L_k)\setminus B(p_k,0,2D')\right)\subset M_k$ is the
center of a strong $\epsilon$-neck in $(M_k,g_k(0))$.  Of course,
for all $k$ sufficiently large, $q_k\in M_k\setminus B(p_k,0,2D'))$.
By Corollary~\ref{cor:strong-neck-centers-open} the subset of points in $M_k\times
\{0\}$ that are centers of strong $\epsilon$-necks is an open set.
Thus, replacing $q_k$ with another point if necessary we can suppose
that it $q_k$ is a closest point to $p_k$ contained in $M_k\setminus
B(p_k,0,2D')$ with the property $q_k$ is not the center of a strong
$\epsilon$-neck. Then $q_k\in M_k\setminus B(p_k,0,L_k)$ and
$(q_k,0)$ is in the closure of the set of points in $M_k$ that are
centers of strong $\epsilon$-necks in $(M_k,g_k(t))$, and hence by
Part 3 of Proposition~\ref{prop:canonical-neighborhood-stability} each $(q_k,0)$ is the center
of a $2\epsilon$-neck in $(M_k,g_k(t))$.

Let $\gamma_k$ be a minimizing geodesic connecting $(p_k,0)$ to
$(q_k,0)$, and let $\mu_k$ be a minimizing geodesic ray from
$(q_k,0)$ to infinity. Set $Q_k=R(q_k,0)$. Since $(q_k,0)$ is the
center of a $2\epsilon$-neck, from Lemma~\ref{neckseparate} we see
that, provided that $\epsilon$ is sufficiently small,  the
$2\epsilon$-neck centered at $q_k$ separates $p$ from $\infty$, so
that $\gamma_k$ and $\mu_k$ exit this $2\epsilon$-neck at opposite
ends. According to Theorem~\ref{thm:kappa-solution-compactness}, after passing to a
subsequence, the based, rescaled flows
$$(M_k,Q_kg(Q_k^{-1}t),(q_k,0))$$
converge geometrically to a limit. Let $(q_\infty,0)$ be the base
point of the resulting limit. By Part 3 of
Proposition~\ref{prop:canonical-neighborhood-stability}, it is the center of a $4\epsilon$-neck
in the limit.

\begin{claim}[Curvature scale separates $p_k$ and $q_k$]\label{claim:kq-scale-separation}
$d_0^2(p_k,q_k)Q_k\rightarrow \infty$ as $k\rightarrow \infty$.
\end{claim}

\begin{proof}
\uses{lem:uniform-curvature-bound-kappa-solution}
Suppose not. Then by passing to a subsequence we can suppose that
these products are bounded independent of $k$. Then since
$d_0(p_k,q_k)\rightarrow\infty$, we see that $Q_k\rightarrow 0$.
Thus, in the rescaled flows $(M_k,Q_kg_k(Q_k^{-1}t))$ the curvature
at $(p_k,0)$ goes to infinity. But this is impossible since the
$Q_kg_k$-distance from $(p_k,0)$ to $(q_k,0)$ is
$\sqrt{Q_k}d_0(p_k,q_k)$ which is bounded independent of $k$ and the
scalar curvature of $(p,0)$ in the metric $Q_kg_k(0)$ is
$R(p_k,0)Q_k^{-1}=Q_k^{-1}$ tends to $\infty$. Unbounded curvature
at bounded distance contradicts Lemma~\ref{lem:uniform-curvature-bound-kappa-solution}, and this
establishes the claim.
\end{proof}

A subsequence of the based flows $(M_k,Q_kg_k(Q_k^{-1}t),(q_k,0))$
converge geometrically to a $\kappa$-solution. According
Theorem~\ref{topsplit} and Corollary~\ref{localprod}, this limiting
flow is the product of a $2$-dimensional $\kappa$-solution with a
line. Since $M$ is orientable, this $2$-dimensional
$\kappa$-solution is an evolving family of round $2$-spheres. This
implies that for all $k$ sufficiently large, $(q_k,0)$ is the center
of a strong $\epsilon$-neck in $(M_k,g_k(t))$. This is a
contradiction and proves the existence of $D<\infty$ as stated  in
the proposition.


Let $(M_k,g_k(t),(p_k,0))$ be a  sequence of non-compact Ricci flows
based at a soul  $p_k$ of $(M_k,g_k(0))$. We rescale so that
$R(p_k,0)=1$. By Lemma~\ref{topiso}, if $\epsilon$ is sufficiently
small, then $p_k$ cannot be the center of an $\epsilon$-neck. It
follows from Proposition~\ref{prop:canonical-neighborhood-stability} that for any limit of a
subsequence the point $p_\infty$, which is the limit of the $p_k$,
is not the center of an $2\epsilon$-neck in the limit. Since the
limit manifold is orientable, it is either contractible with
strictly positive curvature or is a metric product of a round
$2$-sphere and the line. It follows that the limit manifold has
strictly positive curvature at $(p_\infty,0)$, and hence positive
curvature everywhere. The existence of $D_1<\infty$ as required is
now immediate from Theorem~\ref{thm:kappa-solution-compactness}.

The fact that any soul is disjoint from the middle two-thirds of any
$\epsilon$-neck and the fact that the central $2$-spheres of all
$\epsilon$-necks are isotopic in $M\setminus\{p\}$ are contained in
Lemma~\ref{topiso} and Corollary~\ref{neckseparate}.
\end{proof}




\begin{corollary}[Non-compact $\kappa$-solutions are (capped) strong tubes]\label{cor:noncompact-kappa-solution-tube-structure}
\label{noncompkappa}
\uses{prop:uniform-neck-structure-constants,cor:curvature-derivative-bounds-universal,def:strong-epsilon-tube}
There is $\bar\epsilon_2>0$ such that for any $0<\epsilon\le \bar
\epsilon_2$ the following holds.
 There
is $C_0=C_0(\epsilon)$ such that for any $\kappa>0$ and any non-compact
$3$-dimensional $\kappa$-solution not containing an embedded $\Ar P^2$ with
trivial normal bundle, the zero time-slice is either a strong $\epsilon$-tube
or a $C_0$-capped strong $\epsilon$-tube.
\end{corollary}


\begin{proof}
\uses{prop:uniform-neck-structure-constants,cor:curvature-derivative-bounds-universal}
For $\epsilon>0$ sufficiently small let $D(\epsilon)$ and
$D_1(\epsilon)$ be as in the previous corollary. At the expenses of
increasing these, we can assume that they are at least the constant
$C$ in Corollary~\ref{cor:curvature-derivative-bounds-universal}. We set
$$C_0(\epsilon)=\max(D(\epsilon),D_1(\epsilon)).$$
 If the non-compact $\kappa$-solution has
positive curvature, then the corollary follows immediately from
Proposition~\ref{prop:uniform-neck-structure-constants} and Corollary~\ref{cor:curvature-derivative-bounds-universal}. If the
$\kappa$-solution is the product of an evolving round $S^2$ with the
line, then every point of the zero time-slice is the center of a
strong $\epsilon$-neck for every $\epsilon>0$ so that the zero
time-slice of the solution is a strong $\epsilon$-tube. Suppose the
solution is double covered by the product of an evolving round
$2$-sphere and the line. Let $\iota$ be the involution and take the
product coordinates so that $S^2\times\{0\}$ is the invariant
$2$-sphere of $\iota$ in the zero time-slice. Then any point in the
zero time-slice at distance at least $3\epsilon^{-1}$ from
$P=(S^2\times \{0\})/\iota$ is the center of a strong
$\epsilon$-neck. Furthermore, an appropriate neighborhood of $P$ in
the time zero slice is a $(C,\epsilon)$-cap whose core contains the
$3\epsilon^{-1}$ neighborhood of $P$. The derivative bounds in this
case come from the fact that the metric is close in the
$C^{[1/\epsilon]}$-topology to the standard evolving flow. This
proves the corollary in this case and hence completes the proof.
\end{proof}




Now let us consider compact $\kappa$-solutions.



\begin{theorem}[Structure of compact $\kappa$-solutions]\label{thm:compact-kappa-solution-structure}
\label{thmep3}
\uses{cor:noncompact-kappa-solution-tube-structure,def:strong-epsilon-tube}
There is $\bar \epsilon_3>0$ such that for every $0<\epsilon\le \bar
\epsilon_3$  there is $C_1=C_1(\epsilon)<\infty$ such that one of
the following holds for any $\kappa>0$ and any compact
$3$-dimensional $\kappa$-solution $(M,g(t))$.
\begin{enumerate}
\item[(1)] The manifold $M$ is compact and of constant positive sectional curvature.
\item[(2)] The diameter of $(M,g(0))$ is  less than $C_1\cdot (\max_{x\in
M}R(x,0))^{-1/2}$, and $M$ is diffeomorphic to either $S^3$ or $\Ar
P^3$.
\item[(3)] $(M,g(0))$
 is  a double  $C_1$-capped  strong $\epsilon$-tube.
\end{enumerate}
\end{theorem}


\begin{proof}
\uses{prop:universal-kappa-non-round,cor:noncompact-kappa-solution-tube-structure,cor:compact-asymptotic-soliton-classification,thm:gradient-shrinking-soliton-classification}
First notice that if $(M,g(t))$ is not of strictly positive
curvature, then the universal covering of $(M,g(0))$ is a Riemannian
product $S^2\times \Ar$, and hence $(M,g(0))$ is either non-compact
or finitely covered by the product flow on $S^2\times S^1$. The
former case is ruled out since we are assuming that $M$ is compact
and the latter is ruled out because such flows are not
$\kappa$-non-collapsed for any $\kappa>0$. We conclude that
$(M,g(t))$ is of positive curvature. This implies that the
fundamental group of $M$ is finite. If there were an embedded $\Ar
P^2$ in $M$ with trivial normal bundle, that $\Ar P^2$ cannot
separate (since the Euler characteristic of $\Ar P^2$ is one, it is
not the boundary of a compact $3$-manifold). But a non-separating
surface in $M$ induces a surjective homomorphism of $H_1(M)$ onto
$\mathbb{Z}$. We conclude from this that $M$ does not contain an embedded
$\Ar P^2$ with trivial normal bundle.

 We assume that $(M,g(0))$ is not round so that by Proposition~\ref{prop:universal-kappa-non-round}
 there is a universal $\kappa_0>0$ such that $(M,g(0))$ is a
$\kappa_0$-solution. Let $C_0(\epsilon)$ be the constant from
 Corollary~\ref{cor:noncompact-kappa-solution-tube-structure}.


 \begin{claim}[Large-diameter compact solutions are capped tubes]\label{claim:diameter-canonical-neighborhood-criterion}
\uses{cor:noncompact-kappa-solution-tube-structure}
Assuming that $(M,g(0))$ is compact but not of constant positive
sectional curvature, for each $\epsilon>0$ there is $C_1$ such that
if the diameter of $(M,g(0))$ is greater than $C_1(\max_{x\in
M}R(x,0))^{-1/2}$ then every point of $(M,g(0))$ is either contained
in the core of $(C_0(\epsilon),\epsilon)$-cap or is the center of a
strong $\epsilon$-neck in $(M,g(t))$.
\end{claim}

\begin{proof}
\uses{thm:kappa-solution-compactness,cor:noncompact-kappa-solution-tube-structure,prop:canonical-neighborhood-stability}
Suppose that for some $\epsilon>0$ there is no such $C_1$. Then we
take a sequence of constants $C'_k$ that diverges to $+\infty$ as
$k\rightarrow \infty$ and a sequence $(M_k,g_k(t), (p_k,0))$ of
based $\kappa_0$-solutions such that the diameter of $(M_k,0)$ is
greater than $C'_kR^{-1/2}(p_k,0)$ and yet $(p_k,0)$ is not
contained in the core of a $(C_0(\epsilon),\epsilon)$-cap nor is the
center of a strong $\epsilon$-neck.  We scale $(M_k,g_k(t))$ by
$R(p_k,0)$. This allows us to assume that $R(p_k,0)=1$ for all $k$.
According to Theorem~\ref{thm:kappa-solution-compactness}, after passing to a
subsequence we can assume these based $\kappa$-solutions converge to
a based $\kappa$-solution $(M_\infty,g_\infty(t),(p_\infty,0))$.
Since the diameters of the $(M_k,g_k(0))$ go to infinity, $M_\infty$
is non-compact. According to Corollary~\ref{cor:noncompact-kappa-solution-tube-structure} the point
$p_\infty$ is either the center of a strong $\epsilon$-neck, or is
contained in the core of  a $(C_0(\epsilon),\epsilon)$-cap. Since
$R(p_k,0)=1$ for all $k$, it follows from Parts 1 and 4 of
Proposition~\ref{prop:canonical-neighborhood-stability} that for all $k$ sufficiently large,
$(p_k,0)$ is either the center of a strong $\epsilon$-neck in
$(M_k,g_k(t))$ or is contained in the core of a
$(C_0(\epsilon),\epsilon)$-cap. This is a contradiction, proving the
claim.
\end{proof}


Now it follows from Proposition~\ref{epstopology} that if the
diameter of $(M,g(0))$ is greater than $C_1(\max_{x\in
M}R(x,0))^{-1/2}$ and if it is not of constant positive curvature,
then $M$ is diffeomorphic to either $S^3$, $\Ar P^3$,  $\Ar P^3\#\Ar
P^3$ or is a $S^2$-fibration over $S^1$. On the other hand, since
$M$ is compact of positive curvature its fundamental group is
finite, see Theorem 4.1 on p. 154 of \cite{Petersen}. This rules out
the last two cases. This implies that when $(M,g(0))$ has diameter
greater than $C_1(\max_{x\in M}R(x,0))^{-1/2}$ and is not of
constant positive curvature, it is a double $C_0$-capped
$\epsilon$-tube.

We must consider the case when $(M,g(0))$ is not of constant
positive curvature and its diameter is less than or equal to
$C_1(\max_{x\in M}R(x,0))^{-1/2}$. Since $(M,g(0))$ is not
round, by Corollary~\ref{cor:compact-asymptotic-soliton-classification} its asymptotic soliton is not
compact. Thus, by Theorem~\ref{thm:gradient-shrinking-soliton-classification} its asymptotic soliton is
either $S^2\times \Ar$ or is double covered by this product. This
means that for $t$ sufficiently negative the diameter of $(M,g(t))$
is greater than $C_1(\max_{x\in M}R(x,0))^{-1/2}$. Invoking the
previous result for this negative time tells us that $M$ is
diffeomorphic to $S^3$ or $\Ar P^3$.
\end{proof}

\begin{proposition}[Volume and curvature bounds for small-diameter compact solutions]\label{prop:compact-kappa-solution-volume-curvature-bounds}
\uses{thm:compact-kappa-solution-structure,prop:universal-kappa-non-round,thm:kappa-solution-compactness,cor:noncompact-kappa-solution-tube-structure}
Let $\bar\epsilon_2$ and $\bar \epsilon_3$ be as in
Corollary~\ref{cor:noncompact-kappa-solution-tube-structure} and Theorem~\ref{thm:compact-kappa-solution-structure}, respectively.
For each $0<\epsilon\le \min(\bar \epsilon_2,\bar\epsilon_3)$
let $C_1=C_1(\epsilon)$ be as in Theorem~\ref{thm:compact-kappa-solution-structure}. There is
$C_2=C_2(\epsilon)<\infty$ such that for any $\kappa>0$ and any
compact $\kappa$-solution $(M,g(t))$ the following holds. If
$(M,g(0))$ is not of constant positive curvature and if $(M,g(0))$
is of diameter less than $C_1(\max_{x\in M}R(x,0))^{-1/2}$ then
for any $x\in M$ we have
$$C_2^{-1}R(x,0)^{-3/2}< \Vol(M,g(0))< C_2R(x,0)^{-3/2}.$$
In addition, for any $y\in M$ and any $2$-plane $P_y$ in $T_yM$  we
have
$$C_2^{-1}< \frac{K(P_y)}{R(x,0)}< C_2,$$
where $K(P_y)$ is the sectional curvature in the $P_y$-direction.
\end{proposition}


\begin{proof}
\uses{prop:universal-kappa-non-round,thm:kappa-solution-compactness}
The result is immediate from Corollary~\ref{prop:universal-kappa-non-round} and
Theorem~\ref{thm:kappa-solution-compactness}.
\end{proof}

\begin{remark}[No universal volume lower bound for round solutions]\label{rem:round-kappa-solution-volume-remark}
For a round $\kappa$-solution $(M,g(t))$ we have $R(x,0)=R(y,0)$ for
all $x,y\in M$, and the volume of $(M,g(0))$ is bounded above by a
constant times $R(x,0)^{-3/2}$. There is no universal lower bound to
the volume in terms of the curvature. The lower bound takes the form
$C_2|\pi_1(M)|^{-1}R(x,0)^{-3/2}$, where $|\pi_1(M)|$ is the order
of the fundamental group $\pi_1(M)$.
\end{remark}

Let us summarize our results.

\begin{theorem}[Summary structure theorem for $\kappa$-solutions]\label{thm:kappa-solution-structure-summary}
\label{kappasummary}
\uses{thm:compact-kappa-solution-structure,cor:noncompact-kappa-solution-tube-structure,prop:compact-kappa-solution-volume-curvature-bounds,cor:curvature-derivative-bounds-universal}
There is  $\bar \epsilon>0$ such that the following is true for any
$0<\epsilon<\bar \epsilon$.  There is $C=C(\epsilon)$ such that for
any $\kappa>0$ and any $\kappa$-solution $(M,g(t))$ one of the
following holds.
\begin{enumerate}
\item[(1)] $(M,g(t))$ is round for all $t\le 0$. In this case $M$ is
diffeomorphic to the quotient of $S^3$ by a finite subgroup of
$SO(4)$ acting freely.
\item[(2)] $(M,g(0))$ is compact and of positive curvature. For any $x,y\in M$
and any $2$-plane $P_y$ in $T_yM$
 we have \begin{eqnarray*}
 C^{-1/2}R(x,0)^{-1} &  <\diam(M,g(0))
 &  <CR(x,0)^{-1/2} \\
 C^{-1}R(x,0)^{-3/2} & < \Vol(M,g(0)) & < CR(x,0)^{-3/2} \\
C^{-1}R(x,0) & < K(P_y) & < CR(x,0).
\end{eqnarray*}
  In this case $M$ is diffeomorphic either to
 $S^3$ or to $\Ar P^3$.
 \item[(3)] $(M,g(0))$ is of positive curvature and is a double  $C$-capped strong
 $\epsilon$-tube, and in
 particular $M$ is diffeomorphic to $S^3$ or to $\Ar P^3$.
 \item[(4)] $(M,g(0))$ is of positive curvature and is a
 $C$-capped strong $\epsilon$-tube and $M$ is diffeomorphic to $\Ar ^3$.
 \item[(5)] $(M,g(0))$ is isometric to the quotient of the product of
 a round $S^2$ and $\Ar$ by a
 free, orientation-preserving involution. It is a  $C$-capped
 strong $\epsilon$-tube and is diffeomorphic to a punctured $\Ar P^3$.
 \item[(6)] $(M,g(0))$ is isometric to the product of a round $S^2$ and $ \Ar$
  and is a strong $\epsilon$-tube.
 \item[(7)] $(M,g(0))$ is isometric to a product $\Ar P^2\times \Ar$, where the metric on $\Ar P^2$
 is of constant Gaussian curvature.
\end{enumerate}

In particular, in all cases except the first two and the last one,
all points of $(M,g(0))$ are either contained in the core of a
$(C,\epsilon)$-cap or are the centers of a strong $\epsilon$-neck in
$(M,g(0))$.

Lastly, in all cases we have \begin{eqnarray}\sup_{p\in M, t\le
0}\frac{|\nabla
R(p,t)|}{R(p,t)^{3/2}} & <C \label{nablaR} \\
\sup_{p\in M,t\le 0}\frac{\left|\partial R(p,t)/\partial
t\right|}{R(p,t)^2} & <C.\label{dRdt}
\end{eqnarray}
\end{theorem}



An immediate consequence of this result is:

\begin{corollary}[Existence of canonical neighborhoods]\label{cor:canonical-neighborhood-existence}
\label{kappacannbhd}
\uses{thm:kappa-solution-structure-summary,def:canonical-neighborhood}
For every $0<\epsilon\le\bar\epsilon'$ there is
$C=C(\epsilon)<\infty$ such that every point in a $\kappa$-solution
has a strong $(C,\epsilon)$-canonical neighborhood unless the
$\kappa$-solution is a product $\Ar P^2\times \Ar$.
\end{corollary}


\begin{corollary}[Canonical neighborhoods pass to limits of generalized flows]\label{cor:canonical-neighborhood-limit-passage}
\label{limitcannbhd}
\uses{cor:canonical-neighborhood-existence,def:strong-epsilon-neck-canonical-neighborhood,def:generalized-ricci-flow}
Fix $0<\epsilon\le \bar\epsilon'$, and let $C(\epsilon)$ be as in
the last corollary. Suppose that $({\mathcal M}_n,G_n,x_n)$ is a
sequence of based, generalized Ricci flows with ${\bf t}(x_n)=0$ for
all $n$. Suppose that none of the time-slices of the ${\mathcal
M}_n$ contain embedded $\Ar P^2$'s with trivial normal bundle.
Suppose also that there is a smooth limiting flow
$(M_\infty,g_\infty(t),(x_\infty,0))$ defined for $-\infty<t\le 0$
that is a $\kappa$-solution. Then for all $n$ sufficiently large
$x_n$ has a strong $(C,\epsilon)$-canonical neighborhood in
$({\mathcal M}_n,G_n,x_n)$.
\end{corollary}


\begin{proof}
\uses{cor:canonical-neighborhood-existence,prop:canonical-neighborhood-stability}
The limiting manifold $M_\infty$ cannot contain an embedded $\Ar
P^2$ with trivial normal bundle. Hence, by the previous corollary,
the point $(x_\infty,0)$ has a strong $(C,\epsilon)$-canonical
neighborhood in the limiting flow. If the limiting $\kappa$-solution
is round, then for all $n$ sufficiently large $x_n$ is contained in
a component of the zero time-slice that is $\epsilon$-round. If
$(x_\infty,0)$ is contained in a $C$-component of the zero
time-slice of the limiting $\kappa$-solution, then for all $n$
sufficiently large $x_n$ is contained in a $C$-component of the zero
time-slice of ${\mathcal M}_n$. Suppose that $(x_\infty,0)$ is the
center of a strong $\epsilon$-neck in the limiting flow. This neck
extends backwards in the limiting solution some amount past an
interval of normalized time of length  $1$, where by continuity it
is an evolving $\epsilon$-neck defined backwards for an interval of
normalized time of length  greater than $1$. Then by Part 2 of
Proposition~\ref{prop:canonical-neighborhood-stability}, any family of metrics on this neck
sufficiently close to the limiting metric will determine an strong
$\epsilon$-neck. This implies that for all $n$ sufficiently large
$x_n$ is the center of a strong $\epsilon$-neck in $({\mathcal
M}_n,G_n)$. Lastly, if $(x_\infty,0)$ is contained in the core of a
$(C,\epsilon)$-cap in the limiting flow, then by Part 1 of
Proposition~\ref{prop:canonical-neighborhood-stability} for all $n$ sufficiently large $x_n$ is
contained in the core of a $(C,\epsilon)$-cap in $({\mathcal
M}_n,G_n)$.
\end{proof}
